\documentclass[11pt]{article}

\usepackage[utf8]{inputenc}
\usepackage[T1]{fontenc}
\usepackage{lmodern}
\usepackage[spanish,es-noshorthands,es-tabla]{babel}
\usepackage[a4paper,margin=2.15cm]{geometry}
\usepackage{amsmath,amssymb,amsfonts}
\usepackage{booktabs}
\usepackage{longtable}
\usepackage{tabularx}
\usepackage{array}
\usepackage{graphicx}
\usepackage{float}
\usepackage{hyperref}
\usepackage{xcolor}
\usepackage{ragged2e}

\hypersetup{
  colorlinks=true,
  linkcolor=blue!55!black,
  urlcolor=blue!55!black,
  citecolor=blue!55!black
}

\setlength{\emergencystretch}{3em}
\renewcommand{\arraystretch}{1.13}
\setlength{\tabcolsep}{4pt}
\graphicspath{{../library_figures/by_report/unified_chua_fractional/pdf/}}

\newcolumntype{L}[1]{>{\raggedright\arraybackslash}p{#1}}
\DeclareRobustCommand{\codepath}[1]{\begingroup\Urlmuskip=0mu plus 2mu\path{#1}\endgroup}
\makeatletter
\renewcommand*\l@subsection{\@dottedtocline{2}{1.5em}{2.8em}}
\makeatother
\newcommand{\reportinclude}[2][]{\includegraphics[#1]{#2}}
\newcommand{\CaputoD}{\,{}^{C}D_t^q}
\newcommand{\ii}{\mathrm{i}}
\newcommand{\R}{\mathbb{R}}
\newcommand{\TARGET}{\texttt{TARGET}}
\newcommand{\STABLE}{\texttt{estable}}
\newcommand{\EXPERIMENTAL}{\texttt{experimental}}
\newcommand{\INTERNAL}{\texttt{interno}}
\newcommand{\LEGACY}{\texttt{legacy}}
\newcommand{\cmdbreak}[2]{\begin{tabular}[t]{@{}l@{}}\texttt{#1}\\\texttt{#2}\end{tabular}}
\newcommand{\claimtag}[1]{\,\textnormal{\footnotesize[afirmacion]}}
\newcommand{\localcontract}[1]{\,\textnormal{\footnotesize[contrato local]}}
\newcommand{\heuristictag}[1]{\,\textnormal{\footnotesize[heuristica]}}
\newcommand{\diagnostictag}[1]{\,\textnormal{\footnotesize[diagnostico]}}

\title{Hidden Attractors Fractional Order\\
\large Guía de usuario, fundamentos matemáticos, ejemplos y validación}
\author{Proyecto Hidden Attractors Fractional Order}
\date{Actualizado: \today}

\begin{document}
\maketitle
\setcounter{tocdepth}{2}
\tableofcontents

\section{Propósito}

La librería ofrece caracterización de sistemas dinámicos y series temporales,
además de un protocolo específico para estudiar candidatos a atractores
ocultos en sistemas de Chua de orden entero y fraccionario. Integra modelos con
derivada ordinaria, el camino Caputo usado por los casos científicos, un núcleo
experimental con otras definiciones fraccionarias, generación de semillas por función
descriptiva, integración numérica, análisis de trayectorias, pruebas de cuenca,
robustez, diagnósticos de caos y trazabilidad bibliográfica. Para localización
de candidatos, el objetivo es construir:
\[
\text{modelo}\to\text{contrato numérico}\to\text{semilla}\to
\text{trayectoria}\to\text{robustez}\to\text{cuencas}\to\text{decisión}.
\]

Una trayectoria acotada, un gráfico tridimensional, una FFT, un diagrama de Poincaré o un
exponente de Lyapunov estimado no demuestran por sí solos que un atractor sea
oculto. La definición de atractor oculto exige que la cuenca de atracción no
intersecte las vecindades de los puntos de equilibrio relevantes~\cite{LeonovKuznetsov2013,KuznetsovEtAl2017}. Por eso la librería separa
tres niveles:

\begin{center}
\small
\begin{tabular}{@{}L{0.23\textwidth}L{0.39\textwidth}L{0.29\textwidth}@{}}
\toprule
Nivel & Qué significa & Qué no significa \\
\midrule
Diagnóstico visual & Una figura ayuda a inspeccionar geometría, transitorios, espectro o cuencas finitas. & No es una prueba matemática global. \\
Evidencia numérica bajo contrato & Un resultado se obtuvo con \(q\), \(h\), horizonte, memoria, backend y tolerancias declaradas. & No se extrapola fuera de ese contrato. \\
Decisión promovida & Un paquete de evidencia pasó por el protocolo vigente y conserva estado, archivos y trazabilidad. & No borra discrepancias ni convierte pruebas rápidas en validación publicada. \\
\bottomrule
\end{tabular}
\end{center}

\section{Estado de la librería y reglas de lectura}

La versión documentada es \codepath{1.2.0}. La superficie de usuario se divide
por estabilidad:

\begin{itemize}
  \item \STABLE: modelos, registro de sistemas, etiquetas de cuenca e IO
  portable.
  \item \EXPERIMENTAL: análisis, generación de semillas, integradores,
  plotting y workflows.
  \item \INTERNAL: backends nativos, cache, CLI interna y utilidades de
  compilación.
  \item \LEGACY: scripts históricos congelados.
\end{itemize}

Hay cuatro advertencias que deben permanecer visibles en todo el documento:

\begin{enumerate}
  \item La función descriptiva y Nyquist generan semillas. No prueban
  existencia global de un ciclo ni ocultedad
  \cite{GenesioTesi1993,KuznetsovEtAl2017}.
  \item En sistemas autónomos de Caputo no se deben afirmar órbitas periódicas
  exactas como si fueran soluciones ordinarias. La memoria rompe~la
  periodicidad exacta salvo casos degenerados~\cite{TavazoeiHaeri2009}.
  \item Las pruebas rápidas de CI no sustituyen los contratos cuantitativos
  publicados que viven bajo \codepath{validation/}.
  \item Los diagnósticos finito-temporales no se presentan como certificación
  estricta de caos ni como validación automática de ocultedad.
\end{enumerate}

\section{Instalación y primer contacto}

La librería requiere Python 3.11 o superior. Desde la raíz del repositorio:

\begin{verbatim}
python -m pip install -e version_2
\end{verbatim}

Para desarrollo, análisis o documentación:

\begin{verbatim}
python -m pip install -e "version_2[dev,analysis,docs,legacy]"
\end{verbatim}

Las dependencias base son \codepath{numpy}, \codepath{matplotlib},
\codepath{scipy} y \codepath{numba}. El extra \codepath{analysis} añade
métricas externas; el extra \codepath{docs} instala la cadena MkDocs; el extra
\codepath{legacy} conserva dependencias para scripts históricos.

La CLI principal es:

\begin{verbatim}
hidden-attractors --help
\end{verbatim}

Flujo mínimo recomendado:

\begin{verbatim}
hidden-attractors init -e chua_fractional
hidden-attractors inspect-config -p chua_fractional
hidden-attractors run -p chua_fractional
hidden-attractors run -c path/to/config.yaml
\end{verbatim}

Los presets mantenidos por la CLI son:

\begin{center}
\small
\begin{tabular}{@{}L{0.28\textwidth}L{0.47\textwidth}L{0.18\textwidth}@{}}
\toprule
Preset & YAML empaquetado & Lectura \\
\midrule
\codepath{chua_integer} & \codepath{chua_integer_centered_lure_df.yaml} & Ejemplo entero \(q=1\). \\
\codepath{chua_fractional} & \codepath{chua_fractional_centered_lure_df.yaml} & Ejemplo fraccionario no suave. \\
\codepath{chua_bifurcation} & \codepath{chua_fractional_bifurcation.yaml} & Barrido paramétrico. \\
\codepath{chua_basin} & \codepath{chua_fractional_basin.yaml} & Corte de cuenca. \\
\codepath{chua_full_protocol} & \codepath{chua_full_protocol.yaml} & Esqueleto de protocolo completo. \\
\bottomrule
\end{tabular}
\end{center}

\section{Manual de usuario}

El manual de usuario canónico de la librería se encuentra en \codepath{USER_MANUAL.md}. Este archivo contiene las especificaciones detalladas para el uso práctico de la librería. Los metadatos de configuración y sincronización documental se definen en el manifiesto central \codepath{docs/manual_manifest.yaml}, mientras que las afirmaciones científicas oficialmente defendibles están gobernadas por la matriz de validación en \codepath{validation/references/thesis_claims_matrix.md}.

El presente reporte en LaTeX/PDF representa la versión formal extendida para la tesis y artículos del proyecto, mientras que \codepath{USER_MANUAL.md} es la ruta de uso práctico recomendada.

Este reporte no introduce resultados científicos nuevos respecto a los manifiestos de validación. Su función es explicar la metodología, los comandos, las salidas esperadas y la interpretación de la evidencia bajo el contrato numérico declarado.

La versión 1.2.0 reconoce dos usos de la misma base numérica. El primero es la
localización y verificación finito-temporal de candidatos a atractores ocultos.
El segundo es la caracterización independiente de sistemas dinámicos,
trayectorias y series temporales, sin exigir que exista una búsqueda de
ocultedad. En este segundo uso se incluyen equilibrios, Jacobianos,
acotamiento, FFT/PSD, prueba 0--1, secciones de Poincaré, bifurcaciones y
estimadores de Lyapunov.

\begin{quote}
\textbf{Advertencia:} La función descriptiva, el análisis de Nyquist, la continuación numérica y la visualización de una trayectoria acotada no prueban la ocultedad de un atractor. La ocultedad se evalúa rigurosamente analizando las vecindades de todos los puntos de equilibrio bajo los radios, tiempos de integración, integrador numérico y política de memoria declarados.
\end{quote}

\section{Estructura del repositorio}

\begin{center}
\small
\begin{tabular}{@{}L{0.31\textwidth}L{0.59\textwidth}@{}}
\toprule
Ruta & Papel \\
\midrule
\codepath{hidden_attractors/models/} & Parámetros, campos vectoriales, no linealidades, equilibrios y Jacobianos. \\
\codepath{hidden_attractors/systems/} & Registro público de sistemas y forma de Lur'e. \\
\codepath{hidden_attractors/seed_generation/} & Semillas por función descriptiva, Nyquist y Lur'e. \\
\codepath{hidden_attractors/fractional/} & Definiciones, contratos, operadores muestreados y solvers fraccionarios; separa kernel, malla, memoria y backend. \\
\codepath{hidden_attractors/integrations/} & Integradores enteros y fraccionarios, selector de backend y compatibilidad de \(q\). \\
\codepath{hidden_attractors/analysis/} & Métricas genéricas de trayectoria, bifurcación, FFT/PSD, Poincaré, prueba 0--1 y Lyapunov. \\
\codepath{hidden_attractors/plotting/} & Figuras de fase, series, espectros, transferencia, continuación y cuencas. \\
\codepath{hidden_attractors/workflows/} & Protocolo oficial, especificaciones, robustez, cuencas, refinamiento y corridas reproducibles. \\
\codepath{hidden_attractors/native/} & Backends C y clasificadores de cuenca. Son internos y Chua-específicos. \\
\codepath{hidden_attractors/configs/examples/} & YAML listos para usar desde la CLI\@. \\
\codepath{examples/} & Ejemplos pequeños para aprender la API\@. \\
\codepath{validation/} & Evidencia promovida, manifiestos, contratos, comparaciones y generadores editoriales específicos. \\
\codepath{library_figures/} & Figuras promovidas y manifiestos de figuras; \codepath{docs/assets/} queda solo para material documental. \\
\bottomrule
\end{tabular}
\end{center}

\section{Inventario de API de la libreria}

La referencia de la superficie pública mantenida se encuentra en
\codepath{docs/api_reference.md}. Es una guía curada que se contrasta con las
exportaciones de \codepath{hidden_attractors}, \codepath{analysis},
\codepath{integrations} y \codepath{plotting}; no pretende enumerar símbolos
privados ni backends internos. Este reporte sí enumera de forma explícita los
procedimientos numéricos y las funciones gráficas que una persona usuaria
puede invocar.

La lectura pública recomendada se resume así:

\begin{center}
\small
\begin{tabular}{@{}L{0.28\textwidth}L{0.32\textwidth}L{0.30\textwidth}@{}}
\toprule
Familia & Modulos principales & Uso documentado \\
\midrule
Registro de sistemas & \codepath{systems}, \codepath{models} & Definir campos vectoriales, parametros, equilibrios, Jacobianos y forma de Lur'e. \\
Semillas Lur'e/DF & \codepath{seed_generation}, \codepath{lure} & Generar semillas heuristicas por DF/Nyquist; no prueban ocultedad. \\
Integracion & \codepath{integrations}, \codepath{solvers}, \codepath{native} & Ejecutar trayectorias enteras o Caputo bajo contrato de solver y memoria. \\
Workflows & \codepath{workflows}, \codepath{continuation} & Unir semillas, continuacion, referencia dinamica, robustez, cuencas y reportes. \\
Verificacion & \codepath{verification}, \codepath{basins} & Clasificar equilibrios, muestrear vecindades y decidir contactos. \\
Diagnosticos & \codepath{analysis}, \codepath{diagnostics} & Métricas explícitas de trayectoria, FFT/PSD, Poincare, Lyapunov, 0--1 y acotamiento. \\
Figuras e IO & \codepath{plotting}, \codepath{io} & Exportar PNG/PDF/JSON/CSV reproducibles y manifiestos. \\
\bottomrule
\end{tabular}
\end{center}

Para nuevos sistemas tipo Lur'e, los puntos de entrada minimos son \codepath{ChaoticSystem}, \codepath{register_system}, \codepath{LureSystem}, \codepath{WorkflowInputSpec}, \codepath{load_config} e \codepath{integrate}. Cualquier simbolo nuevo anadido al paquete debe registrarse en \codepath{docs/api_reference.md} y, si cambia la ruta de usuario, en \codepath{docs/quick_start.md}, \codepath{docs/getting_started.md}, \codepath{USER_MANUAL.md} y este reporte.

\section{Fundamentos: sistemas dinámicos y caos}

Un sistema dinámico autónomo ordinario se escribe como
\[
\dot{x}=F(x), \qquad x(t)\in\R^n.
\]
Su solución depende de una condición inicial \(x(0)=x_0\). Un conjunto
atractor es, de manera operativa, una región invariante que atrae un conjunto
de condiciones iniciales y dentro de la cual la dinámica permanece acotada.
La librería trabaja siempre con aproximaciones finito-temporales: se integra
hasta un tiempo \(t_f\), se descarta un transitorio \(t_{\rm burn}\), y se
analiza la cola de la trayectoria.

En la práctica numérica, caos suele inspeccionarse mediante sensibilidad a
condiciones iniciales, geometría no trivial, espectro amplio, secciones de
Poincaré y exponentes de Lyapunov. Ninguna de estas piezas aisladas decide
ocultedad. La clasificación de atractor oculto se define por su cuenca: si la
cuenca intersecta vecindades de algún equilibrio, el atractor es
autoexcitado bajo ese contrato; si no se detectan intersecciones en un
muestreo finito, la conclusión debe ser compatible con ocultedad bajo radios
probados, no una prueba global absoluta.

\subsection{Cuándo llamar caótico a un sistema en una validación numérica}

En matemática, ``caos'' puede formalizarse de varias maneras: sensibilidad a
condiciones iniciales, mezcla topológica, órbitas periódicas densas, entropía
positiva u otros criterios. En una librería computacional de sistemas
continuos, la lectura práctica es más modesta: se busca un conjunto acotado,
no periódico y sensible a condiciones iniciales, y se exige que diagnósticos
independientes cuenten la misma historia. Bajo el contrato de este proyecto,
una afirmación razonable es \emph{evidencia fuerte de caos numérico} cuando se
cumplen simultáneamente estas condiciones:

\begin{enumerate}
  \item La trayectoria post-transitorio permanece acotada y no converge a un
  equilibrio ni a una órbita periódica simple.
  \item El mayor exponente de Lyapunov estimado es positivo y se mantiene
  positivo al refinar \(h\), cambiar el horizonte o comparar métodos
  compatibles.
  \item En un flujo autónomo continuo aparece un exponente cercano a cero,
  asociado a la dirección tangente de la trayectoria. En un sistema
  disipativo la suma de exponentes debe ser negativa en promedio.
  \item La sección de Poincaré no se reduce a un número finito pequeño de
  puntos repetidos. Debe mostrar una nube, bandas o una estructura no trivial
  compatible con retornos aperiódicos.
  \item El espectro de la serie temporal no está formado sólo por una
  frecuencia fundamental y sus armónicos. Un espectro amplio o de banda ancha
  apoya la lectura de aperiodicidad.
  \item La geometría del espacio de fases muestra estiramiento y plegamiento,
  con estructura acotada no trivial, y no un transitorio que simplemente aún
  no terminó.
\end{enumerate}

Estas condiciones no prueban ocultedad. Un atractor puede ser caótico y aun
así autoexcitado. La ocultedad se decide por la relación entre cuenca y
equilibrios; el caos se decide por la dinámica sobre el conjunto atrayente.

\begin{center}
\footnotesize
\begin{longtable}{@{}L{0.18\textwidth}L{0.30\textwidth}L{0.23\textwidth}L{0.20\textwidth}@{}}
\toprule
Diagnóstico & Lectura caótica esperada & Lectura no caótica típica & Cuidado principal \\
\midrule
Lyapunov & \(\lambda_{\max}>0\); en flujos autónomos, un exponente cercano a cero; suma negativa si el sistema es disipativo. & Todos los exponentes negativos: equilibrio o ciclo estable; \(\lambda_{\max}\approx0\) sin positivo robusto: cuasiperiodicidad o caso no concluyente. & Estimaciones cortas pueden confundir transitorios con sensibilidad real. \\
Poincaré & Conjunto de cruces disperso, bandas o estructura fractal finita. & Un punto: ciclo de periodo uno; pocos puntos: órbita periódica de periodo finito; curva suave: toro cuasiperiódico. & La sección debe ser transversal y usar una dirección de cruce. \\
Espectro amplio & Energía repartida en un intervalo de frecuencias, con picos mezclados con banda continua. & Líneas discretas dominantes y armónicos limpios. & Ruido numérico o muestreo pobre también ensanchan espectros. \\
Geometría & Atractor acotado con plegamiento, lóbulos, capas o estructura no trivial. & Convergencia a punto, círculo simple o curva cerrada lisa. & Una figura bonita no prueba sensibilidad ni robustez. \\
\bottomrule
\end{longtable}
\end{center}

\subsection{Ejemplo pedagógico: qué sí y qué no decide una trayectoria}

Un ejemplo clásico como Lorenz ayuda a explicar sensibilidad, mezcla y señales no periódicas, pero no pertenece a la metodología activa de este reporte. Por eso no se usan figuras de Lorenz ni galerías externas como evidencia de los casos Chua. La regla de lectura es la misma para todos los sistemas: una trayectoria o un espectro son diagnósticos; la ocultedad requiere vecindades de todos los equilibrios y un contrato numérico reproducible.
\section{Fundamentos del cálculo fraccionario}
\label{sec:fundamentos-calculo-fraccionario}

Para \(0<q<1\), la derivada de Caputo usada por la librería se define
como~\cite{Caputo1967,Podlubny1999}
\[
{}^{C}D_t^q x(t)=
\frac{1}{\Gamma(1-q)}
\int_0^t \frac{x'(\tau)}{{\left(t-\tau\right)}^q}\,d\tau .
\]
El valor de la derivada en \(t\) depende de toda la historia previa. Por eso
un integrador fraccionario no solo necesita \(x_n\), sino una representación
de los estados anteriores. Dos consecuencias prácticas son importantes:

\begin{itemize}
  \item La memoria completa es costosa, pero sirve como contrato de referencia
  para integradores ABM de Caputo.
  \item Las ventanas de memoria o backends acelerados son útiles, pero su
  conclusión depende de la ventana, de \(h\), de \(t_f\) y del error aceptado.
\end{itemize}

\subsection{Forma integral de Volterra}

La forma diferencial de Caputo
\[
{}^{C}D_t^q x(t)=F(t,x(t)),\qquad x(0)=x_0,\qquad 0<q<1,
\]
es equivalente, bajo regularidad suficiente, a la ecuación integral de
Volterra
\begin{equation}
x(t)=x_0+
\frac{1}{\Gamma(q)}
\int_0^t {\left(t-\tau\right)}^{q-1}F(\tau,x(\tau))\,d\tau .
\label{eq:caputo_volterra}
\end{equation}
Esta ecuación muestra por qué los métodos fraccionarios son métodos con
memoria: el estado en \(t\) acumula todo el historial de \(F\) con un kernel
singular \({\left(t-\tau\right)}^{q-1}\). ABM y otros integradores de Caputo discretizan
precisamente esta forma de Volterra.

Para \(m-1<q<m\), la forma general incluye las condiciones iniciales enteras:
\[
x(t)=
\sum_{k=0}^{m-1}x^{(k)}(0)\frac{t^k}{k!}
+\frac{1}{\Gamma(q)}
\int_0^t {\left(t-\tau\right)}^{q-1}F(\tau,x(\tau))\,d\tau .
\]
En esta librería el caso usado en Chua fraccionario es el conmensurado
\(0<q<1\), por lo que basta una condición inicial \(x(0)\).

\subsection{Transformada de Laplace de la derivada de Caputo}

La transformada de Laplace es una razón práctica para usar Caputo en
problemas de valor inicial: las condiciones iniciales aparecen como derivadas
enteras ordinarias. Si \(X(s)=\mathcal{L}\{x(t)\}(s)\) y \(m-1<q<m\),
\begin{equation}
\mathcal{L}\left\{{}^{C}D_t^q x(t)\right\}(s)
=
s^q X(s)-
\sum_{k=0}^{m-1}s^{q-1-k}x^{(k)}(0).
\label{eq:caputo_laplace_general}
\end{equation}
En el caso \(0<q<1\),
\begin{equation}
\mathcal{L}\left\{{}^{C}D_t^q x(t)\right\}(s)
=s^q X(s)-s^{q-1}x(0).
\label{eq:caputo_laplace_q01}
\end{equation}
Esta identidad conecta directamente el problema lineal de Caputo con
funciones de Mittag-Leffler y con el criterio de estabilidad de Matignon.

\subsection{Derivada de Weyl-Liouville}

La derivada de Weyl-Liouville se define sobre una historia extendida, no
sólo desde \(t=0\). Para \(0<q<1\), la forma izquierda puede escribirse como
\begin{equation}
{}_{-\infty}D_t^q x(t)=
\frac{d}{dt}
\left[
\frac{1}{\Gamma(1-q)}
\int_{-\infty}^{t}\frac{x(\tau)}{{\left(t-\tau\right)}^q}\,d\tau
\right].
\label{eq:weyl_liouville_left}
\end{equation}
Su contraparte derecha integra desde \(t\) hasta \(+\infty\). La forma de
Weyl-Liouville es útil para análisis armónico porque, en condiciones
adecuadas,
\[
\mathcal{F}\left\{{}_{-\infty}D_t^q x(t)\right\}(\omega)
= {\left(\ii\omega\right)}^q\widehat{x}(\omega),
\]
con una rama compleja elegida. Esta propiedad explica por qué aparecen
expresiones como \({\left(\ii\omega\right)}^q\) en funciones de transferencia
fraccionarias. Sin embargo, no debe confundirse con un problema de valor
inicial de Caputo: Weyl-Liouville presupone una historia desde
\(-\infty\), mientras que Caputo en la librería parte de una condición inicial
y de memoria acumulada desde \(t=0\).

\subsection{Criterio de Matignon y su origen}

Para sistemas lineales conmensurados
\[
{}^{C}D_t^q x = A x,
\]
el criterio de Matignon \cite{Matignon1996} establece estabilidad asintótica local si todos los
autovalores \(\lambda_i\) de \(A\) satisfacen
\[
|\arg(\lambda_i)|>\frac{q\pi}{2}.
\]
Cuando \(q=1\), esto se reduce al criterio ordinario
\(\operatorname{Re}(\lambda_i)<0\). Matignon clasifica estabilidad local de
equilibrios. No clasifica ocultedad ni caos.

La deducción se ve al aplicar Laplace al sistema lineal. Para \(0<q<1\),
\[
\left(s^q I-A\right)X(s)=s^{q-1}x(0).
\]
Los modos quedan gobernados por los valores propios de \(A\): cada modo
escalar satisface
\[
{}^{C}D_t^q y=\lambda y,\qquad
y(t)=y(0)E_q(\lambda t^q),
\]
donde \(E_q\) es la función de Mittag-Leffler. La estabilidad de ese modo
depende de la región angular donde \(E_q(\lambda t^q)\) decae cuando
\(t\to\infty\). Esa región es el complemento del sector
\(|\arg(\lambda)|\le q\pi/2\). Equivalentemente, los ceros de
\(s^q-\lambda\) no deben producir ramas con parte real positiva. De ahí sale
la condición angular:
\[
|\arg(\lambda_i)|>\frac{q\pi}{2}
\quad\text{para todo valor propio }\lambda_i.
\]
Cuando \(q\) disminuye, el sector inestable se estrecha alrededor del eje real
positivo; por eso algunos valores propios con parte real positiva pueden ser
compatibles con estabilidad fraccionaria si su argumento queda fuera de ese
sector.

\begin{figure}[H]
\centering
\reportinclude[width=0.78\textwidth]{matignon_complex_plane.pdf}
\caption{División del plano complejo según Matignon. Para un orden conmensurado \(q\), los valores propios deben quedar fuera del sector \(|\arg(\lambda)|\le q\pi/2\). La zona sombreada ilustra el caso \(q=0.8\).}
\end{figure}

\subsection{Catálogo de definiciones fraccionarias y frontera de significado}
\label{sec:fundamentos-fraccionarios-diversos}

El núcleo general vive en \codepath{hidden_attractors/fractional/}. Su regla
principal es no identificar cuatro objetos distintos: la definición continua
del operador, la discretización, la semántica de los datos iniciales y el
backend que ejecuta el cálculo. En particular, evaluar una derivada sobre una
historia ya muestreada no resuelve automáticamente una ecuación diferencial
fraccionaria (FDE). Esta separación permite conservar el camino Caputo usado en
los casos Chua y, al mismo tiempo, estudiar otras familias sin cambiar
silenciosamente el modelo matemático.

La búsqueda bibliográfica se apoyó en SciSpace. La bitácora de preguntas,
resultados seleccionados y límites del índice se conserva en
\codepath{docs/scispace_fractional_method_evidence.md}. SciSpace se usó para
localizar y comparar candidatos; las fórmulas se contrastaron con artículos
primarios, libros y páginas editoriales. Algunas referencias, entre ellas la
transformación de Zheng, no fueron recuperadas por el índice exacto y se
conservaron sólo después de contrastar su DOI. Por tanto, SciSpace no se trata
como oráculo matemático ni como evidencia de corrección del código.

\begin{center}
\footnotesize
\begin{longtable}{@{}L{0.18\textwidth}L{0.23\textwidth}L{0.19\textwidth}L{0.29\textwidth}@{}}
\toprule
Familia & Objeto implementado & Estado & Contrato esencial \\
\midrule
Caputo & ABM/PECE y EFORK-3 & Implementado & IVP con datos enteros y memoria declarada; el reporte científico usa principalmente \(0<q<1\). \\
Grünwald--Letnikov (GL) & Pesos, convolución directa, FFT y solver explícito & Experimental salvo kernels de referencia & Historia terminal truncada o desplazamiento Caputo explícito; malla uniforme. \\
Riemann--Liouville (RL) & Operador GL/CQ sobre muestras & Experimental & No convierte \(x(a)\) en un IVP RL; exige semántica \codepath{operator_only}. \\
RL templada & Conjugación exponencial, GL, CQ BDF1/BDF2 e historia rápida FBDF1/GNGF2 & Experimental & GL, CQ e historia rápida son operadores muestreados con semántica \codepath{operator_only}; el parámetro de templado forma parte de la definición. \\
Caputo templada & CQ BDF1/BDF2, historia rápida FBDF1/GNGF2 y ABM/PECE por conjugación & Experimental & CQ e historia rápida son operadores sobre muestras; ABM/PECE es un solver conmensurado. Todos fijan \(\lambda\ge0\) y la condición puntual de forma explícita. \\
Orden variable GL & GL con \(q=q_n\) & Experimental & Operador sobre muestras: el orden se fija en cada salida; otras convenciones no son equivalentes. \\
Caputo de orden variable tipo III & L1 implícito con corrector Picard & Experimental & Solver con perfil temporal prescrito; \(\alpha(t_n)\) se usa en toda la fila histórica. \\
Orden distribuido & Cuadratura en orden más GL & Experimental & Doble discretización: nodos/pesos en \(q\) y malla temporal. \\
Conformable & Derivada local de Khalil & Experimental & Operador local; no posee la memoria de Caputo, RL o GL. \\
Conformable & RK4 en reloj conformable & Experimental & Solver conmensurado sobre \(\tau=(t-a)^q/q\); malla física no uniforme y sin memoria hereditaria. \\
Caputo--Fabrizio & Operador recurrente exponencial & \codepath{research_required} & Normalización explícita; operador de datos, no solver FDE promovido. \\
Hadamard y Caputo--Hadamard & CQ BDF1/BDF2 en tiempo logarítmico & Experimental & Terminal \(a>0\), malla uniforme en \(u=\log(t/a)\), no en tiempo físico. \\
Caputo--Hadamard & ABM/PECE en tiempo logarítmico & Experimental & Solver conmensurado de historia completa y condición puntual \(x(a)\). \\
Atangana--Baleanu--Caputo (ABC) & Operador muestral con kernel Mittag--Leffler & Experimental con crítica explícita & \(0<q\le1/2\), normalización explícita; directo Numba/Python o FFT. \\
Atangana--Baleanu--Caputo (ABC) & Predictor--corrector convencional & Experimental & Solver conmensurado, \(0<q<1\), compatibilidad inicial e historia completa \(\mathcal O(dN^2)\); no SOE. \\
\bottomrule
\end{longtable}
\end{center}

\subsubsection{Riemann--Liouville, Grünwald--Letnikov y desplazamiento Caputo}

Para \(0<q<1\), la derivada izquierda de Riemann--Liouville es
\[
 {}_aD_t^q x(t)=\frac{d}{dt}
 \left[
   \frac{1}{\Gamma(1-q)}
   \int_a^t (t-\tau)^{-q}x(\tau)\,d\tau
 \right].
\]
La definición GL correspondiente se obtiene como el límite
\[
 {}_aD_t^q x(t)=\lim_{h\to0^+}h^{-q}
 \sum_{k=0}^{\lfloor (t-a)/h\rfloor}
 (-1)^k \binom{q}{k}x(t-kh).
\]
Los pesos discretos se generan sin funciones gamma por paso:
\[
 w_0=1,\qquad
 w_k=\left(1-\frac{q+1}{k}\right)w_{k-1}.
\]
Para regularidad suficiente y \(0<q<1\), el desplazamiento
\[
 {}_a^CD_t^q x(t)={}_aD_t^q\bigl[x(t)-x(a)\bigr]
\]
explica la opción \codepath{caputo_shifted}; no autoriza a interpretar la
convolución cruda como un IVP Caputo. Las referencias clásicas son
Podlubny~\cite{Podlubny1999} y Lubich~\cite{Lubich1986}.

\subsubsection{Templada, de orden variable y conformable}

La ruta RL templada implementa la conjugación
\[
 {}_aD_t^{q,\lambda}x(t)=
 e^{-\lambda(t-a)}{}_aD_t^q
 \left(e^{\lambda(\,\cdot-a)}x\right)(t),
 \qquad \lambda\ge0,
\]
y después la discretiza mediante dos rutas distintas: GL en
\codepath{tempered_grunwald_letnikov_derivative}, sólo para RL, y CQ BDF1/BDF2
en \codepath{tempered_convolution_quadrature}, que además ofrece la definición
Caputo conjugada con su corrección explícita. La definición y el parámetro
\(\lambda\) deben conservarse en los metadatos~\cite{SabzikarEtAl2015}.

La ruta Caputo templada declara, en cambio,
\[
 {}_a^CD_t^{q,\lambda}x(t)=e^{-\lambda(t-a)}{}_a^CD_t^q
 \left(e^{\lambda(\,\cdot-a)}x\right)(t),
 \qquad 0<q\leq1,\quad \lambda\ge0.
\]
Este intervalo corresponde al operador CQ; el solver ABM/PECE templado
disponible más abajo conserva el contrato más estricto \(0<q<1\).
Para \(0<q\leq1\), la relación correcta entre ambas definiciones es
\begin{equation}
 {}_a^CD_t^{q,\lambda}x(t)=
 {}_aD_t^{q,\lambda}
 \left[x(t)-e^{-\lambda(t-a)}x(a)\right].
 \label{eq:tempered-caputo-anchor}
\end{equation}
El ancla Caputo es, por tanto, \(e^{-\lambda(t-a)}x(a)\), no la constante
física \(x(a)\). Una señal constante no es anulada en general cuando
\(\lambda>0\), mientras que el ancla exponencial sí lo es. Para
\(m-1<q<m\) se necesita el jet inicial templado completo
\(c_j=(D+\lambda)^jx(a+)\), \(j=0,\ldots,m-1\); la CQ pública actual se
limita deliberadamente a \(0<q\leq1\), donde basta el valor puntual.

Si \(\delta_p(\zeta)^q=\sum_{k\ge0}\omega_k^{(q,p)}\zeta^k\) es la
CQ BDF1 o BDF2 ordinaria, la conjugación produce
\[
 \delta_p(e^{-\lambda h}\zeta)^q
 =\sum_{k\ge0}e^{-\lambda kh}\omega_k^{(q,p)}\zeta^k.
\]
HAFO aplica estos pesos amortiguados sin materializar exponentes positivos y,
para Caputo, resta dentro del corchete de convolución
\[
 x(a)e^{-\lambda nh}\sum_{k=0}^{n}\omega_k^{(q,p)}.
\]
El corchete completo se multiplica después por \(h^{-q}\).
Esta convención no incluye el término \(-\lambda^q x\). Tampoco sustituye en
malla finita el símbolo anterior por
\([\delta_p(\zeta)/h+\lambda]^q\): son discretizaciones distintas y la segunda
permanece registrada como método planificado independiente
\cite{ChenDeng2015,GuoEtAl2019Tempered,JinLiZhou2017}.

La función \codepath{tempered_fast_multistep_history} añade el Fast Method II
real de Guo et al.~\cite{GuoEtAl2019Tempered} para una \emph{historia ya
muestreada}; no integra un lado derecho y no es un solver FDE. Para evitar
confundir el templado físico con los nodos de cuadratura, en lo que sigue se
denota por \(\sigma\ge0\) el primero y por
\(\lambda_{\mathrm{aux},j}>0\) los segundos. HAFO admite, para
\(0<q<1\), los generadores no templados
\begin{equation}
 \Omega_{\mathrm{FBDF1}}(z)=(1-z)^q,
 \qquad
 \Omega_{\mathrm{GNGF2}}(z)=(1-z)^q
 \left(1+\frac q2(1-z)\right).
 \label{eq:tempered-fast-generators}
\end{equation}
El segundo es el generalized Newton--Gregory de orden dos, no la BDF2
fraccionaria \(\bigl(\tfrac32-2z+\tfrac12z^2\bigr)^q\). Sin embargo, en el
límite entero
\[
 \Omega_{\mathrm{GNGF2}}(z)\big|_{q=1}
 =(1-z)\left(1+\frac{1-z}{2}\right)
 =\frac32-2z+\frac12z^2,
\]
por lo que a \(q=1\) reduce exactamente a BDF2; HAFO evalúa ese extremo con
los pesos locales exactos y sin cuadratura auxiliar. La distinción es
necesaria porque el factor real de la BDF2 fraccionaria cambia de signo sobre
el eje de integración y exigiría fijar y verificar una rama compleja que esta
ruta real no implementa. Los generadores siguen la construcción FLMM de
Lubich~\cite{Lubich1986}.

Si \(\omega_\ell^{(q)}\) denota el peso no templado, el cambio de variable
\(r=h\lambda_{\mathrm{aux}}\) da, para \(\ell\ge1\),
\begin{equation}
 \omega_\ell^{(q)}=
 -\frac{\sin(\pi q)}{\pi}
 \int_0^\infty r^q(1+r)^{-\ell-1}F_\Omega(-r)\,dr
 \simeq \sum_{j=1}^{Q}a_j(1+r_j)^{-\ell-1},
 \qquad r_j=h\lambda_{\mathrm{aux},j}.
 \label{eq:tempered-fast-real-axis}
\end{equation}
Aquí \(F_\Omega(-r)=1\) para FBDF1 y
\(F_\Omega(-r)=1-qr/2\) para GNGF2. El signo de la densidad es
\(-\sin(\pi q)/\pi\), no el positivo. En FBDF1 se comprueba directamente con
la identidad beta--reflexión
\[
 -\frac{\sin(\pi q)}{\pi}
 B(q+1,\ell-q)
 =\frac{\Gamma(\ell-q)}{\Gamma(-q)\Gamma(\ell+1)}
 =(-1)^\ell\binom q\ell,
\]
lo que fija algebraicamente el signo de los pesos. La cuadratura trapezoidal
truncada se construye sobre la recta real logarítmica, siguiendo el marco de
convergencia exponencial de Trefethen--Weideman
\cite{TrefethenWeideman2014}; los \(r_j\) guardados por la API son
adimensionales.

La convolución se separa en una ventana local exacta de retardos
\(0\le\ell\le n_0\) y una historia anterior comprimida. Con la normalización
adimensional de estado usada por HAFO, cada nodo se actualiza mediante
\begin{equation}
 y_m^{(j)}=\frac{e^{-\sigma h}}{1+r_j}
 \left(y_{m-1}^{(j)}+u_{m-1}\right),
 \qquad y_0^{(j)}=0.
 \label{eq:tempered-fast-recurrence}
\end{equation}
El factor \(h\) que aparece en otra normalización del estado se absorbe aquí
en los coeficientes externos. Por ello la recurrencia anterior es exactamente
la que ejecutan los backends Python y Numba y sólo contiene factores
contractivos. Para la definición RL templada se agrega la contribución local
con pesos \(h^{-q}e^{-\sigma h\ell}\omega_\ell^{(q)}\). Para la definición
Caputo conjugada se resta, con pesos exactos y no comprimidos, el ancla
\[
 h^{-q}u_0e^{-\sigma nh}\sum_{\ell=0}^{n}\omega_\ell^{(q)}.
\]
Así, la compresión modifica únicamente la historia RL cruda y no aproxima la
corrección inicial de la Ecuación~\eqref{eq:tempered-caputo-anchor}.

La calibración compara \emph{cada} peso comprimido en los retardos finitos
\(n_0+1,\ldots,N-1\) contra la recurrencia exacta del generador y exige
\[
 E_{1,i}^{\mathrm{rel}}=
 \frac{\sum_\ell|\widehat\omega_{\ell,i}-\omega_{\ell,i}|}
      {\sum_\ell|\omega_{\ell,i}|}
 \le \varepsilon_{\mathrm{rel}}.
\]
El refinamiento automático usa \(Q=65,129,257,\ldots\) hasta satisfacer esa
desigualdad o falla al alcanzar \codepath{max_quadrature_points}; cuando hay
cola comprimida y selección automática, el tope debe ser al menos 65. El corte
\codepath{tail_cutoff} localiza el intervalo de la recta real, pero no se
publica como cota analítica independiente. En aritmética exacta, la cantidad
reportada
\[
 B_i=h^{-q_i}\max_n|u_{n,i}|
 \sum_\ell|\widehat\omega_{\ell,i}-\omega_{\ell,i}|
\]
acota sólo el error del reemplazo de pesos sobre esa malla. No incluye
redondeo, error de discretización FLMM/CQ, error de un solver FDE, constantes
desconocidas de la regla trapezoidal infinita ni una cota separada de colas.
Con \(d\) componentes, \(Q\) nodos y ventana \(n_0\), el lote cuesta
\(\mathcal O\!\left(d(Q+n_0)N\right)\) y su memoria activa, excluidas entrada
y salida, es \(\mathcal O\!\left(d(Q+n_0)\right)\).

Con \(v(t)=e^{\lambda(t-a)}x(t)\), el problema
\({}_a^CD_t^{q,\lambda}x=f(t,x)\) se convierte exactamente en
\[
 {}_a^CD_t^qv(t)=e^{\lambda(t-a)}
 f\!\left(t,e^{-\lambda(t-a)}v(t)\right),\qquad v(a)=x(a).
\]
Li--Deng--Zhao establecen el operador, la formulación integral y un método de
Jacobi~\cite{LiDengZhao2019}. HAFO reutiliza ABM/PECE después de la
transformación; no afirma implementar ese predictor--corrector de Jacobi ni
heredar sus órdenes. El parámetro \(\lambda\), la política de memoria y el
backend quedan en el contrato. La identidad \(\lambda=0\) recupera Caputo y
forma parte de las pruebas.

En la convención de orden variable disponible, cada salida usa
\[
 D_h^{q_n}x_n=h^{-q_n}\sum_{k=0}^{n}w_k(q_n)x_{n-k}.
\]
El orden se evalúa en el instante de salida y todos los pesos de esa fila se
recalculan con \(q_n\). Esta es una convención concreta entre varias posibles y
no debe extrapolarse a cualquier definición de orden variable
\cite{SamkoRoss1993}.

La ruta de trayectoria de orden variable fija, por separado, el Caputo tipo III
de Tavares--Almeida--Torres~\cite{TavaresEtAl2016}:
\[
 {}_a^{C,III}D_t^{\alpha(t)}x(t)=
 \frac{1}{\Gamma(1-\alpha(t))}
 \int_a^t(t-s)^{-\alpha(t)}x'(s)\,ds.
\]
No contiene términos con \(\alpha'(t)\) y no se identifica con los tipos I/II.
El L1 usa el valor corriente \(\alpha_n=\alpha(t_n)\) en todos los pesos de la
salida \(n\); el algoritmo rápido de Fang--Sun--Wang~\cite{FangEtAl2020} queda
como backend futuro, mientras la implementación actual conserva la historia
directa. Para un arranque \(C^1\), la ecuación en el terminal exige
\(f(a,x_0)=0\); un arranque no suave se declara explícitamente y no hereda la
tasa L1 suave.

La derivada conformable de Khalil es local:
\[
 T_q^a x(t)=(t-a)^{1-q}x'(t),\qquad t>a,
\]
cuando \(x\) es diferenciable~\cite{KhalilEtAl2014}. Su presencia permite
comparaciones y modelos que la declaren explícitamente, pero no representa un
kernel hereditario ni comparte la semántica de memoria de Caputo.

Para un problema conmensurado
\[
 T_q^a x(t)=f(t,x(t)),
\]
el cambio de reloj
\[
 \tau=\frac{(t-a)^q}{q},\qquad
 \frac{dx}{d\tau}=f\!\left(a+(q\tau)^{1/q},x\right)
\]
convierte el modelo en una ODE ordinaria. El solver
\codepath{integrate_conformable_rk4} aplica RK4 clásico con paso uniforme en
\(\tau\), devuelve por separado el reloj y el tiempo físico y rechaza órdenes
distintos por componente porque no existiría un reloj único. La malla física es
generalmente no uniforme. Esta transformación es local y no crea memoria
hereditaria; sus trayectorias no deben agregarse a evidencia Caputo, RL, GL o
Hadamard como si compartieran el mismo operador.

\subsubsection{Caputo--Fabrizio, orden distribuido y ABC}

Con una normalización \(M(q)\) declarada, Caputo--Fabrizio usa el kernel
exponencial
\[
 {}_a^{CF}D_t^q x(t)=
 \frac{M(q)}{1-q}\int_a^t x'(\tau)
 \exp\!\left[-\frac{q}{1-q}(t-\tau)\right]d\tau .
\]
\codepath{caputo_fabrizio_derivative} actualiza exactamente la contribución
exponencial por intervalo; \codepath{caputo_fabrizio_derivative_reference}
conserva una suma directa para contraste. Ambas son funciones sobre datos
muestreados, no un solver promovido. La definición original y la crítica de
consistencia se citan juntas~\cite{CaputoFabrizio2015,DiethelmEtAl2020}; no se
oculta esta controversia detrás de una etiqueta de implementación.

Un operador de orden distribuido toma la forma
\[
 \mathcal D_w x(t)=\int_{q_{\min}}^{q_{\max}}
 w(q)D^q x(t)\,dq
 \approx \sum_{r=1}^{R}\nu_rD^{q_r}x(t).
\]
\codepath{distributed_order_gl_derivative} recibe los nodos \(q_r\), los pesos
de cuadratura \(\nu_r\) y la historia temporal. Su costo suma el de las
convoluciones temporales para todos los nodos; el resultado depende de ambas
resoluciones~\cite{CaputoDistributed2001,DiethelmFord2009}.

El solver específico \codepath{integrate_distributed_order_caputo_l1} restringe
la base a Caputo y la medida a masas no negativas. Aplica L1 a cada nodo y
agrega primero los coeficientes en un kernel único, por lo que la historia cuesta
\(\mathcal O(RN+N^2d)\) en lugar de la suma ingenua
\(\mathcal O(RN^2d)\). Una masa en \(q=1\) se trata como backward Euler
exacto. El corrector Picard vectorial y esta agregación son adaptaciones HAFO;
los análisis publicados de cuadratura/L1 para difusión lineal no prueban
estabilidad ni dinámica oculta del flujo no lineal~\cite{HuEtAl2014Distributed,LinXu2007L1}.

La familia ABC reemplaza el kernel exponencial por uno de Mittag--Leffler,
esquemáticamente
\[
 {}_a^{ABC}D_t^q x(t)=\frac{B(q)}{1-q}
 \int_a^t x'(\tau)
 E_q\!\left(-\frac{q}{1-q}(t-\tau)^q\right)d\tau .
\]
Sobre una malla uniforme, HAFO interpola linealmente la historia y usa
\[
 {}_a^{ABC}D_{t_n}^q x\approx
 \frac{B(q)}{1-q}\sum_{j=1}^{n}(x_j-x_{j-1})w_{n-j},
 \qquad
 w_k=\frac{F((k+1)h)-F(kh)}{h},
\]
\[
 F(s)=sE_{q,2}\!\left(-\frac{q}{1-q}s^q\right).
\]
La función \codepath{abc_piecewise_linear_weights} evalúa la serie definitoria
de \(E_{q,2}\) y rechaza falta de convergencia, cancelación excesiva o pesos que
violen positividad/monotonía. La función
\codepath{atangana_baleanu_caputo_derivative} aplica la convolución directa
Numba/Python o FFT en lote. Se conserva la restricción analizada
\(0<q\le1/2\) de Yadav--Pandey--Shukla~\cite{YadavEtAl2019}; la definición
abstracta no se declara restringida por ello.

La normalización siempre aparece en metadatos. El valor predeterminado es
\(B(q)=1\), mientras \codepath{atangana_baleanu_normalization} ofrece de forma
explícita \(B(q)=1-q+q/\Gamma(q)\). La prueba manufacturada de orden \(1/2\)
usa la identidad \(E_{1/2}(-z)=e^{z^2}\operatorname{erfc}(z)\) y las rutas
Numba, Python y FFT se contrastan entre sí. Esto valida un cálculo finito, no
un solver ni una trayectoria ABC.

La definición y la crítica de kernels no singulares se citan juntas
\cite{AtanganaBaleanu2016,DiethelmEtAl2020}. HAFO mantiene tanto el operador
como el solver en estado experimental y conserva esa controversia en el alcance
del resultado; la disponibilidad numérica no resuelve las objeciones teóricas.

El solver \codepath{integrate_abc_predictor_corrector} implementa las
ecuaciones (9)--(14) convencionales de Lee--Kim--Jang
\cite{LeeKimJang2024}. Para \(0<q<1\), usa la ecuación integral
\[
 x(t)=x_0+c_q f(t,x(t))+
 \frac{d_q}{\Gamma(q)}\int_a^t(t-s)^{q-1}f(s,x(s))\,ds,
 \qquad
 c_q=\frac{1-q}{B(q)},\quad d_q=\frac{q}{B(q)}.
\]
Una condición inicial clásica regular debe satisfacer \(f(a,x_0)=0\); HAFO
calcula ese residuo y rechaza un problema incompatible. La publicación supone
un primer valor con error \(\mathcal O(h^2)\), pero no da su algoritmo. HAFO
añade un arranque implícito de producto--trapecio resuelto por punto fijo y
reporta tanto las iteraciones como la falta de convergencia. Esta elección de
arranque no se atribuye al artículo.

La recurrencia publicada y su evidencia numérica son escalares. La aplicación
del mismo balance a cada coordenada de un estado vectorial con orden común es
una extensión de HAFO. La implementación conserva toda la historia,
\(\mathcal O(dN^2)\) de trabajo y \(\mathcal O(dN)\) de almacenamiento; no usa
la variante rápida por sumas de exponenciales. Un refinamiento de malla observa
orden dos en un problema escalar cuadrático suave; ese resultado no se
transfiere a Chua no suave ni permite inferir estabilidad, caos,
atracción u ocultedad de una trayectoria finita.

\subsubsection{Hadamard y Caputo--Hadamard}

Con \(a>0\), el integral de Hadamard usa medida \(ds/s\):
\[
 ({}_a^HI_t^\alpha x)(t)=\frac{1}{\Gamma(\alpha)}
 \int_a^t\left(\log\frac{t}{s}\right)^{\alpha-1}x(s)\frac{ds}{s}.
\]
La derivación natural emplea
\[
 \delta=t\frac{d}{dt},\qquad
 u=\log(t/a),\qquad
 \delta x(t)=\frac{d}{du}x(ae^u).
\]
La variante tipo RL aplica \(\delta\) después del integral de orden \(1-q\);
la variante Caputo--Hadamard aplica el integral a \(\delta x\)
\cite{JaradEtAl2012}. Esta transformación justifica una malla exponencial
\(t_n=ae^{n\bar h}\), pero no identifica Hadamard con RL en tiempo físico.
La discretización y el solver se detallan en la
Sección~\ref{sec:nucleo-numerico-fraccionario}.

\subsubsection{Qué demuestra la verificación actual}

El caso independiente de Wolfram
\codepath{validation/wolfram/cases/gl_fractional_operator_validation.wl}
reconstruye pesos GL/RL, identidades de gamma para monomios, la constante RL y
el límite \(q=1\). El caso
\codepath{validation/wolfram/cases/hadamard_fractional_operator.wl} verifica,
sin leer la implementación de HAFO ni este reporte, la transformación
logarítmica, identidades Hadamard y Caputo--Hadamard, CQ BDF1/BDF2, el límite
\(q=1\) y una integración ABM manufacturada. El resumen retenido en
\codepath{validation/outputs/wolfram/hadamard_fractional_operator/}
marca \codepath{passed=true}; la peor diferencia de operador fue
\(3.019806626980426\times10^{-14}\) y la máxima diferencia de estado ABM fue
\(8.881784197001252\times10^{-16}\).

El caso CQ templado
\codepath{validation/wolfram/cases/tempered_convolution_quadrature.wl}
construye los coeficientes BDF1/BDF2, los pesos amortiguados y la corrección
Caputo mediante recurrencias de alta precisión independientes. El resumen
promovido en
\codepath{validation/outputs/wolfram/tempered_convolution_quadrature_verified/}
aprobó 18 de 18 aserciones. La reconstrucción Python independiente difirió
como máximo \(1.30841518175551\times10^{-14}\), y la comparación obligatoria
con el núcleo público de HAFO como máximo
\(4.44089209850063\times10^{-15}\). Estas cifras sólo demuestran concordancia
algebraico-numérica en las mallas ensayadas.

El caso separado
\codepath{validation/wolfram/cases/tempered_fast_multistep_history.wl}
reconstruye a 80 dígitos, sin importar HAFO, el signo beta--reflexión, ambos
generadores, el límite GNGF2--BDF2, la recurrencia real, las definiciones RL
templada y Caputo conjugada y el ancla exponencial. Aprobó 13 de 13 aserciones;
el máximo residuo interno Wolfram entre historia rápida y directa fue
\(1.72700092683619\times10^{-26}\). El comparador Python independiente tuvo
una diferencia máxima \(8.881784197001252\times10^{-15}\), y la comparación
obligatoria con el núcleo público
\codepath{hidden_attractors.fractional.tempered_fast_multistep_history}
alcanzó \(1.2434497875801753\times10^{-14}\). Los artefactos se conservan en
\codepath{validation/outputs/wolfram/tempered_fast_multistep_history_verified/}.
Estas magnitudes verifican identidades y una malla finita concreta; no son una
cota general de la regla trapezoidal, un teorema de solver FDE ni evidencia de
caos, atracción u ocultedad.

El caso
\codepath{validation/wolfram/cases/atangana_baleanu_operator.wl} reconstruye
independientemente, para \(q=1/2\), el kernel, los pesos y tres señales de
prueba. Su resumen retenido marca \codepath{passed=true} y el comparador de
\codepath{tests/test_atangana_baleanu_wolfram.py} contrasta las rutas Python,
Numba y FFT. Esta evidencia pertenece sólo al operador muestral en malla finita;
no cubre por sí misma el predictor--corrector ABC, su arranque ni su
compatibilidad inicial. Un segundo caso independiente,
\codepath{validation/wolfram/cases/abc_predictor_corrector.wl}, deriva por
integración simbólica los pesos lineales y reconstruye una trayectoria
manufacturada. El resumen retenido da residuo simbólico cero, discrepancia máxima
de pesos Wolfram--Python de \(3.5388358909926865\times10^{-16}\) y discrepancia
de trayectoria de \(2.220446049250313\times10^{-16}\). El error finito
\(4.952013688854784\times10^{-4}\) frente a la solución Volterra se registra
sólo como diagnóstico, no como prueba de convergencia.

El caso
\codepath{validation/wolfram/cases/variable_order_caputo_type3_l1.wl}
deriva independientemente los pesos L1 mediante integración simbólica, la
identidad tipo III para potencias, una recurrencia manufacturada y la reducción
a orden constante. El residuo simbólico fue cero, la máxima diferencia
Wolfram--HAFO fue \(1.5543122344752192\times10^{-15}\) y la diferencia de
trayectoria \(2.220446049250313\times10^{-16}\). Los errores finitos del operador
\(2.0245592\times10^{-2}\) y de la recurrencia
\(1.4766213\times10^{-2}\) son diagnósticos de una malla, no tasas de
convergencia.

Estas validaciones son comparaciones algebraicas y de malla finita
independientes. No
demuestran convergencia general, estabilidad no lineal, caos u ocultedad, ni
se extienden automáticamente a Caputo--Fabrizio, a la ruta conformable, al orden
variable tipo I/II, a otras convenciones de orden variable o al orden
distribuido. Las demás familias se
sostienen por soluciones manufacturadas,
invariantes discretos, paridad de backends y referencias primarias, con su
alcance declarado en cada resultado.


\section{Modelos de Chua implementados}

\subsection{Chua no suave}

El sistema de Chua no suave usado por la librería es
\begin{equation}
\begin{aligned}
\CaputoD x &= \alpha(y-x-f(x)),\\
\CaputoD y &= x-y+z,\\
\CaputoD z &= -\beta y-\gamma z,
\end{aligned}
\label{eq:chua_nonsmooth}
\end{equation}
con
\begin{equation}
f(x)=m_1x+\frac{m_0-m_1}{2}\left(|x+1|-|x-1|\right).
\label{eq:nonsmooth_f}
\end{equation}
Para \(q=1\), \(\CaputoD\) se interpreta como derivada ordinaria. Para
\(0<q<1\), se integra el mismo campo vectorial con memoria de Caputo.

Los parámetros por defecto del proyecto son
\[
\alpha=8.4562,\quad \beta=12.0732,\quad \gamma=0.0052,\quad
m_0=-0.1768,\quad m_1=-1.1468.
\]
Los equilibrios se obtienen imponiendo \(F(E)=0\). Como Caputo anula
constantes, las ecuaciones algebraicas de equilibrio son las mismas para el
caso entero y para el fraccionario:
\[
y^\ast=\frac{\gamma}{\beta+\gamma}x^\ast,\qquad
z^\ast=-\frac{\beta}{\beta+\gamma}x^\ast,\qquad
-\frac{\beta}{\beta+\gamma}x^\ast=f(x^\ast).
\]
La librería los calcula con \codepath{equilibria_nonsmooth}. El Jacobiano
regional es
\[
J(x)=
\begin{bmatrix}
-\alpha(1+f'(x)) & \alpha & 0\\
1 & -1 & 1\\
0 & -\beta & -\gamma
\end{bmatrix},
\qquad
f'(x)=
\begin{cases}
m_0, & |x|<1,\\
m_1, & |x|>1.
\end{cases}
\]
En \(x=\pm1\) el campo no es diferenciable; por tanto el Jacobiano regional no
debe interpretarse en esas superficies.

\subsection{Chua con no linealidad arctan}

El sistema arctan sigue la misma estructura dinámica, pero reemplaza la
característica no lineal por una función suave:
\[
f(x)=a_1x+a_2\arctan(\rho x).
\]
La implementación local usa \codepath{rhs_arctan}, \codepath{equilibria_arctan}
y \codepath{jacobian_arctan}. Este modelo está conectado a Wu et al.
2023~\cite{Wu2023ArctanChua}, pero la reproducción registrada no forma parte de
la evidencia promovida de ocultedad porque no contiene cuencas completas ni
una referencia dinámica promovida.

\subsection{Comparación de no linealidades: trozos frente a arctan}

El artículo arctan reemplaza la característica definida a trozos por una
característica suave. En la convención de Wu et al.\ se escribe
\[
f_{\arctan}(x)=m x+(n-m)\arctan(x),
\qquad m=0.4,\qquad n=-1.1585.
\]
La librería parametriza la misma expresión como
\[
a_1=m,\qquad a_2=n-m=-1.5585,\qquad \rho=1,
\]
de modo que \(f_{\arctan}(x)=a_1x+a_2\arctan(\rho x)\). Para comparar contra
una característica no suave con las mismas pendientes interior y exterior se
usa
\[
f_{\rm pwl}(x)=m x+\frac{n-m}{2}\left(|x+1|-|x-1|\right).
\]
La diferencia visual importante es que \(f_{\rm pwl}\) tiene cambios bruscos
de pendiente en \(x=\pm1\), mientras que \(f_{\arctan}\) conserva una
transición suave: su pendiente
\[
f'_{\arctan}(x)=a_1+\frac{a_2\rho}{1+\rho^2x^2}
\]
coincide con \(n\) cerca de \(x=0\) y tiende a \(m\) cuando \(|x|\) crece.

\begin{figure}[H]
\centering
\reportinclude[width=0.94\textwidth]{chua_nonlinearity_piecewise_vs_arctan.pdf}
\caption{Comparación entre la característica definida a trozos y la característica arctan usada por el caso Wu et al. 2023. La curva arctan suaviza los cambios de pendiente del modelo por tramos.}
\end{figure}

Esta figura reproduce la idea cualitativa de la comparación publicada:
sustituir los quiebres de una característica lineal por tramos por una
función suave con pendientes límite comparables. La figura no valida la
dinámica; sólo explica la no linealidad que entra al campo vectorial.

\subsection{Lectura uniforme de modelos y figuras}

Las figuras de atractores, series, transferencia y continuación se colocan en las secciones de los tres casos Chua. En la parte conceptual sólo se conserva la comparación de no linealidades, porque es una figura pedagógica generada por el script editorial del reporte. Esta separación evita mezclar corridas antiguas con evidencia de los ejemplos activos.
\section{Lur'e, función de transferencia y función descriptiva}

La forma de Lur'e separa el sistema en un bloque lineal y una realimentación
no lineal escalar:
\[
\CaputoD X=P X+b\,\psi(r^\top X),\qquad \sigma=r^\top X.
\]
Para Chua no suave, se toma la pendiente externa como parte lineal y se deja
la saturación residual en \(\psi\). Con \(r={\left(1,0,0\right)}^\top\),
\[
P=
\begin{bmatrix}
-\alpha(1+m_1) & \alpha & 0\\
1 & -1 & 1\\
0 & -\beta & -\gamma
\end{bmatrix},
\qquad
b={\left(-\alpha,0,0\right)}^\top .
\]
La transferencia usada por el balance armónico se define usando la convención explícita:
\[
\widehat W_q\!\left(\lambda\right)=r^\top\!{\left(\lambda I-P\right)}^{-1}b, \qquad W_q\!\left(\ii\omega\right)\equiv \widehat W_q\!\left({\left(\ii\omega\right)}^q\right), \qquad \lambda={\left(\ii\omega\right)}^q.
\]
Si la implementación usa \(W_{\rm code}=r^\top\!{\left(P-\lambda I\right)}^{-1}b\), entonces \(W_{\rm code}=-\widehat W_q\!\left(\lambda\right)\), y la condición residual debe escribirse con el signo correspondiente. Para \(q=1\), esto coincide con la transferencia ordinaria. Para \(0<q<1\), la librería usa la convención local \(\lambda={\left(\ii\omega\right)}^q\). Esta operación es una herramienta de semillas, no una equivalencia exacta con una órbita periódica de Caputo.

La función descriptiva de la saturación residual con \(\psi(\sigma)=(m_0-m_1)\operatorname{sat}(\sigma)\) es
\begin{equation}
N(A)=
\begin{cases}
m_0-m_1, & A\le 1,\\[4pt]
\displaystyle
\frac{2(m_0-m_1)}{\pi}
\left[
\arcsin\left(\frac{1}{A}\right)
+\frac{1}{A}\sqrt{1-\frac{1}{A^2}}
\right], & A>1.
\end{cases}
\label{eq:df_nonsmooth}
\end{equation}
El cruce armónico busca \(A\) y \(\omega\) tales que
\[
1-N(A)\widehat W_q\!\left(\lambda\right)=0, \qquad \lambda={\left(\ii\omega\right)}^q.
\]
La solución produce una semilla inicial para integración y continuación. No
prueba que exista un atractor oculto.

\subsection{Función descriptiva sesgada y cierre continuo}

La ruta no suave corregida admite una oscilación desplazada
\(\sigma=c+A\cos\theta\). Para no confundir el balance de continua con el
primer armónico, define por separado
\begin{align}
\Psi_0(c,A)
  &=\frac{1}{2\pi}\int_0^{2\pi}
    \psi(c+A\cos\theta)\,d\theta,\\
\Psi_1(c,A)
  &=\frac{1}{\pi}\int_0^{2\pi}
    \psi(c+A\cos\theta)\cos\theta\,d\theta,
&
N_1(c,A)&=\frac{\Psi_1(c,A)}{A}.
\label{eq:biased_df}
\end{align}
Si \(\det P\neq0\), la raíz BDF debe cerrar simultáneamente la componente
continua y la armónica:
\begin{equation}
\bar X=-P^{-1}b\Psi_0,\qquad
c=r^\top\bar X,\qquad
1-N_1(c,A)\widehat W_q\!\left({(\ii\omega)}^q\right)=0.
\label{eq:biased_dc_harmonic_closure}
\end{equation}
Este sistema determina una raíz \((A,c,\omega)\) y una semilla sesgada; sigue
siendo un mecanismo de localización. La aceptación dinámica requiere después
integración causal, robustez y sondeos de cuenca bajo contratos separados.

Para arctan, la función descriptiva implementada en la forma de Lur'e es
\[
N_{\arctan}(A)=
a_2\,\frac{2\left(\sqrt{1+\rho^2 A^2}-1\right)}{\rho A^2}.
\]
\section{Integradores y contratos numéricos}

\subsection{Orden entero}

Para \(q=1\), el problema deja de ser fraccionario y se integra como una EDO
ordinaria
\[
\dot{x}=F(t,x),\qquad x(t_n)=x_n,\qquad t_{n+1}=t_n+h .
\]
La librería acepta RK4 y el límite entero \codepath{efork_q1}. Aunque
son métodos de orden entero, se documentan dentro del mismo contrato porque
los ejemplos de Chua entero se usan como referencia de control antes de pasar
a \(0<q<1\). El contrato mínimo sigue siendo: paso \(h\), horizonte \(t_f\),
condición inicial, umbrales de divergencia, transitorio descartado, versión
del integrador y artefactos generados.

\subsubsection{RK4 clásico}

RK4 es el método explícito de Runge-Kutta de cuarto orden usado por
\codepath{hidden_attractors/integrations/rk4.py} y por el despachador general
cuando se solicita \codepath{integrator=\texttt{rk4}} con \(q=1\). Con pendientes no
escaladas:
\begin{align}
k_1 &= F(t_n,x_n),\\
k_2 &= F\!\left(t_n+\frac{h}{2},\,x_n+\frac{h}{2}k_1\right),\\
k_3 &= F\!\left(t_n+\frac{h}{2},\,x_n+\frac{h}{2}k_2\right),\\
k_4 &= F(t_n+h,\,x_n+h k_3),\\
x_{n+1} &= x_n+\frac{h}{6}\left(k_1+2k_2+2k_3+k_4\right).
\label{eq:rk4_formula}
\end{align}
Sus coeficientes de Butcher son:
\begin{center}
\small
\begin{tabular}{ccccc}
\toprule
\(0\) & \(0\) & \(0\) & \(0\) & \(0\) \\
\(\frac{1}{2}\) & \(\frac{1}{2}\) & \(0\) & \(0\) & \(0\) \\
\(\frac{1}{2}\) & \(0\) & \(\frac{1}{2}\) & \(0\) & \(0\) \\
\(1\) & \(0\) & \(0\) & \(1\) & \(0\) \\
\midrule
& \(\frac{1}{6}\) & \(\frac{1}{3}\) & \(\frac{1}{3}\) & \(\frac{1}{6}\) \\
\bottomrule
\end{tabular}
\end{center}
En problemas suaves el error global esperado es \(O(h^4)\). En trayectorias
caóticas esta propiedad no garantiza coincidencia punto a punto a tiempos
largos; solamente indica que el método reproduce correctamente problemas de
orden entero bien condicionados bajo refinamiento de malla.

\subsection{ABM predictor-corrector de Caputo}

El integrador ABM de historia completa se basa en fórmulas predictor-corrector
para ecuaciones diferenciales fraccionarias~\cite{DiethelmFordFreed2002}. De
forma explícita, para \(f_j=f(t_j,x_j)\), el predictor PECE usado es
\begin{align}
x_{n+1}^{P}
  &=x_0+\frac{h^q}{\Gamma(q+1)}
    \sum_{j=0}^{n}b_{j,n+1}f_j,\\
b_{j,n+1}
  &=(n+1-j)^q-(n-j)^q,
\label{eq:abm_predictor}
\end{align}
y el corrector es
\begin{equation}
x_{n+1}=x_0+
\frac{h^q}{\Gamma(q+2)}
\left[
f(t_{n+1},x_{n+1}^{P})+
\sum_{j=0}^{n} a_{j,n+1} f_j
\right],
\label{eq:abm_corrector}
\end{equation}
con
\[
a_{j,n+1}=
\begin{cases}
n^{q+1}-(n-q)(n+1)^q, & j=0,\\
(n-j+2)^{q+1}-2(n-j+1)^{q+1}+(n-j)^{q+1},
&1\le j\le n.
\end{cases}
\]
Las sumas conservan todos los valores desde el tiempo inicial declarado. Por
ello, la implementación directa cuesta \(O(N^2)\) operaciones y \(O(N)\)
almacenamiento para \(N\) pasos. La implementación local
\codepath{caputo_abm_integrate} y los backends de historia completa se usan
como referencia numérica bajo contrato.

\subsubsection{Historia única y reinicio autónomo}

La continuidad del estado no implica por sí sola continuidad de la memoria.
En la cadena no suave de Paper07, los 21 niveles de la homotopía forman un
único problema de Volterra con un parámetro por tramos: la historia de Caputo
se conserva al cambiar de nivel. Su extremo se usa después como condición
inicial de un \emph{nuevo} problema autónomo; la memoria de esa evaluación
comienza en ese nuevo tiempo inicial. En la ruta c590, cada orden y cada
reinicio son problemas de valor inicial independientes con historia completa
desde su propio origen; no se transfiere historia entre órdenes o reinicios.

ABM sí tiene validación propia en la librería. No se compara contra una tabla
publicada de Chua, sino contra problemas fraccionarios con solución exacta o
manufacturada: decaimiento lineal con solución de Mittag-Leffler,
soluciones \(t^m\) forzadas analíticamente y sistemas lineales diagonales. El
estado promovido en los artefactos locales es
\codepath{method_validated_against_exact_solution}.

\subsection{Sondeo de vecindades y clasificador D90}

Los protocolos del reporte distinguen dos geometrías de muestreo. En una bola
tridimensional alrededor del equilibrio \(e_i\),
\[
X_{0,i\ell m}=e_i+r_\ell U_{i\ell m}^{1/3}d_{i\ell m},
\qquad U_{i\ell m}\sim\mathcal U(0,1),\quad
d_{i\ell m}\sim\mathcal U(\mathbb S^2),
\]
donde \(U^{1/3}\) hace uniforme el muestreo respecto del volumen. En una
superficie esférica se usa
\[
X_{0,i\ell m}=e_i+r_\ell d_{i\ell m},\qquad
d_{i\ell m}\in\mathbb S^2.
\]
La auditoría c590 combina los seis semiejes cartesianos con direcciones
deterministas de Fibonacci esférica; la cadena no suave extendida usa bolas
pseudoaleatorias con el generador y la semilla registrados.

Sea \(Y=\{Y_j\}\) la cola postransitoria de una sonda y
\(\mathcal C_h=\{A_k\}\) la nube objetivo finita. Después del submuestreo
declarado, \(\widetilde Y=S_Y(Y)\) y
\(\widetilde{\mathcal C}_h=S_C(\mathcal C_h)\), la regla de contacto es
\begin{align}
D_{90}(\widetilde Y,\widetilde{\mathcal C}_h)
  &=Q_{0.90}\!\left(
    \left\{\min_k\|\widetilde Y_j-\widetilde A_k\|_2\right\}_j
    \right),\\
Y\leadsto_\delta\mathcal C_h
  &\Longleftrightarrow
    D_{90}(\widetilde Y,\widetilde{\mathcal C}_h)\le\delta .
\label{eq:d90_contact}
\end{align}
Es un clasificador numérico dependiente de nube, submuestreo y tolerancia, no
una identidad exacta entre atractores.

Cuando el protocolo busca el primer radio de contacto, procesa los radios en
orden ascendente para todos los equilibrios. Si aparece al menos un contacto,
completa todas las sondas previstas en ese radio y omite los radios mayores.
Así, el primer radio no se selecciona con una corrida incompleta y las filas
posteriores al corte no se mezclan con el experimento declarado.

\subsection{EFORK-3 y backends nativos}

EFORK-3 es un método Runge-Kutta fraccionario usado como referencia
manufacturada y como base de backends rápidos~\cite{Ghoreishi2023}. La
librería separa:
\begin{itemize}
  \item \codepath{efork3_caputo_integrate}: implementación pequeña de
  referencia para validación del método publicado.
  \item \codepath{FractionalChuaBackend}: backend nativo Chua-específico para
  corridas más pesadas.
  \item \codepath{BasinBackend}: clasificador nativo para cortes, bolas y
  cascarones de cuenca.
\end{itemize}

Para una ecuación de Caputo
\[
{}^{C}D_t^q x(t)=F(t,x(t)),\qquad 0<q<1,
\]
EFORK-3 no evalúa simplemente \(F\); evalúa un campo modificado que resta la
historia conocida. Si el paso base es \(n\), se escribe
\[
F_n(t,z)=F(t,z)-H_n(t),
\]
donde, para memoria completa, el término discreto de historia usado por la
implementación de referencia tiene la forma
\[
H_n(t)=
\frac{1}{h\,\Gamma(2-q)}
\sum_{j=s}^{n-1}
(x_{j+1}-x_j)
\left[{\left(t-t_j\right)}^{1-q}-{\left(t-t_{j+1}\right)}^{1-q}\right].
\]
En memoria completa \(s=0\); en memoria truncada \(s\) se mueve a la ventana
activa. Con \(h_q=h^q\), las tres etapas validadas son:
\begin{align}
K_1 &= h_q F_n(t_n,x_n),\\
K_2 &= h_q F_n(t_n+c_2h,\,x_n+a_{21}K_1),\\
K_3 &= h_q F_n(t_n+c_3h,\,x_n+a_{31}K_1+a_{32}K_2),\\
x_{n+1} &= x_n+w_1K_1+w_2K_2+w_3K_3.
\label{eq:efork3_formula}
\end{align}
La forma tabular es
\begin{center}
\small
\begin{tabular}{cccc}
\toprule
0 & \(0\) & \(0\) & \(0\) \\
$c_2$ & $a_{21}$ & 0 & 0 \\
$c_3$ & $a_{31}$ & $a_{32}$ & 0 \\
\midrule
& $w_1$ & $w_2$ & $w_3$ \\
\bottomrule
\end{tabular}
\end{center}
Los coeficientes dependen de \(q\). Definiendo
\[
g_1=\Gamma(1+q),\qquad
g_2=\Gamma(1+2q),\qquad
g_3=\Gamma(1+3q),\qquad
\Delta=2g_2^2-g_3,
\]
la librería calcula:
\begin{align}
c_2 &= {\left(\frac{1}{2g_1}\right)}^{1/q},
&
c_3 &= {\left(\frac{1}{4g_1}\right)}^{1/q},
&
a_{21} &= \frac{1}{2g_1^2},
\\
a_{31} &=
\frac{g_1^2g_2+2g_2^2-g_3}{4g_1^2\Delta},
&
a_{32} &= -\frac{g_2}{4\Delta},
\\
w_1 &=
\frac{8g_1^3g_2^2-6g_1^3g_3+g_2g_3}{g_1g_2g_3},
&
w_2 &=
\frac{2g_1^2(4g_2^2-g_3)}{g_2g_3},
&
w_3 &=
-\frac{8g_1^2\Delta}{g_2g_3}.
\label{eq:efork3_coefficients}
\end{align}
El detalle crítico es la tercera etapa:
\[
K_3\ \text{usa}\ a_{31}K_1+a_{32}K_2.
\]
Ese orden fue revisado porque una adaptación anterior lo había invertido. En
el límite entero \(q=1\), la librería usa la ruta \codepath{efork_q1}; sus
coeficientes racionales son
\[
a_{21}=\frac{1}{2},\qquad
a_{31}=\frac{1}{2},\qquad
a_{32}=-\frac{1}{4},\qquad
w_1=\frac{2}{3},\qquad
w_2=\frac{5}{3},\qquad
w_3=-\frac{4}{3}.
\]

Un backend rápido no se interpreta como evidencia fuerte si no conserva su
contrato de memoria, tolerancias y comparación contra referencia. La
validación publicada de EFORK-3 aplica a la implementación de referencia y al
orden de etapas; las trayectorias nativas Chua-específicas se aceptan por
comparación operacional contra Python, por pruebas de contrato y por
artefactos de validación separados, no por una nueva tabla publicada.

\subsection{Núcleo numérico fraccionario general}
\label{sec:nucleo-numerico-fraccionario}

Las funciones de esta sección se importan desde
\codepath{hidden_attractors.fractional}; no se asume que estén reexportadas por
la raíz \codepath{hidden_attractors}. Los resultados estructurados conservan
definición, método, órdenes, terminal, malla, política de memoria, backend,
fuentes y alcance de evidencia. El opt-in \codepath{allow_experimental=True}
autoriza ejecución, pero no eleva una capacidad experimental a validación
científica.

\subsubsection{Inventario función por función}

\begin{center}
\scriptsize
\begin{longtable}{@{}L{0.23\textwidth}L{0.23\textwidth}L{0.23\textwidth}L{0.23\textwidth}@{}}
\toprule
Función & Objeto y semántica & Costo/backend & Límites, prueba y fuente \\
\midrule
\endfirsthead
\toprule
Función & Objeto y semántica & Costo/backend & Límites, prueba y fuente \\
\midrule
\endhead
\codepath{list_fractional_derivatives},
\codepath{list_fractional_methods} & Enumeran contratos registrados y permiten filtrar por estado. & \(\mathcal O(K)\) sobre un registro pequeño; Python. & Metadatos, no cálculo numérico; \codepath{test_fractional_contracts_and_gl.py}. \\
\addlinespace
\codepath{get_fractional_derivative},
\codepath{get_fractional_method} & Recuperan una definición o método por nombre estable. & \(\mathcal O(1)\); Python. & Rechazan nombres desconocidos; registro ejercitado por las pruebas de contrato y problema. \\
\addlinespace
\codepath{normalize_fractional_orders} & Convierte un orden escalar o uno por componente en un vector de dimensión declarada. & \(\mathcal O(d)\); NumPy. & Valida finitud, intervalo y longitud; \codepath{test_fractional_contracts_and_gl.py}. \\
\addlinespace
\codepath{validate_fractional_method} & Comprueba compatibilidad entre definición, método, orden, memoria y opt-in experimental. & \(\mathcal O(d)\); Python. & \codepath{test_fractional_contracts_and_gl.py}; no es validación de convergencia ni estabilidad. \\
\addlinespace
\codepath{FractionalProblem} & Contrato serializable de derivada, método, orden, estado inicial, paso, intervalo, coordenada de malla, memoria y kernel. & Construcción \(\mathcal O(d)\). & Los operadores muestreados se registran pero no pasan como solvers; \codepath{test_fractional_problem.py}. \\
\addlinespace
\codepath{solve_fractional_problem} & Despacha un \codepath{FractionalProblem} ejecutable sobre un RHS. & Hereda costo y backend del método. & \codepath{test_fractional_problem.py}; devuelve evidencia finita, no clasifica caos u ocultedad. \\
\addlinespace
\codepath{solve_fractional_system} & Adaptador de un sistema HAFO registrado al mismo contrato; puente usado por Toolbox Chaos. & Overhead Python por RHS más el solver seleccionado. & \codepath{test_fractional_problem.py}; mapas y ecuaciones en diferencias requieren otro tipo de problema. \\
\addlinespace
\codepath{get_fractional_reference} & Recupera por clave una referencia primaria o crítica del registro. & \(\mathcal O(1)\); Python. & \codepath{test_fractional_problem.py}; procedencia no valida por sí misma el método. \\
\addlinespace
\codepath{grunwald_letnikov_weights} & Genera \(w_k=(-1)^k\binom{q}{k}\) por recurrencia. & \(\mathcal O(N)\) tiempo y memoria; Numba. & \codepath{test_fractional_contracts_and_gl.py}; Podlubny y Lubich~\cite{Podlubny1999,Lubich1986}. \\
\addlinespace
\codepath{grunwald_letnikov_derivative} & Convolución GL directa sobre muestras, cruda o desplazada. & \(\mathcal O(dN^2)\), memoria \(\mathcal O(dN)\); Numba/Python. & Operador, no FDE; Wolfram sólo verifica GL/RL en malla finita. \\
\addlinespace
\codepath{native_grunwald_letnikov_weights} & Pesos GL mediante la ABI C. & \(\mathcal O(N)\); C. & Paridad con la recurrencia canónica; backend interno opcional. \\
\addlinespace
\codepath{native_grunwald_letnikov_convolution} & Convolución causal de pesos e historia mediante C/OpenMP. & \(\mathcal O(dN^2)\); C/OpenMP según construcción. & Acelera el mismo algoritmo, no cambia su error matemático; \codepath{test_native_grunwald_letnikov.py}. \\
\addlinespace
\codepath{native_grunwald_letnikov_derivative} & Envoltura completa de pesos, convolución y metadatos nativos. & \(\mathcal O(dN^2)\); C/OpenMP. & Fallback y metadatos de compilación explícitos; benchmark de host, no promesa universal. \\
\addlinespace
\codepath{gl_linear_convolution_fft_length} & Calcula el relleno cero mínimo para convolución lineal sin alias circular. & \(\mathcal O(1)\). & Utilidad de tamaño; no calcula una derivada. \\
\addlinespace
\codepath{fast_grunwald_letnikov_derivative} & Convolución GL por FFT en lote. & \(\mathcal O(dN\log N)\), memoria \(\mathcal O(dN)\); NumPy FFT o selección automática. & No es online y puede perder ventaja en mallas cortas; \codepath{test_fast_grunwald_letnikov.py}, Matusiak~\cite{Matusiak2020}. \\
\addlinespace
\codepath{integrate_gl_explicit_numba} & Recurrencia explícita de una FDE discreta GL con RHS compilado. & \(\mathcal O(dN^2+NC_F)\); Numba \codepath{nopython}. & Historia completa y convención inicial explícita; no es ABM Caputo. \\
\addlinespace
\codepath{integrate_gl_explicit} & Envoltura con RHS genérico y selección Numba/Python. & Igual orden asintótico; Python/Numba. & Trayectoria de un modelo discreto declarado; \codepath{test_fractional_contracts_and_gl.py}. \\
\addlinespace
\codepath{lubich_bdf_weights} & Coeficientes de \(\delta(\zeta)^q\) para BDF1 o BDF2. & \(\mathcal O(N)\); Numba. & Sólo \(0<q\le1\); \codepath{test_lubich_convolution_quadrature.py}, Lubich~\cite{Lubich1986,Lubich2004}. \\
\addlinespace
\codepath{lubich_convolution_quadrature} & CQ RL cruda o desplazada Caputo sobre una historia uniforme. & Directa \(\mathcal O(dN^2)\) en Python/Numba; FFT \(\mathcal O(dN\log N)\) en lote. & No es solver FDE; BDF2 sin correcciones de arranque~\cite{JinLiZhou2017}. \\
\addlinespace
\codepath{hadamard_convolution_quadrature} & CQ Hadamard/Caputo--Hadamard tras \(u=\log(t/a)\). & Directa \(\mathcal O(dN^2)\); FFT \(\mathcal O(dN\log N)\); Python/Numba/FFT. & Exige \(a>0\) y uniformidad en \(u\); operador, no solver; prueba Python y comparación Wolfram finita independiente~\cite{JaradEtAl2012,YinEtAl2024}. \\
\addlinespace
\codepath{integrate_caputo_hadamard_abm} & Solver ABM/PECE conmensurado en tiempo logarítmico y RHS en tiempo físico. & \(\mathcal O(dN^2)\), historia \(\mathcal O(dN)\); NumPy o historial C con callback Python. & Malla física no uniforme, historia completa; solución manufacturada y comparación Wolfram finita~\cite{Zheng2021,GreenLiuYan2021}. \\
\addlinespace
\codepath{riemann_liouville_gl_derivative} & Envoltura RL explícita sobre muestras. & \(\mathcal O(dN^2)\); Numba/Python. & Semántica \codepath{operator_only}; no acepta una condición puntual como IVP RL. \\
\addlinespace
\codepath{tempered_grunwald_letnikov_derivative} & RL templada por conjugación exponencial y GL. & \(\mathcal O(dN^2)\); Numba/Python. & No es Caputo templada; \codepath{test_fractional_contracts_and_gl.py}~\cite{SabzikarEtAl2015}. \\
\addlinespace
\codepath{tempered_convolution_quadrature} & CQ BDF1/BDF2 para RL templada o Caputo templada conjugada sobre muestras uniformes. & Directa \(\mathcal O(dN^2)\) en Python/Numba; FFT lineal en lote \(\mathcal O(dN\log N)\). & Operador, no solver ni historia rápida; corrección inicial Caputo exacta y sin correcciones de arranque BDF2~\cite{ChenDeng2015,GuoEtAl2019Tempered,JinLiZhou2017}. \\
\addlinespace
\codepath{tempered_fast_multistep_history} & Operador muestral RL templado o Caputo conjugado con ventana local exacta e historia real recurrente FBDF1/GNGF2. & \(\mathcal O(d(Q+n_0)N)\), memoria activa \(\mathcal O(d(Q+n_0))\); Python/Numba. & Calibración \(L^1\) de todos los pesos comprimidos en la malla finita; no es solver FDE ni BDF2 fraccionaria~\cite{GuoEtAl2019Tempered,Lubich1986}. \\
\addlinespace
\codepath{integrate_tempered_caputo_abm} & Solver Caputo templado por conjugación y ABM/PECE conmensurado. & \(\mathcal O(dN^2+NC_F)\); historia física amortiguada en C con callback RHS o NumPy. & \(\lambda\ge0\) explícito; reducción exacta a Caputo para \(\lambda=0\), historia completa o ventana declarada~\cite{LiDengZhao2019}. \\
\addlinespace
\codepath{variable_order_grunwald_letnikov_derivative} & GL con \(q_n\) fijado en cada salida. & \(\mathcal O(dN^2)\); Numba/Python. & Convención de salida concreta, no universal~\cite{SamkoRoss1993}. \\
\addlinespace
\codepath{variable_order_l1_weight}, \codepath{integrate_variable_order_caputo_type3_l1} & Pesos L1 y solver Caputo tipo III con orden temporal prescrito y corrector Picard. & Historia \(\mathcal O(N^2)\), trabajo \(\mathcal O(dN^2)\), almacenamiento \(\mathcal O(dN)\); Numba/Python. & Sólo \(0<\alpha(t)<1\), historia completa y regularidad declarada; pruebas manufacturadas y Wolfram finito~\cite{TavaresEtAl2016,FangEtAl2020}. \\
\addlinespace
\codepath{conformable_khalil_derivative} & Aproxima \(T_q^a x=(t-a)^{1-q}x'(t)\). & \(\mathcal O(dN)\); Numba/Python. & Operador local sobre malla; no memoria hereditaria~\cite{KhalilEtAl2014}. \\
\addlinespace
\codepath{integrate_conformable_rk4} & Integra un flujo conmensurado en el reloj \(\tau=(t-a)^q/q\). & \(\mathcal O(N(C_F+d))\), salida \(\mathcal O(dN)\); Numba/Python. & RK4 clásico sobre una ODE local reparametrizada; malla física no uniforme y sin memoria hereditaria~\cite{KhalilEtAl2014}. \\
\addlinespace
\codepath{caputo_fabrizio_derivative} & Recurrencia exacta por intervalo para el kernel exponencial. & \(\mathcal O(dN)\); Numba/Python. & Estado de definición \codepath{research_required}; operador, no solver~\cite{CaputoFabrizio2015,DiethelmEtAl2020}. \\
\addlinespace
\codepath{caputo_fabrizio_derivative_reference} & Suma directa del mismo operador como oráculo interno. & \(\mathcal O(dN^2)\); Python. & Referencia de prueba, no ruta recomendada de producción. \\
\addlinespace
\codepath{atangana_baleanu_normalization} & Evalúa la convención explícita \(B(q)=1-q+q/\Gamma(q)\). & \(\mathcal O(1)\); Python. & Opción, no normalización universal; extremos cero y uno documentados. \\
\addlinespace
\codepath{abc_linear_product_weights} & Devuelve los dos pesos lineales del predictor--corrector por retardo, con índice cero sin uso. & \(\mathcal O(N)\), memoria \(\mathcal O(N)\); NumPy. & Sólo \(0<q<1\); ecuaciones (9)--(14) de Lee--Kim--Jang y prueba algebraica dedicada~\cite{LeeKimJang2024}. \\
\addlinespace
\codepath{abc_piecewise_linear_weights} & Construye pesos de intervalo para el kernel \(E_q\) mediante la serie de \(E_{q,2}\). & \(\mathcal O(NK)\), memoria \(\mathcal O(N)\); Numba/Python. & Sólo \(0<q\le1/2\); verifica convergencia, cancelación, positividad y monotonía~\cite{YadavEtAl2019}. \\
\addlinespace
\codepath{atangana_baleanu_caputo_derivative} & Aplica el operador ABC a datos uniformes con interpolación lineal. & Directa \(\mathcal O(dN^2)\) o FFT \(\mathcal O(dN\log N)\), más pesos \(\mathcal O(NK)\). & Experimental con crítica explícita; operador, no solver; Numba/Python/FFT~\cite{AtanganaBaleanu2016,DiethelmEtAl2020}. \\
\addlinespace
\codepath{atangana_baleanu_caputo_derivative_reference} & Ruta directa Python transparente del mismo esquema. & \(\mathcal O(dN^2+NK)\). & Oráculo de implementación; la identidad \(q=1/2\) aporta el contraste matemático independiente. \\
\addlinespace
\codepath{integrate_abc_predictor_corrector} & Solver ABC conmensurado convencional de Lee--Kim--Jang, con compatibilidad y arranque auditados. & \(\mathcal O(dN^2+NC_F)\), historia \(\mathcal O(dN)\); Numba/Python. & Sólo \(0<q<1\) e historia completa; extensión vectorial HAFO, no SOE ni garantía de orden dos para Chua no suave~\cite{LeeKimJang2024,DiethelmEtAl2020}. \\
\addlinespace
\codepath{distributed_order_gl_derivative} & Suma ponderada de operadores GL en nodos de orden. & \(\mathcal O(RdN^2)\); Numba/Python. & Error temporal y de cuadratura en orden; \codepath{test_distributed_order_operator.py}~\cite{CaputoDistributed2001,DiethelmFord2009}. \\
\addlinespace
\codepath{integrate_distributed_order_caputo_l1} & Solver Caputo de orden distribuido con kernel L1 combinado y Picard. & \(\mathcal O(RN+N^2d)\), memoria \(\mathcal O(Nd+N+R)\); Numba/Python. & Medida no negativa, historia completa y errores temporal/de orden separados; \(q=1\) usa backward Euler~\cite{DiethelmFord2009,HuEtAl2014Distributed,LinXu2007L1}. \\
\addlinespace
\codepath{integrate_multi_term_caputo_l1} & Fachada para una suma finita de derivadas Caputo con coeficientes de ecuación no normalizados. & Canonización \(\mathcal O(R\log R)\), seguida por el mismo kernel \(\mathcal O(RN+N^2d)\); Numba/Python. & Medida atómica discreta, coalescencia exacta y metadatos originales; no es cuadratura de una densidad continua~\cite{DiethelmFord2004MultiOrder,RenSun2014MultiTerm,SheLiSun2022TransformedL1,ZakyMachado2020Atomic}. \\
\bottomrule
\end{longtable}
\end{center}

\subsubsection{CQ templada BDF1/BDF2 por conjugación exponencial}

Para \(p\in\{1,2\}\), HAFO usa
\[
 \delta_1(\zeta)=1-\zeta,
 \qquad
 \delta_2(\zeta)=\frac32-2\zeta+\frac12\zeta^2,
 \qquad
 \delta_p(\zeta)^q=\sum_{k\ge0}\omega_k^{(q,p)}\zeta^k.
\]
La conjugación exponencial sustituye \(\zeta\) por
\(e^{-\lambda h}\zeta\), de modo que los pesos realmente aplicados son
\[
 \omega_{k,\lambda}^{(q,p)}=
 e^{-\lambda kh}\omega_k^{(q,p)}.
\]
En \(t_n=a+nh\), la ruta RL templada evalúa
\begin{equation}
 \mathcal D_{h,p}^{q,\lambda}x_n=
 h^{-q}\sum_{k=0}^{n}
 \omega_k^{(q,p)}e^{-\lambda kh}x_{n-k},
 \label{eq:tempered-cq-rl}
\end{equation}
y la ruta Caputo conjugada implementa exactamente
\begin{equation}
 {}^C\mathcal D_{h,p}^{q,\lambda}x_n=
 h^{-q}\left[
 \sum_{k=0}^{n}\omega_k^{(q,p)}e^{-\lambda kh}x_{n-k}
 -x_0e^{-\lambda nh}\sum_{k=0}^{n}\omega_k^{(q,p)}
 \right].
 \label{eq:tempered-cq-caputo}
\end{equation}
La segunda expresión discretiza la Ecuación~\eqref{eq:tempered-caputo-anchor};
no usa el desplazamiento incorrecto \(x-x_0\). Sólo se materializan
exponentes no positivos. Un factor que subdesborda a cero descarta un término
ya irrepresentable, mientras que el estado transformado de crecimiento
\(e^{+\lambda(t-a)}x(t)\) nunca se construye.

La convención es la derivada conjugada no normalizada: no se resta
\(\lambda^q x\). El símbolo discreto es
\(h^{-q}\delta_p(e^{-\lambda h}\zeta)^q\), no
\([\delta_p(\zeta)/h+\lambda]^q\). HAFO registra el segundo como
\codepath{tempered_symbol_shift_cq}, todavía planificado, porque sus pesos,
constantes de error y arranque finitos no son intercambiables
\cite{SabzikarEtAl2015,ChenDeng2015}. BDF2 usa historia truncada desde el
terminal, sin una etapa BDF1 ni correcciones de arranque para datos poco
regulares; por ello no se le atribuyen los órdenes de los esquemas corregidos
de Jin--Li--Zhou~\cite{JinLiZhou2017}.

La API exige declarar si se evalúa RL como operador o Caputo con valor puntual:
\begin{verbatim}
import numpy as np
from hidden_attractors.fractional import (
    TEMPERED_CAPUTO_INITIAL_CONDITION,
    tempered_convolution_quadrature,
)

times = np.linspace(0.7, 1.7, 513)
tau = times - times[0]
orders = np.array([0.58, 0.84])
tempering = np.array([0.35, 1.10])
samples = np.column_stack((
    np.exp(-tempering[0] * tau) * (0.8 + 1.2 * tau**3),
    np.exp(-tempering[1] * tau) * (-0.3 + 0.7 * tau**4),
))

result = tempered_convolution_quadrature(
    samples,
    orders,
    tempering=tempering,
    bdf_order=2,
    definition="tempered_caputo",
    times=times,
    lower_terminal=0.7,
    initial_condition_semantics=TEMPERED_CAPUTO_INITIAL_CONDITION,
    backend="fft",
)
\end{verbatim}
El ejemplo reproducible
\codepath{examples/tempered_convolution_quadrature.py} completa este caso con
el sistema manufacturado no lineal
\[
 {}_a^CD_t^{q_i,\lambda_i}x_i=x_i^2+g_i(t),\qquad
 x_i=e^{-\lambda_i(t-a)}(x_{0,i}+A_i(t-a)^{\beta_i}),
\]
donde \(g_i\) se obtiene de
\[
 e^{-\lambda_i(t-a)}A_i
 \frac{\Gamma(\beta_i+1)}{\Gamma(\beta_i+1-q_i)}
 (t-a)^{\beta_i-q_i}-x_i^2.
\]
Con 512 intervalos, BDF2 y FFT, la ejecución retenida dio errores absolutos
finales \(4.21911\times10^{-6}\) y \(5.52861\times10^{-6}\). Es consistencia
manufacturada finita del
operador, no validación de un solver FDE no lineal ni evidencia de caos,
atracción u ocultedad.

El oráculo independiente
\codepath{validation/wolfram/cases/tempered_convolution_quadrature.wl}
reconstruye los coeficientes BDF y la conjugación sin importar HAFO. La corrida
retenida aprobó 18 de 18 aserciones. La reconstrucción Python independiente
difirió como máximo \(1.30842\times10^{-14}\), y la comparación obligatoria
con la API pública como máximo \(4.44089\times10^{-15}\). Son verificaciones
algebraico-numéricas de malla finita, no teoremas de convergencia o estabilidad.
\begin{flushleft}
\small
Artefactos canónicos:
\codepath{validation/reference_cases/tempered_convolution_quadrature/}\\
Resumen SHA-256:
\texttt{1B7DEADA82462BCE}\allowbreak\texttt{5ED31534A6D8054F}\allowbreak
\texttt{BF7DF95171A879F8}\allowbreak\texttt{3612089411D240E6}.\\
Comparación SHA-256:
\texttt{0F28C92F47EA4ACB}\allowbreak\texttt{E9871157B9C21C71}\allowbreak
\texttt{FBB6E863A0C992D1}\allowbreak\texttt{A83B98D59D5334EB}.
\end{flushleft}

El benchmark reproducible
\codepath{benchmarks/bench_tempered_convolution_quadrature.py} calienta Numba
fuera de la medición y compara Python directo, Numba directo y FFT para RL y
Caputo, BDF1/BDF2 y tres cargas. Numba coincidió exactamente con Python y la
diferencia FFT máxima fue \(1.95\times10^{-13}\). FFT tuvo la menor mediana en
12 de 12 casos, con razones Numba/FFT de \(1.09\)--\(1.11\) para
\((N,d)=(64,2)\), \(1.01\)--\(1.10\) para \((256,3)\) y
\(1.23\)--\(1.30\) para \((768,4)\). Python/Numba varió respectivamente entre
\(2.68\)--\(2.93\), \(45.4\)--\(47.8\) y \(312\)--\(359\).
\begin{flushleft}
\small
Evidencia retenida:
\codepath{validation/reference_cases/tempered_convolution_quadrature/benchmark_20260803.json}
\\
SHA-256:
\texttt{B0CC97F37E1EDA87}\allowbreak\texttt{55049F7CF109542D}\allowbreak
\texttt{727199568371E52F}\allowbreak\texttt{FBA540811AE16E5D}.
\end{flushleft}
Esta medición justifica conservar Numba y FFT para la CQ por lotes, pero no se
transfiere automáticamente al operador recurrente siguiente ni justifica por
sí sola un backend C o Julia.

\subsubsection{Historia rápida multistep templada FBDF1/GNGF2}

\codepath{tempered_fast_multistep_history} implementa el Fast Method II de eje
real de Guo et al.~\cite{GuoEtAl2019Tempered} como \textbf{operador sobre
muestras uniformes}. Recibe la historia completa \(u_n\); no recibe un RHS, no
avanza una FDE y no convierte una trayectoria de entrada en solución de un
modelo fraccionario. Admite órdenes por componente \(0<q_i\le1\), templados
\(\sigma_i\ge0\), las definiciones
\codepath{tempered_riemann_liouville} y \codepath{tempered_caputo}, y los
backends \codepath{python} y \codepath{numba}.

Los dos generadores base son los de la
Ecuación~\eqref{eq:tempered-fast-generators}. Para FBDF1, los pesos locales
satisfacen
\[
 \omega_0=1,\qquad
 \omega_\ell=\frac{\ell-1-q}{\ell}\omega_{\ell-1}.
\]
Para GNGF2 se combina la misma sucesión GL \(g_\ell\) como
\[
 \omega_0=1+\frac q2,\qquad
 \omega_\ell=\left(1+\frac q2\right)g_\ell
              -\frac q2g_{\ell-1},
 \quad
 g_0=1,\qquad g_\ell=\frac{\ell-1-q}{\ell}g_{\ell-1}.
\]
Ésta es GNGF2; no es
\(\bigl(\tfrac32-2z+\tfrac12z^2\bigr)^q\). Sólo cuando \(q=1\) su generador
se vuelve exactamente \(\tfrac32-2z+\tfrac12z^2\), y el código evita entonces
la integral real degenerada y devuelve los pesos BDF2 exactos. Esta decisión
no elimina la CQ BDF2 directa/FFT disponible en
\codepath{tempered_convolution_quadrature}; mantiene dos objetos matemáticos
distintos~\cite{Lubich1986}.

Para \(0<q<1\), la representación de eje real usa el signo
\(-\sin(\pi q)/\pi\) fijado en la
Ecuación~\eqref{eq:tempered-fast-real-axis}. La identidad beta--reflexión
recupera exactamente \((-1)^\ell\binom q\ell\) en FBDF1 y constituye una
prueba algebraica del signo. La API devuelve
\codepath{quadrature_nodes} como \(r_j=h\lambda_{\mathrm{aux},j}\), no como
tasas dimensionales. Para cada nodo y componente actualiza
\[
 y_m^{(j)}=\frac{e^{-\sigma h}}{1+r_j}
 \left(y_{m-1}^{(j)}+u_{m-1}\right),
\]
que coincide con la Ecuación~\eqref{eq:tempered-fast-recurrence}; no se
materializa \(e^{+\sigma t}\). En el instante \(n\), la parte local usa
exactamente \(0\le\ell\le\min(n,n_0)\). Para \(n>n_0\), la parte antigua se
reconstruye con los estados anteriores y los coeficientes
\(h^{-q}e^{-n_0\sigma h}a_j(1+r_j)^{-n_0-1}\). Por tanto,
\[
 \mathcal D_{h,\mathrm{fast}}^{q,\sigma}u_n=
 h^{-q}\sum_{\ell=0}^{\min(n,n_0)}
 e^{-\sigma h\ell}\omega_\ell u_{n-\ell}
 +h^{-q}e^{-n_0\sigma h}\sum_{j=1}^{Q}
 a_j(1+r_j)^{-n_0-1}y_{n-n_0}^{(j)}.
\]
En Caputo conjugada se resta aparte
\(h^{-q}u_0e^{-\sigma nh}\sum_{\ell=0}^{n}\omega_\ell\) con la suma parcial
exacta. La compresión nunca aproxima esa ancla.

La selección automática comienza en \(Q=65\), usa la malla anidada
\(65,129,257,\ldots\) y aumenta \(Q\) hasta que el error relativo \(L^1\) de \emph{todos} los pesos de
retardos \(n_0+1,\ldots,N-1\) no supera
\codepath{relative_tolerance}; una \codepath{quadrature_points} explícita que
no cumple la tolerancia produce error y un tope automático menor que 65 se
rechaza cuando existe cola comprimida. La cota
\codepath{operator_absolute_error_bound} multiplica el error absoluto \(L^1\)
de esos pesos por \(h^{-q_i}\max_n|u_{n,i}|\). Es una cota de compresión en
aritmética exacta para la malla finita: no cubre redondeo, truncamiento
FLMM/CQ, error de solver, convergencia general ni constantes analíticas de la
regla trapezoidal o sus colas. \codepath{tail_cutoff} sólo determina el
intervalo truncado y sus metadatos dicen expresamente que no hay certificado
separado. El costo del lote es \(\mathcal O(d(Q+n_0)N)\); la memoria activa de
historia es \(\mathcal O(d(Q+n_0))\), más \(\mathcal O(dN)\) para la salida.

La llamada mínima conserva definición y semántica inicial explícitas:
\begin{verbatim}
from hidden_attractors.fractional import (
    TEMPERED_CAPUTO_INITIAL_CONDITION,
    tempered_fast_multistep_history,
)

fast = tempered_fast_multistep_history(
    samples,
    orders=[0.48, 0.67, 0.83],
    tempering=[0.10, 0.35, 0.70],
    multistep_method="gngf2",
    definition="tempered_caputo",
    step=0.005,
    initial_condition_semantics=TEMPERED_CAPUTO_INITIAL_CONDITION,
    local_history_steps=50,
    relative_tolerance=1e-9,
    backend="numba",
)
assert fast.compression_tolerance_satisfied
\end{verbatim}

El ejemplo \codepath{examples/tempered_fast_history_chua.py} genera con
DOP853 una trayectoria Chua \(q=1\), uniformemente muestreada, y la usa sólo
como una señal realista de tres canales. Aplica después GNGF2 Caputo templada y
la compara con una convolución directa independiente
\(\mathcal O(dN^2)\). Es \emph{postprocesamiento}: no resuelve un Chua
fraccionario y no aporta evidencia de caos, atracción u ocultedad. Su prueba
dedicada exige 801 muestras finitas, tolerancia de compresión satisfecha y una
diferencia rápida--directa no mayor que \(5\times10^{-10}\).

El oráculo independiente
\codepath{validation/wolfram/cases/tempered_fast_multistep_history.wl}
trabaja a 80 dígitos y aprobó 13/13 aserciones. Verificó el signo, los
coeficientes, el límite entero, las cuatro combinaciones
FBDF1/GNGF2--RL/Caputo, el ancla y la amortiguación de la recurrencia. El mayor
residuo rápido--directo dentro de Wolfram fue
\(1.72700092683619\times10^{-26}\); el comparador independiente en precisión
doble alcanzó \(8.881784197001252\times10^{-15}\), y la comparación contra el
núcleo público HAFO, \(1.2434497875801753\times10^{-14}\). La evidencia
retenida vive en
\codepath{validation/reference_cases/tempered_fast_multistep_history/}.
Se reproduce con:
\begin{verbatim}
python examples/tempered_fast_history_chua.py
python -m pytest tests/test_tempered_fast_history.py \
  tests/test_tempered_fast_history_edges.py \
  tests/test_tempered_fast_history_example.py \
  tests/test_tempered_fast_history_wolfram.py \
  tests/test_tempered_fast_history_benchmark.py -q
\end{verbatim}
Las pruebas cubren API pública, contrato/capability, paridad Python--Numba,
FBDF1 frente a CQ directa, GNGF2 frente a una convolución independiente,
ancla Caputo, límite \(q=1\), ventana puramente local, cota finita y rechazo de
contratos inválidos. Estas verificaciones no demuestran un teorema general de
convergencia ni una trayectoria FDE.

El benchmark vertical
\codepath{benchmarks/bench_tempered_fast_history.py} separó el JIT y la
selección automática de \(Q\), y midió tres repeticiones de las 12 combinaciones
RL/Caputo, FBDF1/GNGF2 y tres tamaños. Seleccionó \(Q=65\) para
\((N,d,n_0)=(128,2,16)\) y \((512,3,32)\), y \(Q=129\) para
\((2048,4,50)\). Python y Numba rápidos coincidieron bit a bit; las máximas
diferencias directa--FFT y rápida--directa fueron
\(5.0004\times10^{-12}\) y \(4.4384\times10^{-8}\), respectivamente. La
segunda quedó dentro de la cota de compresión más un margen de acumulación
flotante registrado por separado.

El modelo analítico de memoria activa del evaluador fue 3392, 5472 y
14016 bytes, frente a 4096, 24576 y 131072 bytes para los dos juegos completos
de pesos. La ventaja de almacenamiento creció hasta \(9.35\times\) en
\(N=2048\), excluidas entrada y salida. Las rutas FFT por lotes obtuvieron la
menor mediana en los doce casos retenidos; esta medición finita no invalida el contrato
recurrente, pues FFT almacena la historia de pesos y no es \emph{streaming}.
No se implementó ni midió un candidato C o Julia con semántica idéntica. La
decisión registrada es
\codepath{insufficient_evidence_to_add_native_c_or_julia}: C se reconsidera
sólo si el perfil HAFO/Toolbox muestra un cuello de botella material, y Julia
sólo como adaptador por lote fijado, nunca por paso.

El artefacto portable se conserva en
\codepath{validation/outputs/benchmarks/tempered_fast_history_backends_20260803_v2.json}.
Su SHA-256 es
\texttt{325D9C4BF2254C6F}\allowbreak\texttt{EA9C295080157345}\allowbreak
\texttt{485B1A5426329A9D}\allowbreak\texttt{E0F7C3D949417D4F}.
El puente del repositorio hermano \emph{Toolbox Chaos}, definido en
\codepath{core/hidden_engine.py}, carga de forma diferida esta API, conserva
el resultado tipado y consulta el catálogo original mediante
\codepath{engine_capability}, sin duplicar estados. Cuatro pruebas focales y
el archivo completo de integración aprobaron juntas 28 de 28 pruebas sobre
ausencia opcional, delegación, capacidad y rechazo explícito de FFT. Es evidencia de integración
de software, no de dinámica fraccionaria.

\subsubsection{Solver Caputo templado por conjugación}

El solver usa la identidad
\[
 {}_a^CD_t^{q,\lambda}x(t)=e^{-\lambda(t-a)}{}_a^CD_t^q
 \left(e^{\lambda(\,\cdot-a)}x\right)(t),
 \qquad v(t)=e^{\lambda(t-a)}x(t),
\]
para resolver
\[
 {}_a^CD_t^qv(t)=e^{\lambda(t-a)}
 f\!\left(t,e^{-\lambda(t-a)}v(t)\right),\qquad v(a)=x(a).
\]
\codepath{integrate_tempered_caputo_abm} usa esa conjugación para derivar los
factores \(e^{-\lambda(t_n-t_j)}\) y los aplica directamente al historial
ABM/PECE. Los kernels Python y C almacenan \(x\), no el estado potencialmente
enorme \(v\).
Li--Deng--Zhao~\cite{LiDengZhao2019} proporcionan la definición, la formulación
de Volterra y un predictor--corrector de Jacobi; la discretización disponible
en HAFO es el ABM amortiguado algebraicamente equivalente a la transformación y
no se etiqueta como su algoritmo de Jacobi.

El paso es uniforme en tiempo físico, \(0<q<1\), \(\lambda\ge0\) y un único
orden común. \codepath{FractionalProblem} exige
\codepath{kernel_parameters=\{"tempering": lambda\}}. La historia completa es
la ruta de referencia; una ventana reinicia en \(t_s\) con el ancla exacta
\(e^{-\lambda(t-t_s)}x(t_s)\) y se marca como modelo numérico aproximado. El
backend C acelera pesos e historia física, aunque el RHS genérico cruza un
callback Python. El criterio de divergencia se aplica durante la recurrencia al
estado físico \(x\), antes de evaluar la siguiente derivada. Como sólo se usan
exponentes no positivos, un horizonte grande no obliga a almacenar \(v\);
\codepath{transformed_states} se omite si su reconstrucción desbordaría
\codepath{float64}.

Una solución manufacturada reproducible es
\[
 x_*(t)=e^{-\lambda\tau}(x_0+\tau^p),\qquad
 f_*(t)=e^{-\lambda\tau}
 \frac{\Gamma(p+1)}{\Gamma(p+1-q)}\tau^{p-q},\qquad \tau=t-a.
\]
La prueba también exige que \(\lambda=0\) reproduzca la trayectoria Caputo.
Ambas son comprobaciones finitas del solver, no evidencia de estabilidad,
caos, atracción ni ocultedad.

\begin{verbatim}
from hidden_attractors.fractional import FractionalProblem
from hidden_attractors.fractional import solve_fractional_problem

problem = FractionalProblem(
    derivative="tempered_caputo",
    method="tempered_caputo_abm_pece_transform",
    orders=0.7,
    initial_state=[1.0],
    step=0.01,
    t_span=(0.0, 1.0),
    memory_policy="full_history",
    kernel_parameters={"tempering": 0.4},
    allow_experimental=True,
)
result = solve_fractional_problem(problem, rhs)
\end{verbatim}

\subsubsection{Solver Caputo de orden variable tipo III por L1}

Para \(t_n=a+nh\) y \(\alpha_n=\alpha(t_n)\), la aproximación implementada es
\[
 \frac{h^{-\alpha_n}}{\Gamma(2-\alpha_n)}
 \sum_{j=0}^{n-1}a_{n-j-1}^{(n)}(X_{j+1}-X_j),\qquad
 a_k^{(n)}=(k+1)^{1-\alpha_n}-k^{1-\alpha_n}.
\]
Al separar el último incremento, \codepath{integrate_variable_order_caputo_type3_l1}
resuelve
\[
 X_n=X_{n-1}-H_n+\Gamma(2-\alpha_n)h^{\alpha_n}f(t_n,X_n)
\]
mediante Picard. La suma histórica directa usa Numba sin \codepath{fastmath} o
una ruta Python; el corrector conserva tolerancias absoluta/relativa, máximo de
iteraciones, residuo y diagnóstico del paso fallido. Una iteración de pulido sólo
se acepta si reduce un residuo que ya satisfacía la tolerancia. No se comprueba
contractividad ni se infiere regularidad de \(\alpha(t)\).

El programa de orden puede tener firma \codepath{alpha(t)},
\codepath{alpha(t, initial_state)} o
\codepath{alpha(t, initial_state, parameters)}. El argumento de estado siempre
es una copia fija de \(X_0\), nunca el estado evolucionado. El perfil completo
solicitado se valida puntualmente en \((0,1)\) antes de integrar y el orden
nominal de \codepath{FractionalProblem} debe coincidir con \(\alpha(a)\). Esta
decisión mantiene el operador como orden prescrito en el tiempo y evita convertir
silenciosamente el modelo en \(\alpha(t,X(t))\).

\begin{verbatim}
problem = FractionalProblem(
    derivative="caputo_variable_type3",
    method="vo_caputo_type3_l1",
    orders=0.55,
    initial_state=[1.0],
    step=0.01,
    t_span=(0.0, 1.0),
    memory_policy="full_history",
    kernel_parameters={
        "order_function": alpha,
        "order_function_name": "alpha(t)=0.55+0.2*t",
    },
    method_options={"initial_regularity": "smooth"},
    allow_experimental=True,
)
result = solve_fractional_problem(problem, rhs, parameters)
\end{verbatim}

El oráculo Wolfram independiente deriva los pesos por \codepath{Integrate} y
reconstruye una recurrencia sin leer HAFO. Retiene residuo simbólico cero,
diferencia máxima Wolfram--HAFO de
\(1.5543122344752192\times10^{-15}\) y diferencia de trayectoria de
\(2.220446049250313\times10^{-16}\). Los errores finitos del operador y de la
recurrencia son diagnósticos de malla; no prueban una tasa global, estabilidad,
caos, atracción ni ocultedad.

\subsubsection{Solver Caputo de orden distribuido por kernel L1 combinado}

Para nodos \(0<\alpha_r\le1\) y masas efectivas \(\Omega_r\ge0\), el modelo
discreto aproxima
\[
 \sum_{r=1}^{R}\Omega_r\,{}_a^CD_t^{\alpha_r}X(t)=f(t,X(t)).
\]
Con
\[
 \rho_r=\frac{\Omega_rh^{-\alpha_r}}{\Gamma(2-\alpha_r)},\qquad
 K_k=\sum_r\rho_r\left[(k+1)^{1-\alpha_r}-k^{1-\alpha_r}\right],
\]
la ecuación de paso es
\[
 \sum_{k=0}^{n-1}K_k\Delta X_{n-k}=f(t_n,X_n),\qquad K_0>0.
\]
HAFO precomputa sólo \(K_k\) y resuelve
\[
 X_n=X_{n-1}-K_0^{-1}\sum_{k=1}^{n-1}K_k\Delta X_{n-k}
      +K_0^{-1}f(t_n,X_n)
\]
por Picard. Así evita un tensor orden--tiempo--estado y reduce la estructura
ingenua \(\mathcal O(RN^2d)\) a \(\mathcal O(RN+N^2d)\). Los metadatos
separan Numba solicitado, intentado, realmente usado y fallback; un estado no
finito fallido no se incorpora al prefijo aceptado.

La API distingue masas directas de pesos de cuadratura multiplicados por una
densidad explícita. Rechaza medidas firmadas, no normaliza por defecto y conserva
masa cruda/efectiva. El átomo \(\alpha=1\) usa exactamente
\(K_0=\Omega/h\), \(K_{k>0}=0\), es decir, backward Euler. En presencia de ese
átomo no se exige \(f(a,X_0)=0\); sin él, el chequeo sólo genera advertencia
cuando se declara un arranque \codepath{smooth}.

\begin{verbatim}
problem = FractionalProblem(
    derivative="caputo_distributed_order",
    method="distributed_order_caputo_l1",
    orders=[0.35, 0.80],
    initial_state=[1.0],
    step=0.01,
    t_span=(0.0, 1.0),
    memory_policy="full_history",
    kernel_parameters={
        "order_weights": [0.40, 0.60],
        "weight_semantics": "nonnegative_mass",
        "order_quadrature_name": "two_atom_demo",
    },
    method_options={"initial_regularity": "smooth"},
    allow_experimental=True,
)
result = solve_fractional_problem(problem, rhs)
\end{verbatim}

Aquí \codepath{orders} son los nodos de la medida, independientes de la dimensión,
y el contrato usa \codepath{order_mode=distributed}; no existe un orden nominal
ignorado. Las fuentes sustentan el marco de cuadratura y L1, principalmente en
problemas lineales~\cite{CaputoDistributed2001,DiethelmFord2009,HuEtAl2014Distributed,LinXu2007L1}.
El kernel combinado y Picard son adaptaciones HAFO. La trayectoria finita no
demuestra convergencia general, estabilidad, caos, atracción ni ocultedad.

El oráculo Wolfram independiente construye cada peso mediante
\codepath{Integrate}, agrega tres nodos y resuelve una recurrencia lineal afín
sin leer HAFO. La corrida focal aprobó seis pruebas: residuo simbólico, residuo
de recurrencia y error afín iguales a cero en el artefacto; la diferencia máxima
Wolfram--HAFO fue \(1.5543122344752192\times10^{-15}\), con tolerancia
\(8\times10^{-12}\). Esta identidad finita no constituye una tasa de
convergencia.

\subsubsection{Fachada Caputo multitérmino como medida atómica}

La ecuación
\[
 \sum_{j=1}^{R}c_j\,{}_a^CD_t^{\alpha_j}x(t)=f(t,x(t)),
 \qquad 0<\alpha_j\le1,\quad c_j\ge0,
\]
corresponde exactamente a la medida
\(
\mu=\sum_jc_j\delta_{\alpha_j}
\).
Por ello, \codepath{integrate_multi_term_caputo_l1} reutiliza el solver anterior
con \codepath{nonnegative_mass} y \codepath{normalization=none}; no interpreta
\(c_j\) como pesos de una cuadratura continua ni infiere una densidad. Esta
distinción concuerda con las formulaciones multiorden y multitérmino, el uso de
L1 término a término y la reducción por deltas de Dirac
~\cite{DiethelmFord2004MultiOrder,RenSun2014MultiTerm,SheLiSun2022TransformedL1,ZakyMachado2020Atomic}.

Antes de delegar, la fachada valida todos los órdenes y coeficientes, elimina
ceros sólo bajo una política explícita, ordena y agrupa únicamente órdenes
\codepath{float64} exactamente iguales. Para un orden canónico
\(\widehat\alpha_r\), calcula
\[
 \widehat c_r=\operatorname{fsum}
 \{c_j:\alpha_j=\widehat\alpha_r\}.
\]
No usa tolerancia: fusionar órdenes cercanos alteraría el kernel completo. La
salida conserva órdenes y coeficientes originales, grupos de índices, términos
nulos, suma de coeficientes y términos canónicos. En particular, una suma
\(\sum_jc_j\ne1\) permanece sin normalizar y no pierde sus unidades.

\begin{verbatim}
from hidden_attractors.fractional import integrate_multi_term_caputo_l1

result = integrate_multi_term_caputo_l1(
    rhs,
    initial_state=[0.8],
    orders=[1/3, 2/3, 2/3, 0.55, 1.0],
    coefficients=[0.4, 0.2, 0.5, 0.0, 0.75],
    step=0.01,
    n_steps=100,
    initial_regularity="nonsmooth",
)
\end{verbatim}

Los términos canónicos son
\((1/3,0.4),(2/3,0.7),(1,0.75)\), cuya masa es \(1.85\), no uno. El átomo
\(\alpha=1\) conserva la rama backward Euler. La fachada añade
\(\mathcal O(R\log R)\) una sola vez y no reconstruye el kernel combinado ni
la historia \(\mathcal O(N^2d)\). Por ello no se añadió un segundo backend C o
Julia: el trabajo costoso ya está compilado con Numba, mientras duplicar el
solver aumentaría el contrato de paridad sin acelerar esta capa semántica.

El benchmark del 3 de agosto de 2026 calentó Numba y realizó siete repeticiones
medidas en Python 3.14.3, NumPy 2.4.5, Numba 0.65.1 y un host de 16 hilos
lógicos. Para \(R=128\) términos reducidos exactamente a ocho, \(N=256\) y
dimensión tres, la fachada y el solver canónico produjeron tiempos y estados
idénticos; la razón de medianas fachada/directo fue \(1.0120\). Al aislar la
construcción compilada de 4096 lags, la razón no coalescido/coalescido fue
\(15.0220\); la historia sobre un kernel ya combinado dio \(0.9909\), pues ya
no depende de \(R\). Son mediciones de ese host y esa carga, no aceleraciones
universales. El JSON de trabajo tuvo SHA-256
\codepath{32ac0c43b388ce9664ce4b9cd49f936bfd2e83d704bf463c8c39c6ca219438c9}.

SciSpace se consultó con dos preguntas completas: una sobre ecuaciones Caputo
multitérmino y métodos L1/predictor--corrector, y otra sobre la equivalencia con
una medida atómica frente al orden distribuido continuo. Se conservaron los
identificadores \codepath{3tp7pod1yv}, \codepath{30zlukzu},
\codepath{1o5cvoa44t}, \codepath{3hf80rfehx} y
\codepath{3ued1yt7yp}; la columna \codepath{methods_used} se obtuvo para
\(2/2\) y \(3/3\), respectivamente. La canonización, la política de ceros y el
wrapper estructurado son adaptaciones HAFO, no resultados atribuidos a esos
artículos.

El caso Wolfram independiente usa coeficientes racionales cuya suma es
\(37/20\), incluye \(\alpha=1\), deriva los pesos por integración y verifica
permutación, coalescencia y una recurrencia afín sin importar HAFO. El ejemplo
\codepath{multi_term_caputo_relaxation.py} usa otra solución afín manufacturada.
La ejecución viva marcó \codepath{passed=true}: diferencia máxima
Wolfram--HAFO de \(3.552713678800501\times10^{-15}\) frente a tolerancia
\(8\times10^{-12}\), diferencia máxima de kernel del mismo orden, diferencia
de trayectoria de \(2.220446049250313\times10^{-15}\) e identidad cuadrática
con diferencia cero. La suite focal aprobó 31 pruebas y omitió una prueba
opcional porque no se retuvo el artefacto temporal; la corrida usó y eliminó un
directorio nuevo fuera del repositorio.
Ambos son controles finitos de implementación; no demuestran una tasa general,
contractividad de Picard, estabilidad de un flujo caótico, atracción u
ocultedad.

\subsubsection{Solver conformable en su reloj natural}

Para la derivada conformable desplazada,
\[
 T_{q,a}x(t)=(t-a)^{1-q}x'(t)=f(t,x(t)),
\]
la transformación
\[
 \tau=\frac{(t-a)^q}{q},\qquad
 t=a+(q\tau)^{1/q}
\]
produce la ODE \(dx/d\tau=f(a+(q\tau)^{1/q},x)\).
\codepath{integrate_conformable_rk4} usa RK4 clásico en una malla uniforme de
\(\tau\) y devuelve \codepath{clock_times} junto con los tiempos físicos. El
backend Numba requiere el ABI numérico de RHS; el fallback Python admite un RHS
genérico y ambos reportan su identidad exacta. \codepath{FractionalProblem}
exige \codepath{grid_coordinate=conformable_clock}, memoria
\codepath{none} y un orden común. Esta ruta es una reparametrización local, no
una aproximación de memoria Caputo, RL, GL o Hadamard.

\subsubsection{Operador Atangana--Baleanu--Caputo sobre muestras}

Con \(\lambda=q/(1-q)\), la integral exacta del kernel de Mittag--Leffler
sobre cada intervalo lineal usa
\[
 \int E_q(-\lambda s^q)\,ds
 =sE_{q,2}(-\lambda s^q).
\]
Así, la convolución causal usa los incrementos \(\Delta x_j\), no pesos ABM
de Caputo ni la recurrencia exponencial de Caputo--Fabrizio. Un uso reproducible
es:
\begin{verbatim}
import numpy as np
from hidden_attractors.fractional import (
    atangana_baleanu_caputo_derivative,
    atangana_baleanu_normalization,
)

t = np.linspace(0.0, 1.0, 1025)
x = 1.0 + t**3
abc = atangana_baleanu_caputo_derivative(
    x, t[1] - t[0], 0.5,
    normalization=atangana_baleanu_normalization,
    normalization_name="B(q)=1-q+q/Gamma(q)",
    backend="fft",
)
\end{verbatim}
El backend FFT es una evaluación en lote y no se selecciona automáticamente
sin un benchmark específico. \codepath{tests/test_atangana_baleanu_operator.py}
comprueba constante, linealidad, identidad cerrada de orden \(1/2\), pesos y
paridad Numba--Python--FFT. Este operador conserva el dominio numérico auditado
\(0<q\le1/2\). No resuelve una FDE y su backend FFT no se reutiliza como si
fuera el predictor--corrector o una historia online.

El artefacto Wolfram retenido en
\codepath{validation/outputs/wolfram/atangana_baleanu_operator/} reconstruye
independientemente el caso \(q=1/2\) sin leer HAFO y marca
\codepath{passed=true}. \codepath{tests/test_atangana_baleanu_wolfram.py}
contrasta después las tres rutas públicas. Esta es evidencia del operador en
una malla finita, no del solver ABC, su arranque, compatibilidad o convergencia.

\subsubsection{Predictor--corrector ABC convencional}

Para \(0<q<1\), \codepath{integrate_abc_predictor_corrector} implementa la
recurrencia de historia completa de las ecuaciones (9)--(14) de
Lee--Kim--Jang~\cite{LeeKimJang2024} para
\[
 x(t)=x_0+c_q f(t,x(t))+
 \frac{d_q}{\Gamma(q)}\int_a^t(t-s)^{q-1}f(s,x(s))\,ds,
 \quad c_q=\frac{1-q}{B(q)},\quad d_q=\frac{q}{B(q)}.
\]
La normalización \(B(q)>0\) es un parámetro explícito escalar o funcional; si no
se suministra otra convención, HAFO registra \(B(q)=1\). El valor y la
descripción evaluados se guardan en el resultado. Una condición inicial clásica
regular debe cumplir \(f(a,x_0)=0\). El solver mide ese residuo antes de integrar
y no oculta una incompatibilidad mediante una tolerancia no registrada.

El artículo supone un valor \(x_1\) con error \(\mathcal O(h^2)\), pero no
prescribe su construcción. HAFO resuelve para \(x_1\) la ecuación implícita de
producto--trapecio del primer intervalo por punto fijo; registra tolerancia,
máximo e iteraciones, y expone \codepath{startup_no_convergence} si falla. Este
arranque es una decisión de HAFO. La evidencia publicada es escalar; aplicar el
mismo balance por coordenada a un estado vectorial con un solo orden es también
una extensión de HAFO.

La ruta convencional cuesta \(\mathcal O(dN^2)\), conserva
\(\mathcal O(dN)\) valores y declara \codepath{fast_soe_used=False}. No
implementa la aproximación rápida por suma de exponenciales del mismo artículo
ni admite
\codepath{fast_history}. Los backends Numba y Python comparten los pesos y el
contrato de parada. Un refinamiento de malla observa orden dos en un problema
escalar cuadrático suave; las soluciones manufacturadas y la paridad de backends
verifican casos finitos, pero no garantizan ese orden para el campo no suave de Chua,
ni estabilidad general, caos, atracción u ocultedad. La definición se mantiene
experimental y se presenta junto con la crítica de kernels no singulares de
Diethelm et al.~\cite{DiethelmEtAl2020}.

Un problema compatible mínimo puede escribirse como:
\begin{verbatim}
import numpy as np
from hidden_attractors.fractional import integrate_abc_predictor_corrector

def rhs(time, state):
    del state
    return np.array([time**2])  # f(0, x0) = 0

result = integrate_abc_predictor_corrector(
    rhs,
    initial_state=np.array([1.0]),
    order=0.7,
    lower_terminal=0.0,
    step=0.01,
    n_steps=100,
    normalization=1.0,
    normalization_name="B(q)=1",
    use_acceleration=False,
)
\end{verbatim}
La trayectoria producida sólo demuestra la ejecución declarada del IVP.

La validación independiente en
\codepath{validation/wolfram/cases/abc_predictor_corrector.wl} reconstruye las
integrales de las bases lineales y una recurrencia manufacturada sin importar
HAFO. El artefacto retenido reporta residuo simbólico cero, discrepancia máxima
Wolfram--Python de pesos \(3.5388358909926865\times10^{-16}\) y discrepancia de
trayectoria \(2.220446049250313\times10^{-16}\). El error finito
\(4.952013688854784\times10^{-4}\) frente a la solución Volterra exacta es un
diagnóstico de discretización; no certifica convergencia general, estabilidad,
caos, atracción ni ocultedad.

\subsubsection{CQ de Lubich BDF1/BDF2}

La función generadora del método multipaso define
\begin{equation}
 \delta(\zeta)^q=\sum_{k=0}^{\infty}\omega_k\zeta^k,
 \qquad
 D_h^q x_n=h^{-q}\sum_{k=0}^{n}\omega_kx_{n-k}.
 \label{eq:lubich-cq}
\end{equation}
HAFO admite
\[
 \delta_1(\zeta)=1-\zeta,
 \qquad
 \delta_2(\zeta)=\frac32-2\zeta+\frac12\zeta^2.
\]
BDF1 coincide exactamente con los pesos GL canónicos. En BDF2, los pesos se
obtienen en \(\mathcal O(N)\) de la identidad
\(\delta W'=q\delta'W\). La convención
\codepath{riemann_liouville} aplica la convolución cruda y exige
\codepath{operator_only_no_ivp}; \codepath{caputo_shifted} opera sobre
\(x-x(a)\) y exige \codepath{point_value_shift_x_minus_x0}.

Las rutas directa Python/Numba cuestan \(\mathcal O(dN^2)\). La ruta FFT usa
relleno cero, calcula una convolución lineal en \(\mathcal O(dN\log N)\) y sólo sirve cuando toda
la historia ya está disponible. BDF2 usa historia truncada en el terminal, sin
prehistoria inventada ni correcciones iniciales. Por ello HAFO no hereda los
órdenes corregidos para datos de baja regularidad analizados en
\cite{JinLiZhou2017}. El test dedicado verifica BDF1=GL, expansión BDF2 con alta
precisión, límite entero, monomios, órdenes por componente, mallas y paridad
Python--Numba--FFT.

Ejemplo sobre muestras uniformes de una coordenada de Chua:
\begin{verbatim}
from hidden_attractors.fractional import (
    CAPUTO_SHIFTED_INITIAL_CONDITION,
    lubich_convolution_quadrature,
)

cq = lubich_convolution_quadrature(
    trajectory[:, 1],
    0.92,
    bdf_order=2,
    definition="caputo_shifted",
    times=trajectory[:, 0],
    lower_terminal=float(trajectory[0, 0]),
    initial_condition_semantics=CAPUTO_SHIFTED_INITIAL_CONDITION,
    backend="fft",
)
\end{verbatim}
Este código calcula un operador sobre la señal de Chua; no vuelve a integrar el
sistema ni prueba que la señal sea caótica.

\subsubsection{CQ Hadamard y Caputo--Hadamard}

Con \(u_n=n\bar h\) y \(t_n=ae^{u_n}\), HAFO aplica la
Ecuación~\eqref{eq:lubich-cq} a la señal transformada \(y(u)=x(ae^u)\). La ruta
Hadamard tipo RL opera sobre \(x\); la ruta Caputo--Hadamard opera sobre
\(x-x(a)\). La salida incluye tanto \codepath{times} físicos como
\codepath{log_times}, pesos, transformación, semántica inicial, costo y fuentes.
Yin et al. implementan y analizan órdenes BDF mayores y correcciones para
fuentes singulares~\cite{YinEtAl2024}; HAFO conserva sólo BDF1/BDF2 sin dichas
correcciones.

\begin{verbatim}
from hidden_attractors.fractional import (
    CAPUTO_HADAMARD_INITIAL_CONDITION,
    hadamard_convolution_quadrature,
)

operator = hadamard_convolution_quadrature(
    samples,
    orders=(0.75, 0.9),
    bdf_order=2,
    definition="caputo_hadamard",
    log_step=0.01,
    lower_terminal=1.0,
    initial_condition_semantics=CAPUTO_HADAMARD_INITIAL_CONDITION,
    backend="numba",
)
\end{verbatim}

La comparación independiente en
\codepath{validation/outputs/wolfram/hadamard_fractional_operator/}, archivo
\codepath{hadamard_python_comparison.json}, recalcula la transformación
logarítmica, identidades analíticas y CQ
BDF1/BDF2 sin importar HAFO. El artefacto retenido marca
\codepath{passed=true} y una diferencia numérica máxima de
\(3.019806626980426\times10^{-14}\). Es una comprobación finita del operador y
de la implementación de la malla, no una cota de error uniforme ni un
resultado de estabilidad.

\subsubsection{Solver ABM/PECE Caputo--Hadamard en tiempo logarítmico}

El IVP
\[
 {}_a^{CH}D_t^q x(t)=f(t,x(t)),\qquad x(a)=x_0,
\]
se transforma exactamente mediante~\cite{Zheng2021}
\[
 {}_0^CD_u^q y(u)=f(ae^u,y(u)),
 \qquad y(u)=x(ae^u),\quad u=\log(t/a).
\]
\codepath{integrate_caputo_hadamard_abm} reutiliza el ABM/PECE canónico sobre
una malla uniforme en \(u\). Admite actualmente \(0<q<1\) conmensurado,
condición puntual, historia completa y \codepath{log_step}. El RHS siempre
recibe tiempo físico. El backend acelerado usa el historial ABM en C, pero el
cambio \(u\mapsto ae^u\) y un RHS genérico cruzan un callback Python; se registra
como tal y no se anuncia como campo puramente C.

Ejemplo de uso del campo real Chua no suave bajo esta definición distinta:
\begin{verbatim}
import numpy as np
from hidden_attractors import chua_nonsmooth_parameters, rhs_nonsmooth
from hidden_attractors.fractional import integrate_caputo_hadamard_abm

parameters = chua_nonsmooth_parameters()

def physical_rhs(time, state, p):
    del time
    return rhs_nonsmooth(state, p)

result = integrate_caputo_hadamard_abm(
    physical_rhs,
    initial_state=np.array([0.1, 0.0, 0.0]),
    order=0.92,
    parameters=parameters,
    lower_terminal=1.0,
    upper_terminal=float(np.exp(2.0)),
    log_step=0.002,
    use_acceleration=True,
)
\end{verbatim}
La trayectoria resultante es una demostración finita de la llamada. No reproduce
los casos Caputo del reporte, no valida caos y no aporta evidencia de ocultedad.
El test del solver usa forzamiento constante, una potencia manufacturada de
\(\log(t/a)\), equivalencia con Caputo transformado, tiempo físico del RHS y
paridad C--Python. Las mallas graduadas de Green--Liu--Yan
\cite{GreenLiuYan2021} permanecen pendientes. El mismo caso Wolfram compara un
paso ABM manufacturado y retiene una diferencia máxima de estado de
\(8.881784197001252\times10^{-16}\). Esta paridad de precisión finita tampoco
demuestra convergencia general, estabilidad, caos o existencia de atractores.

\subsubsection{Contrato unificado para HAFO y Toolbox Chaos}

\begin{verbatim}
from hidden_attractors.fractional import (
    FractionalProblem,
    solve_fractional_system,
)

problem = FractionalProblem(
    derivative="caputo_hadamard",
    method="caputo_hadamard_abm_pece",
    orders=0.92,
    initial_state=[0.1, 0.0, 0.0],
    step=0.002,
    t_span=(1.0, 7.38905609893065),
    grid_coordinate="log_t_over_lower_terminal",
    memory_policy="full_history",
    lower_terminal=1.0,
    allow_experimental=True,
)
result = solve_fractional_system(problem, "chua-nonsmooth")
\end{verbatim}
El adaptador preserva el mismo resultado estructurado que un RHS directo. Un
método cuyo contrato dice \codepath{sampled_operator} se rechaza en este
despacho: registrarlo no lo convierte en integrador.

\subsubsection{Rendimiento, C, Numba y Julia}

C/OpenMP se reserva para convoluciones e historiales con bucles estables y ABI
numérica simple. Numba conserva RHS definidos por la persona usuaria y rutas
verificables sin cruzar un intérprete por paso. FFT se usa sólo para operadores
en lote suficientemente largos. Julia puede actuar como extensión opcional por
trayectoria o como oráculo cruzado para catálogos especializados; una llamada
Julia en cada paso destruiría la ventaja de compilación y no forma parte del
número base. Cualquier razón de aceleración citada por los benchmarks es una
medición del host, compilador, tamaño de malla y calentamiento registrados, no
una superioridad universal.

\subsubsection{Comandos de prueba del núcleo}

\begin{verbatim}
cd version_2
python -m pytest tests/test_fractional_contracts_and_gl.py \
  tests/test_fast_grunwald_letnikov.py \
  tests/test_native_grunwald_letnikov.py \
  tests/test_lubich_convolution_quadrature.py \
  tests/test_hadamard_convolution_quadrature.py \
  tests/test_caputo_hadamard_solver.py \
  tests/test_hadamard_wolfram.py \
  tests/test_atangana_baleanu_operator.py \
  tests/test_abc_predictor_corrector.py \
  tests/test_abc_predictor_corrector_wolfram.py \
  tests/test_tempered_caputo_solver.py \
  tests/test_tempered_caputo_regressions.py \
  tests/test_tempered_convolution_quadrature.py \
  tests/test_tempered_convolution_quadrature_example.py \
  tests/test_tempered_convolution_quadrature_wolfram.py \
  tests/test_tempered_convolution_quadrature_benchmark.py \
  tests/test_variable_order_caputo_type3_solver.py \
  tests/test_variable_order_caputo_type3_wolfram.py \
  tests/test_distributed_order_caputo_solver.py \
  tests/test_distributed_order_caputo_wolfram.py \
  tests/test_conformable_solver.py \
  tests/test_distributed_order_operator.py -q
python benchmarks/bench_fractional_gl_kernels.py
python benchmarks/bench_tempered_convolution_quadrature.py
\end{verbatim}
Estas pruebas cubren contratos, fórmulas discretas, soluciones manufacturadas y
paridad de backends. No son una campaña de dinámica no lineal ni una
certificación de atractores.


\subsection{Cómo se validaron ABM, RK4 y EFORK-3}

La validación de integradores responde una pregunta limitada: si el código del
método integra correctamente problemas con solución conocida o reproduce una
referencia publicada. No certifica que una trayectoria sea caótica, no prueba
ocultamiento y no reemplaza los análisis de cuenca o robustez.

\begin{center}
\footnotesize
\begin{longtable}{@{}L{0.15\textwidth}L{0.25\textwidth}L{0.28\textwidth}L{0.22\textwidth}@{}}
\toprule
Código & Qué se comparó & Publicación o recurso usado & Estado documentado \\
\midrule
\codepath{abm} &
\codepath{hidden_attractors/integrations/abm.py} y ruta Caputo de historia completa. &
Solución exacta \(y(t)=y_0E_q\!\left(\lambda t^q\right)\), solución manufacturada \(y(t)=t^m\) y sistema lineal diagonal \(D^q X = A X\). &
\codepath{method_validated_against_exact_solution}; error decreciente para Mittag-Leffler y solución manufacturada, y validación vectorial lineal. \\
\codepath{rk4} &
\codepath{hidden_attractors/integrations/rk4.py} y ruta general \(q=1\). &
Soluciones exactas de \(y'=-y\), oscilador armónico, sistema lineal diagonal y comparación con \codepath{scipy.solve_ivp} usando DOP853~\cite{SciPy2020}. &
\codepath{method_validated_against_exact_and_solve_ivp}; orden observado \(4\), deriva de energía decreciente y consistencia contra \codepath{solve_ivp}. \\
\codepath{efork3} &
\codepath{hidden_attractors/integrations/efork.py}, \codepath{hidden_attractors/solvers/efork_published.py} y validación de orden de etapas. &
Ghoreishi, Ghaffari y Saad (2023), \emph{Fractional Order Runge-Kutta Methods}, tablas 3, 4, 9 y 10; scripts \codepath{constantes_efork.py} y \codepath{ejemplo1.py} archivados como procedencia. &
\codepath{passed_published_efork3_tables}; diferencia absoluta máxima \(5.5489224\times10^{-9}\) contra valores impresos, tolerancia \(6\times10^{-9}\). \\
\bottomrule
\end{longtable}
\end{center}

La evidencia de convergencia de integradores se conserva en los tests y contratos de validación. En este reporte unificado no se incluye una figura independiente de errores de convergencia si no forma parte del paquete editorial generado para los tres casos Chua.


Los artefactos principales son:
\begin{itemize}
  \item \codepath{version_2/validation/outputs/integrator_method_validation/integrator_method_validation_summary.json}, que registra ABM como validado contra soluciones exactas y RK4 como validado contra soluciones exactas y \codepath{solve_ivp}.
  \item \codepath{version_2/validation/integrator_method_validation/abm_caputo.yaml}, que declara las pruebas ABM contra Mittag-Leffler, solución manufacturada y sistema lineal diagonal.
  \item \codepath{version_2/validation/integrator_method_validation/rk4_integer.yaml}, que declara los problemas de prueba, mallas y criterios de aceptación para RK4.
  \item \codepath{version_2/tests/test_integrator_selector.py} y \codepath{version_2/tests/test_general_solver.py}, que cubren compatibilidad \(q=1\), rechazo de integradores enteros para \(q<1\) y el decaimiento exacto usado como prueba operacional.
  \item {\footnotesize \begin{tabular}[t]{@{}l@{}}\codepath{version_2/validation/reference_cases/}\\\codepath{efork3_ghoreishi_ghaffari/}\\\codepath{03_integrators/efork3_validation_summary.json}\end{tabular}}, que registra la reproducción de las tablas publicadas de EFORK-3.
  \item \codepath{version_2/docs/efork3_validation.md}, que explica la procedencia de los scripts independientes y la corrección del orden \(K_3=a_{31}K_1+a_{32}K_2\).
\end{itemize}

\section{Diagnósticos matemáticos independientes de la búsqueda de ocultedad}
\label{sec:diagnosticos-independientes}

Las funciones de esta sección aceptan una trayectoria o una serie temporal
suministrada por la persona usuaria. No ejecutan por sí mismas una búsqueda de
atractores ocultos y sus salidas no contienen una decisión de cuenca. Para una
trayectoria uniformemente muestreada se usa
\[
 t_n=t_0+n\Delta t,\qquad
 X_n=(x_{n,1},\ldots,x_{n,d})^\mathsf{T},\qquad n=0,\ldots,N-1.
\]
Cuando se declara un tiempo de descarte \(t_{\rm burn}\), todas las métricas se
calculan sobre el conjunto de índices
\(\mathcal I=\{n:t_n\geq t_{\rm burn}\}\). Las etiquetas devueltas son
diagnósticos finito-temporales; no son pruebas asintóticas de acotamiento,
caos, periodicidad u ocultedad.

\subsection{Métricas de trayectoria y acotamiento}

Para cada componente, \codepath{compute_trajectory_metrics} calcula el rango y
la varianza de la cola:
\[
 R_j=\max_{n\in\mathcal I}x_{n,j}
     -\min_{n\in\mathcal I}x_{n,j},
 \qquad
 s_j^2=\frac{1}{|\mathcal I|}
 \sum_{n\in\mathcal I}(x_{n,j}-\bar x_j)^2.
\]
Si se proporcionan equilibrios \(E_\ell\), también calcula
\[
 d_{\min}(X_{N-1})=\min_\ell\|X_{N-1}-E_\ell\|_2
\]
y sólo registra contacto final cuando esa distancia no excede
\codepath{equilibrium_tol}. La etiqueta \codepath{bounded_nontrivial} significa
únicamente que la trayectoria es finita, no rebasó el radio numérico de
divergencia y no terminó dentro de la tolerancia de un equilibrio.

\codepath{compute_boundedness_metrics} forma la serie
\[
 r_n=
 \begin{cases}
 \|X_n\|_2,&\texttt{norm="euclidean"},\\
 \|X_n\|_\infty,&\texttt{norm="max"},
 \end{cases}
\]
y, con una fracción de ventana \(\gamma\in(0,0.5]\), usa
\[
 G=\frac{\operatorname{mean}(r_{N-m},\ldots,r_{N-1})}
 {\max\{\operatorname{mean}(r_0,\ldots,r_{m-1}),\epsilon_{\rm mach}\}},
 \qquad m=\max\{2,\operatorname{round}(\gamma N)\}.
\]
La implementación clasifica \codepath{unbounded_candidate} si se rebasa
\codepath{divergence_radius} o si \(G\geq100\); clasifica
\codepath{bounded_candidate} si el máximo es finito y \(G<10\); los valores
intermedios son inconclusos. Estos umbrales son reglas operativas de la
librería, no teoremas de estabilidad.

\begin{verbatim}
from hidden_attractors import (
    compute_boundedness_metrics,
    compute_trajectory_metrics,
    trajectory_metrics_for_system,
)

trajectory_metrics = compute_trajectory_metrics(
    times, states, equilibria=equilibria, t_start=40.0
)
system_metrics = trajectory_metrics_for_system(
    trajectory, system=system, h=0.01, t_start=40.0
)
boundedness = compute_boundedness_metrics(
    times, states, burn_time=40.0,
    divergence_radius=120.0,
)
\end{verbatim}

\codepath{trajectory_metrics_for_system} es una fachada para las mismas
ecuaciones: acepta una matriz \((t,X)\) y obtiene los equilibrios del objeto
\codepath{system}, o acepta el diccionario \codepath{equilibria} de forma
explícita. Si \codepath{has_time=False}, construye
\(t_n=nh\); si la matriz ya contiene tiempo, \(h\) no modifica sus muestras.
La función exige una de las dos fuentes de equilibrios y no ejecuta búsqueda
de semillas ni una prueba de ocultedad.

\begin{verbatim}
trajectory_metrics_for_system(
    traj, *, system=None, equilibria=None, h, t_start,
    divergence_norm=120.0, equilibrium_tol=1.0e-3,
    has_time=True,
)
\end{verbatim}

\subsection{DFT, FFT, potencia, entropía espectral y Welch}

Después de retirar tendencia o media y aplicar una ventana \(w_n\), la
transformada discreta implementada por NumPy es
\[
 Y_k=\sum_{n=0}^{N-1}w_n y_n
       \exp\!\left(-2\pi\ii\frac{kn}{N}\right),
 \qquad
 f_k=\frac{k}{N\Delta t},
 \qquad 0\leq k\leq\left\lfloor\frac N2\right\rfloor .
\]
\codepath{fft_spectrum} devuelve la amplitud unilateral
\[
 A_k=\frac{|Y_k|}{\sum_n w_n},
\]
mientras que \codepath{compute_fft_psd} usa potencia no normalizada
\(P_k=|Y_k|^2\) y, por defecto, la normaliza como
\(p_k=P_k/\sum_jP_j\). De esa distribución obtiene
\[
 H_{\rm spec}=
 -\frac{\sum_{k:p_k>0}p_k\log p_k}{\log K},
 \qquad
 D_{\rm peak}=\frac{\max_k p_k}{\sum_kp_k}.
\]
La fracción de ancho de banda es el número mínimo de bins de mayor potencia
que acumulan el \(90\) por ciento, dividido por el número total de bins. La
clasificación de pico dominante, candidato cuasiperiódico o espectro de banda
ancha usa umbrales explícitos del código y no certifica caos. La FFT rápida es
el algoritmo computacional para evaluar la DFT
\cite{CooleyTukey1965}; la entropía normalizada deriva de la entropía de
Shannon \cite{Shannon1948}.

Para \codepath{psd_welch}, cada segmento solapado \(r\), de longitud \(L\), se
centra y multiplica por una ventana de Hann. La estimación implementada es
\[
 \widehat S_{yy}(f_k)=
 \frac1{M}\sum_{r=1}^{M}
 c_k\frac{\Delta t\left|\sum_{n=0}^{L-1}
 w_n y_{r,n}\exp(-2\pi\ii kn/L)\right|^2}
 {\sum_{n=0}^{L-1}w_n^2},
\]
donde \(c_k=2\) en los bins interiores positivos y \(c_k=1\) en continua y,
para longitud par, Nyquist. Esta es la densidad unilateral de Welch
\cite{Welch1967}; si \(y\) tiene unidades \(U\) y el tiempo unidades \(T\), la
salida tiene unidades \(U^2/\mathrm{Hz}=U^2T\). La implementación se contrasta
con \codepath{scipy.signal.welch} usando la misma ventana explícita, solapamiento
y detrend. La entrada debe estar uniformemente muestreada; aumentar el
solapamiento no crea información y la ventana intercambia fuga espectral por
resolución.

\begin{verbatim}
from hidden_attractors import compute_fft_psd
from hidden_attractors.analysis import fft_spectrum, psd_welch

summary = compute_fft_psd(
    times, states[:, 0], burn_time=40.0,
    window="hann", normalize_power=True,
)
fft = fft_spectrum(states[:, 0], h=times[1] - times[0])
welch = psd_welch(
    states[:, 0], h=times[1] - times[0],
    nperseg=512, overlap=0.5,
)
\end{verbatim}

\subsection{Cruces de una sección de Poincaré}

\par\smallskip\noindent
Para la sección \(x_s=c\), la función
\codepath{detect_poincare_crossings} acepta un cruce positivo cuando
\[
 x_{n,s}<c\leq x_{n+1,s},
 \qquad x_{n+1,s}-x_{n,s}>0,
\]
y un cruce negativo cuando
\[
 x_{n,s}>c\geq x_{n+1,s},
 \qquad x_{n+1,s}-x_{n,s}<0.
\]
Estas desigualdades evitan contar dos veces una muestra situada exactamente
en la sección. Con interpolación lineal calcula
\[
 \theta_n=\frac{c-x_{n,s}}{x_{n+1,s}-x_{n,s}},\qquad
 t_\star=t_n+\theta_n(t_{n+1}-t_n),\qquad
 X_\star=X_n+\theta_n(X_{n+1}-X_n).
\]
En \codepath{derivative_mode="integer_rhs"}, además se exige que
\(F_s(t_\star,X_\star)\) tenga el signo de cruce. En
\codepath{derivative_mode="geometric_fractional"} sólo se usa la orientación
de las muestras. Por tanto, para Caputo la salida es una sección geométrica
numérica y no un mapa de retorno clásico exacto. Tampoco el modo entero es un
mapa exacto: sigue siendo una interpolación lineal entre muestras; la
evaluación de \(F\) únicamente aporta la prueba clásica de orientación. En
consecuencia, la API registra \codepath{exact_poincare_map=False} en todos los
modos y distingue el caso entero mediante
\codepath{classical_integer_section_interpretation=True}. La separación entre
localización del evento y orientación sigue la práctica de cálculo numérico de
mapas de Poincaré \cite{Henon1982}.

\begin{verbatim}
from hidden_attractors import detect_poincare_crossings

section = detect_poincare_crossings(
    times,
    states,
    section_variable="x",
    section_value=0.0,
    direction="positive",
    derivative_mode="geometric_fractional",
    burn_time=40.0,
)
\end{verbatim}

\subsection{Prueba 0--1 de Gottwald--Melbourne}

Para una señal escalar \(\phi_j\) y un valor
\(c\in(\pi/5,4\pi/5)\), la implementación construye
\[
 p_c(n)=\sum_{j=1}^{n}\phi_j\cos(jc),\qquad
 q_c(n)=\sum_{j=1}^{n}\phi_j\sin(jc).
\]
Para un retardo \(\ell\), calcula el desplazamiento cuadrático medio
\[
 M_c(\ell)=\frac1{N-\ell}
 \sum_{j=1}^{N-\ell}
 \left([p_c(j+\ell)-p_c(j)]^2+
       [q_c(j+\ell)-q_c(j)]^2\right),
\]
retira el término oscilatorio
\[
 V_{\rm osc}(\ell)=
 \bar\phi^2\frac{1-\cos(\ell c)}{1-\cos c},
 \qquad
 D_c(\ell)=M_c(\ell)-V_{\rm osc}(\ell),
\]
y usa la correlación
\[
 K_c=\operatorname{corr}\bigl((1,\ldots,L),
 (D_c(1),\ldots,D_c(L))\bigr).
\]
El valor reportado es la mediana de \(K_c\) sobre valores reproducibles de
\(c\), siguiendo la formulación robusta de
Gottwald y Melbourne \cite{GottwaldMelbourne2004,GottwaldMelbourne2009}.
La librería usa \(K>0.8\) como candidato caótico, \(K<0.2\) como candidato
regular y la zona intermedia como inconclusa. Ruido, tendencia, sobremuestreo
y transitorios pueden sesgar el resultado.

\begin{verbatim}
from hidden_attractors import zero_one_test

result_01 = zero_one_test(
    states[:, 0],
    n_c=100,
    random_seed=12345,
    detrend=True,
    normalize=True,
)
\end{verbatim}

\subsection{Postprocesamiento para diagramas de bifurcación}

La función \codepath{bifurcation_points_from_trajectories} no realiza
continuación de ramas. Para cada parámetro \(\mu_r\), descarta el transitorio y
selecciona máximos locales de un observable \(y_n\):
\[
 y_n-y_{n-1}\geq0,\qquad y_n-y_{n+1}\geq0.
\]
Los modos \codepath{minima}, \codepath{both} y \codepath{sample} cambian esa
regla. El diagrama es el conjunto
\(\{(\mu_r,y_{r,n}):n\in\mathcal M_r\}\); no identifica automáticamente
puntos de bifurcación, estabilidad de ramas ni multiplicadores de Floquet.
Para un conjunto no vacío \(\mathcal B\) de puntos extraídos,
\codepath{bifurcation_summary} sólo resume
\[
 n_{\mathcal B}=|\mathcal B|,\qquad
 \mu_{\min}=\min_{b\in\mathcal B}\mu_b,\quad
 \mu_{\max}=\max_{b\in\mathcal B}\mu_b,
\]
\[
 y_{\min}=\min_{b\in\mathcal B}y_b,\qquad
 y_{\max}=\max_{b\in\mathcal B}y_b.
\]
Si \(\mathcal B\) está vacío devuelve únicamente
\codepath{n_points=0}; este resumen tampoco localiza bifurcaciones.
Su firma pública es \codepath{bifurcation_summary(points)}.

\begin{verbatim}
from hidden_attractors import (
    bifurcation_points_from_trajectories,
    bifurcation_summary,
)
from hidden_attractors.plotting import plot_bifurcation_diagram

points = bifurcation_points_from_trajectories(
    scans, observable="x", t_start=40.0, mode="maxima"
)
point_summary = bifurcation_summary(points)
plot_bifurcation_diagram(
    points, "bifurcation.pdf",
    parameter_label="mu",
    observable_label="maximos locales de x",
)
\end{verbatim}

\subsection{Exponentes de Lyapunov desde las ecuaciones}

\subsubsection{QR--Benettin para \(q=1\)}

Para \(\dot X=F(X)\), la matriz variacional satisface
\[
 \dot\Phi(t)=J(X(t))\Phi(t),\qquad
 J(X)=\frac{\partial F}{\partial X},\qquad \Phi(0)=I.
\]
Después de cada intervalo de reortogonalización,
\(\Phi_k=Q_kR_k\), y el estimador finito-temporal es
\[
 \widehat\lambda_i(T)=
 \frac1T\sum_{k=1}^{m}\log |(R_k)_{ii}|.
\]
Ésta es la ruta \codepath{integer_qr_benettin}, basada en el método de
Benettin y la práctica de reortogonalización para espectros numéricos
\cite{Benettin1980a,Benettin1980b,Wolf1985}. Sólo es válida para \(q=1\).
Si no se proporciona Jacobiano analítico, la librería usa diferencias finitas
y registra la advertencia correspondiente. La discretización implementada
debe quedar explícita:
\[
 X_{n+1}=\operatorname{EFORK}_{q=1}(F,X_n,h),\qquad
 \Phi_{n+1}=\Phi_n+hJ(X_n)\Phi_n.
\]
Es decir, el estado usa el paso de tres etapas
\codepath{efork_q1_step}, mientras que la base tangente usa Euler explícito
antes de cada QR. La exactitud de los exponentes debe comprobarse refinando
\(h\); la etiqueta QR--Benettin identifica la construcción variacional y la
reortogonalización, no un mismo integrador de alto orden para ambos bloques.
\codepath{integer_qr_benettin_lyapunov_exponents} es la fachada canónica
directa para \(F\), \(J\) y \(X_0\); reutiliza el mismo núcleo y no constituye
un segundo método. Impone el contrato entero
\(\lvert q-1\rvert\leq10^{-9}\).
\codepath{integer_system_lyapunov_exponents} adapta un
\codepath{ChaoticSystem}: usa su \codepath{evaluate}, emplea
\codepath{jacobian_matrix} cuando el sistema declara un Jacobiano analítico y,
en caso contrario, aproxima
\[
 J_{ij}(X)\approx
 \frac{F_i(X+\varepsilon e_j)-F_i(X-\varepsilon e_j)}
 {2\varepsilon}.
\]
La fachada intenta inferir el orden declarado del sistema y lo rechaza cuando
\(\lvert q_{\mathrm{system}}-1\rvert>10^{-9}\). Si el objeto no declara orden, lo acepta por
compatibilidad; esa aceptación no demuestra que el modelo sea entero.

\begin{verbatim}
integer_qr_benettin_lyapunov_exponents(
    rhs, jacobian, x0, *, h, t_final, t_burn=0.0,
    reorthonormalize_every=10, jacobian_eps=1.0e-6,
    div_threshold=None, q=1.0,
)
integer_system_lyapunov_exponents(
    system, x0, *, h, t_final, t_burn=0.0,
    reorthonormalize_every=10, jacobian_eps=1.0e-6,
    div_threshold=None,
)
\end{verbatim}

\begin{verbatim}
from hidden_attractors import (
    integer_qr_benettin_lyapunov_exponents,
    integer_system_lyapunov_exponents,
)

direct = integer_qr_benettin_lyapunov_exponents(
    rhs, jacobian, x0, h=0.01, t_final=20.0,
    t_burn=2.0, q=1.0,
)
registered = integer_system_lyapunov_exponents(
    system, x0, h=0.01, t_final=20.0, t_burn=2.0,
)
\end{verbatim}

\subsubsection{Sistema variacional Caputo ABM--QR}

Para \(0<q<1\), la ruta
\codepath{fractional_variational_abm_qr} integra el sistema extendido
\[
 \CaputoD X=F(X),\qquad
 \CaputoD\Phi=J(X)\Phi,\qquad \Phi(0)=I
\]
con ABM. Como el operador de Caputo depende de la historia, después de
\(\Phi_k=Q_kR_k\) la implementación coherente transforma cada bloque
variacional almacenado:
\[
 \Phi_j\leftarrow \Phi_jR_k^{-1},
 \qquad
 G_j\leftarrow
 \begin{bmatrix}F(X_j)\\J(X_j)\Phi_j\end{bmatrix},
 \quad j\leq k.
\]
Con \codepath{memory_mode="full"} se conserva toda la historia; una ventana
finita cambia el problema numérico. El sistema original--variacional se apoya
en la formulación discutida por Danca y Kuznetsov
\cite{DancaKuznetsov2018}; la transformación de \emph{toda} la historia
almacenada después de QR es una extensión propia de esta librería y no una
reproducción literal del algoritmo publicado. La implementación tiene pruebas
sintéticas, pero no se presenta como validación cuantitativa publicada para
todos los modelos, órdenes o superficies no suaves.

\subsubsection{Dinámica clonada}

Las dos rutas clonadas no requieren Jacobiano. En cada bloque se integran una
trayectoria fiducial y \(d\) copias
\[
 X^{(i)}_0=X_0+\delta e_i,\qquad i=1,\ldots,d.
\]
Al final del bloque se forma
\[
 V_k=\left[
 X^{(1)}_k-X^{(0)}_k,\ldots,
 X^{(d)}_k-X^{(0)}_k
 \right],
\]
se ortogonaliza mediante Gram--Schmidt o QR, y se acumulan las normas
\(\rho_{k,i}\):
\[
 \widehat\lambda_i=
 \frac1{K\,t_{\rm clone}}
 \sum_{k=1}^{K}\log\frac{\rho_{k,i}}{\delta}.
\]
\codepath{fractional_cloned_dynamics_abm_gs_published} sigue la ruta de
Fischer, Zourmba y Mohamadou \cite{Fischer2020}; la variante
\codepath{fractional_cloned_dynamics_abm_qr} sustituye GS por QR. Ambas usan
la misma semántica de reinicio ABM por bloques para la integración
fraccionaria. La ejecución registra
\codepath{effective_memory_protocol="published_block_restart"}; los nombres
anteriores, como \codepath{experimental_qr_block_restart}, sólo se conservan
en \codepath{requested_memory_protocol} como alias de compatibilidad. La
diferencia efectiva de la variante es la ortogonalización QR, no otra historia
de Caputo. Ninguna de las dos es un estimador Caputo de memoria completa.

\begin{center}
\footnotesize
\begin{longtable}{@{}L{0.27\textwidth}L{0.12\textwidth}L{0.13\textwidth}L{0.20\textwidth}L{0.20\textwidth}@{}}
\toprule
Método solicitado & Orden & Jacobiano & Memoria & Estado de evidencia \\
\midrule
\codepath{integer_qr_benettin} & \(q=1\) & Opcional; DF centradas & No aplica & Controles lineales exactos e intercomparación interna; sin reproducción cuantitativa de un espectro publicado.\\
\codepath{fractional_variational_abm_qr} & \(0<q<1\) & Opcional; DF centradas & Completa o ventana declarada & Implementado; validación sintética, sin claim cuantitativo publicado general.\\
\codepath{fractional_cloned_dynamics_abm_gs_published} & \(0<q\leq1\) & No & Reinicio por bloques & Reproducción con discrepancia de benchmark registrada.\\
\codepath{fractional_cloned_dynamics_abm_qr} & \(0<q\leq1\) & No & Reinicio por bloques & Variante numérica para comparación; sin claim publicado.\\
\bottomrule
\end{longtable}
\end{center}

\subsubsection{Contrato explícito de validación del método}

\codepath{validate_lyapunov_method_request} comprueba una
\codepath{LyapunovComputationRequest} sin ejecutar el integrador y devuelve
exactamente la terna
\[
 (\mathrm{ok},\mathrm{status},\mathrm{warnings}).
\]
Para una solicitud compatible devuelve
\codepath{(True, "compatible", warnings)}; una incompatibilidad ordinaria se
representa con \codepath{ok=False} y no mediante una excepción. Primero exige
una fuente y sólo una para el campo vectorial: exactamente una de
\codepath{system} o \codepath{rhs}:
\[
 \mathbf{1}_{\{\mathrm{system}\neq\mathrm{None}\}}+
 \mathbf{1}_{\{\mathrm{rhs}\neq\mathrm{None}\}}=1.
\]
También exige
\[
 X_0\in\mathbb{R}^{d},\ d\geq1,\quad
 h>0,\quad T_{\mathrm{final}}>0,\quad T_{\mathrm{burn}}\geq0,\quad
 \varepsilon_J>0,
\]
con \(X_0\), \(q\), los tiempos, \(h\) y \(\varepsilon_J\) finitos.
Si se proporciona, \codepath{reorthonormalize_every} debe ser un entero
positivo. \codepath{memory_window} sólo se inspecciona en la ruta variacional
con \codepath{memory_mode="window"} y entonces también debe ser un entero
positivo. Asimismo,
\codepath{reorthonormalization_time} y \codepath{div_threshold} deben ser
finitos y positivos.

La validación específica reproduce el registro de métodos: QR--Benettin
entero requiere \(q=1\) y memoria \codepath{not_applicable}; ABM--QR
variacional requiere \(0<q<1\), memoria \codepath{full} o
\codepath{window}, y comprueba la igualdad entre el \(q\) solicitado y el
declarado por el sistema cuando ambos existen. Las dos rutas de dinámica
clonada requieren \(0<q\leq1\) y aceptan
\codepath{not_applicable}, \codepath{published_block_restart} o el alias
\codepath{experimental_qr_block_restart}. La ausencia de Jacobiano en los
métodos variacionales produce
\codepath{analytic_jacobian_missing_finite_difference_used}; las rutas
clonadas producen \codepath{cloned_dynamics_no_jacobian_required}. Son
advertencias, no certificaciones de caos, ocultedad o validez fraccionaria.

Los estados de incompatibilidad que puede devolver el código actual son:
\begin{itemize}
  \item \codepath{unknown_method};
  \item \codepath{invalid_parameter};
  \item \codepath{method_not_valid_for_fractional_caputo};
  \item \codepath{memory_mode_not_applicable_for_integer_method};
  \item \codepath{method_not_valid_for_integer_or_out_of_range_q};
  \item \codepath{memory_mode_must_be_full_or_window_for_fractional_method};
  \item \codepath{request_q_does_not_match_system_q};
  \item \codepath{method_not_valid_for_out_of_range_q};
  \item \codepath{memory_mode_must_be_block_restart_for_cloned_dynamics}.
\end{itemize}

\begin{verbatim}
from hidden_attractors import (
    LyapunovComputationRequest,
    validate_lyapunov_method_request,
)

request = LyapunovComputationRequest(
    system=system, rhs=None, jacobian=None, x0=x0,
    q=1.0, method="integer_qr_benettin",
    h=0.01, t_final=20.0, t_burn=2.0,
    reorthonormalize_every=10,
    memory_mode="not_applicable",
)
ok, status, warnings = validate_lyapunov_method_request(request)
\end{verbatim}

\begin{verbatim}
import numpy as np
from hidden_attractors import compute_lyapunov_spectrum

rhs = lambda x: np.array([-x[0], -2.0*x[1]])
jac = lambda x: np.array([[-1.0, 0.0], [0.0, -2.0]])

summary = compute_lyapunov_spectrum(
    rhs=rhs,
    jacobian=jac,
    x0=np.array([1.0, 1.0]),
    q=1.0,
    method="integer_qr_benettin",
    h=0.01,
    t_final=20.0,
    t_burn=2.0,
)
print(summary.result.exponents)
\end{verbatim}

\subsection{Lyapunov desde una serie temporal escalar}

Esta ruta es distinta de los métodos variacionales. A partir de una señal
\(s_n\), una reconstrucción por retardos forma \cite{Takens1981}
\[
 Y_i=(s_i,s_{i+\tau},\ldots,s_{i+(m-1)\tau})^\mathsf{T}.
\]
Rosenstein selecciona para cada \(Y_i\) un vecino temporalmente separado
\(Y_{j(i)}\) y calcula \cite{Rosenstein1993}
\[
 d(k)=\left\langle
 \log\|Y_{i+k}-Y_{j(i)+k}\|_2
 \right\rangle_i,\qquad
 \widehat\lambda_{\max}=
 \frac1{\Delta t}\frac{\mathrm d d(k)}{\mathrm d k}.
\]
Eckmann ajusta mapas lineales locales de la forma
\[
 Y_{i+1}-Y_{j+1}\approx
 A_i(Y_i-Y_j)
\]
y obtiene un espectro finito mediante productos y
reortogonalización \cite{Eckmann1986}. Ordenados
\(\lambda_1\geq\cdots\geq\lambda_m\), la dimensión de Kaplan--Yorke es
\cite{Frederickson1983}
\[
 D_{\rm KY}=j+
 \frac{\sum_{i=1}^{j}\lambda_i}{|\lambda_{j+1}|},
 \qquad
 j=\max\left\{r:\sum_{i=1}^{r}\lambda_i\geq0\right\}.
\]
La función conserva intervalo de muestreo, unidades, parámetros de
reconstrucción, versión del backend, calidad del ajuste, advertencias y un
límite explícito para la matriz de distancias. La reconstrucción depende del
observable, del muestreo, de \(m\), \(\tau\), la exclusión temporal y la
ventana de ajuste.

\begin{verbatim}
from hidden_attractors import estimate_time_series_lyapunov

result = estimate_time_series_lyapunov(
    signal,
    sample_interval=0.01,
    time_unit="s",
    observable="x",
    rosenstein_emb_dim=8,
    eckmann_emb_dim=9,
    eckmann_matrix_dim=3,
)
print(result.largest_exponent, result.exponent_unit)
print(result.spectrum, result.kaplan_yorke_dimension)
\end{verbatim}

\subsection{Medidas de complejidad mediante backends opcionales}

\codepath{compute_complexity_measures} es un adaptador: no reimplementa
\codepath{nolds} ni \codepath{antropy}. Según el backend instalado puede
solicitar:
\begin{itemize}
  \item entropía muestral, \(\operatorname{SampEn}(m,r)
  =-\log(A_{m+1}/B_m)\) \cite{RichmanMoorman2000};
  \item dimensión de correlación,
  \(D_2=\lim_{r\to0}\mathrm d\log C(r)/\mathrm d\log r\)
  \cite{GrassbergerProcaccia1983};
  \item exponente de Rosenstein y exponente de Hurst
  \cite{Hurst1951};
  \item DFA, \(F(s)\propto s^\alpha\) \cite{Peng1994};
  \item entropía de permutación normalizada \cite{BandtPompe2002} y
  dimensión fractal de Higuchi \cite{Higuchi1988}.
\end{itemize}
Las convenciones exactas, valores predeterminados y correcciones pertenecen a
la versión registrada del backend. Por ello el resultado debe citar tanto la
librería adaptadora como el backend que realizó el cálculo.

\begin{verbatim}
from hidden_attractors.integrations import (
    available_complexity_backends,
    compute_complexity_measures,
)

print(available_complexity_backends())
complexity = compute_complexity_measures(
    signal,
    backend="auto",
    sample_rate=100.0,
    measures=[
        "sample_entropy",
        "permutation_entropy",
        "spectral_entropy",
        "dfa",
    ],
)
\end{verbatim}

\subsection{ADM Wu2023: ecuación implementada y frontera científica}

La ruta \codepath{adm_wu2023_integrate} es una reproducción específica del
Chua arctan de Wu et al. \cite{Wu2023ArctanChua}. En cada paso aproxima
\[
 X_{n+1}=X_n+
 \sum_{r=1}^{4}C^{(r)}
 \frac{h^{rq}}{\Gamma(rq+1)}.
\]
Sea \(g(x)=\arctan(\rho x)\), \(s=1+\rho^2x_n^2\), y
\[
 g'_n=\frac{\rho}{s},\quad
 g''_n=-\frac{2\rho^3x_n}{s^2},\quad
 g'''_n=\frac{8\rho^5x_n^2}{s^3}-\frac{2\rho^3}{s^2}.
\]
El primer coeficiente es \(C^{(1)}=F(X_n)\). Para la primera componente, los
polinomios de Adomian usados por el código son
\[
\begin{aligned}
 A_1={}&g'_n C^{(1)}_1,\\
 A_2={}&g'_n C^{(2)}_1+
 \frac{g''_n}{2}(C^{(1)}_1)^2
 \frac{\Gamma(2q+1)}{\Gamma(q+1)^2},\\
 A_3={}&g'_n C^{(3)}_1+
 g''_nC^{(1)}_1C^{(2)}_1
 \frac{\Gamma(3q+1)}{\Gamma(q+1)\Gamma(2q+1)}
 +\frac{g'''_n}{6}(C^{(1)}_1)^3
 \frac{\Gamma(3q+1)}{\Gamma(q+1)^3}.
\end{aligned}
\]
Para \(r=1,2,3\), el siguiente nivel satisface
\[
\begin{aligned}
 C^{(r+1)}_1&=-\alpha(1+a_1)C^{(r)}_1
 +\alpha C^{(r)}_2-\alpha a_2 A_r,\\
 C^{(r+1)}_2&=C^{(r)}_1-C^{(r)}_2+C^{(r)}_3,\\
 C^{(r+1)}_3&=-\beta C^{(r)}_2-\gamma C^{(r)}_3.
\end{aligned}
\]
Cada actualización usa únicamente \(X_n\). No contiene la convolución de
historia de Caputo y no es equivalente a ABM o EFORK-3 con memoria completa.
Su uso defendible es la reproducción metodológica, no sustituir una
integración Caputo.

\begin{verbatim}
import numpy as np
from hidden_attractors.integrations import adm_wu2023_integrate

times, states, status, info = adm_wu2023_integrate(
    params={
        "alpha": 10.0, "beta": 14.87, "gamma": 0.0,
        "a1": -1.27, "a2": 1.08, "rho": 1.0,
    },
    x0=np.array([0.1, 0.0, 0.0]),
    q=0.99,
    h=0.005,
    N=2000,
)
\end{verbatim}

\subsection{Integradores: selector y llamadas}

\begin{center}
\footnotesize
\begin{longtable}{@{}L{0.25\textwidth}L{0.16\textwidth}L{0.25\textwidth}L{0.26\textwidth}@{}}
\toprule
Entrada & Orden & Llamada mínima & Contrato \\
\midrule
\codepath{integrate(..., integrator="rk4")} & \(q=1\) &
\codepath{integrate(rhs,x0,1,h,tf,"rk4")} & RK4 clásico.\\
\codepath{integrate(..., integrator="heun")} & \(q=1\) &
\codepath{integrate(rhs,x0,1,h,tf,"heun")} & Trapecio explícito de orden dos.\\
\codepath{integrate(..., integrator="efork_q1")} & \(q=1\) &
\codepath{integrate(rhs,x0,1,h,tf,"efork_q1")} & Límite entero de coeficientes EFORK.\\
\codepath{integrate(..., integrator="abm")} & \(0<q<1\) &
\codepath{integrate(rhs,x0,q,h,tf,"abm")} & ABM Caputo; declarar memoria.\\
\codepath{integrate(..., integrator="efork3")} & \(0<q\leq1\) &
\codepath{integrate(rhs,x0,q,h,tf,"efork3")} & EFORK-3; en \(q=1\) redirige al límite entero.\\
\codepath{adm_wu2023_integrate} & \(0<q\leq1\) &
\codepath{adm_wu2023_integrate(params,x0,q,h,N)} & Reproducción arctan local, sin memoria Caputo.\\
\codepath{integrate_fractional_abm} & órdenes por componente &
\codepath{integrate_fractional_abm(rhs,x0,orders,h,N)} & Ruta experimental de reinicio por bloques; no es contrato público general no conmensurado.\\
\bottomrule
\end{longtable}
\end{center}

Las funciones de bajo nivel \codepath{integrate_general},
\codepath{caputo_abm_integrate}, \codepath{efork_integrate},
\codepath{rk4_integrate} y \codepath{fractional_integrate} existen para control
avanzado y pruebas. Para código nuevo se prefiere \codepath{integrate}, salvo
la reproducción ADM, que requiere sus parámetros específicos y se invoca de
forma directa.


\subsection{Entropía de cuencas e incertidumbre de estado final}
\label{sec:basin-entropy-uncertainty}

Estas funciones reciben etiquetas de destino ya calculadas. El mismo código se
puede aplicar a cuencas de un sistema entero o a etiquetas producidas por un
solver fraccionario, pero el segundo experimento debe conservar definición de
derivada, orden, terminal, prehistoria, memoria, integrador, horizonte y
clasificador. Compartir una métrica de postproceso no vuelve Markoviano al
sistema fraccionario ni prueba ocultedad global.

Para una caja no vacía \(i\), con proporciones \(p_{ij}\) de cada destino,
\[
 S_i=-\sum_jp_{ij}\log p_{ij},\qquad
 S_b=\frac1{N_c}\sum_{i=1}^{N_c}S_i.
\]
Si \(\mathcal B\) contiene las cajas con al menos dos destinos, la entropía de
frontera es
\[
 S_{bb}=\frac1{|\mathcal B|}\sum_{i\in\mathcal B}S_i.
\]
\codepath{basin_entropy} divide una rejilla bidimensional entera en cajas no
solapadas. Por defecto exige cajas completas de igual área;
\codepath{partial_boxes="drop"} descarta franjas y
\codepath{"include_equal"} conserva cajas parciales con el mismo peso. Las
etiquetas ignoradas no forman probabilidades. El kernel Numba cuesta
\(\mathcal O(N)\)
en el número de celdas y devuelve conteos de cajas, muestras ignoradas,
fracción de frontera, base y política parcial. El criterio
\(S_{bb}>\log 2\) se reporta con margen y tolerancia como condición suficiente
bajo las hipótesis de Daza et al.~\cite{DazaEtAl2016}; no certifica Wada o
fractalidad sin convergencia en escala.

Para pares de condiciones de referencia y perturbadas,
\[
 f(\varepsilon)=
 \frac{\#\{k:c_k\ne c'_k\}}{N_{\mathrm{valido}}}.
\]
\codepath{uncertainty_fraction} cuesta \(\mathcal O(N)\), usa operaciones
vectorizadas NumPy y devuelve la fracción y un intervalo de Wilson. Ese
intervalo supone pares Bernoulli independientes; en
rejillas espacialmente correlacionadas es sólo descriptivo y no incluye error
de integración o clasificación. En varias escalas,
\codepath{estimate_uncertainty_exponent} ajusta por mínimos cuadrados
\[
 \log f(\varepsilon)=\alpha\log\varepsilon+\log C.
\]
Su costo es \(\mathcal O(S)\) para \(S\) escalas y usa
\codepath{numpy.linalg.lstsq}. Devuelve \(\alpha\), \(C\), error estándar,
\(R^2\), escalas y, si se declara la dimensión de muestreo \(D\), la
estimación \(D-\alpha\). Un buen \(R^2\) no demuestra un régimen asintótico
\cite{McDonaldEtAl1985}.

\begin{verbatim}
from hidden_attractors.analysis.basin_uncertainty import (
    basin_entropy,
    estimate_uncertainty_exponent,
    uncertainty_fraction,
)

entropy = basin_entropy(
    chua_basin_labels,
    box_size=(8, 8),
    ignored_labels=(-1,),
    partial_boxes="require_complete",
)
fractions = [
    uncertainty_fraction(base, shifted, ignored_labels=(-1,)).fraction
    for shifted in perturbed_label_grids
]
scaling = estimate_uncertainty_exponent(
    perturbation_scales,
    fractions,
    sampling_space_dimension=2,
)
\end{verbatim}
Las matrices del ejemplo deben proceder de corridas Chua con su contrato
guardado; los nombres no representan datos incluidos automáticamente por la
librería. \codepath{tests/test_basin_uncertainty.py} verifica cajas exactas,
 políticas parciales, etiquetas ignoradas, criterio \(\log 2\), intervalo de
Wilson y ajustes regulares o degenerados.

\subsection{Reconstrucción generalizada por retardos}
\label{sec:delay-embedding-advanced}

Para pares explícitos \((c_k,\tau_k)\),
\codepath{generalized_delay_embedding} construye
\[
 Y_k(i)=x_{c_k}[i-\tau_k].
\]
Los retardos positivos miran al pasado y los negativos al futuro. La función
interseca exactamente los índices válidos y devuelve vectores, índices ancla y
tiempos alineados. Su costo y memoria son \(\mathcal O(Mp)\) para \(M\) anclas y \(p\)
coordenadas; usa NumPy. Con tiempos irregulares conserva los tiempos reales,
pero no inventa una duración física única para un retardo expresado en índices.
Takens fundamenta la reconstrucción clásica de dimensión finita
\cite{Takens1981}; aplicarla a una proyección de una FDE hereditaria no prueba
que toda su historia admita un embedding difeomorfo finito.

\codepath{estimate_delay_autocorrelation} calcula por FFT con relleno cero
\[
 \rho(k)=
 \frac{\sum_{i=0}^{N-k-1}(x_i-\bar x)(x_{i+k}-\bar x)}
      {\sum_{i=0}^{N-1}(x_i-\bar x)^2}.
\]
Su costo es \(\mathcal O(N\log N)\) con \codepath{numpy.fft}. Selecciona el primer cruce positivo/no positivo o
el primer mínimo local; si el criterio no aparece, registra
\codepath{criterion_not_found}.

\codepath{estimate_delay_mutual_information} usa el estimador histográfico
plug-in
\[
 I(k)=\sum_{a,b}p_{ab}(k)
 \log\frac{p_{ab}(k)}{p_a(k)p_b(k)}
\]
en nats. Reutiliza bordes globales fijos mediante
\codepath{numpy.histogram2d} y cuesta aproximadamente \(\mathcal O(NK)\) para
\(K\) retardos, además del histograma. Adopta el primer mínimo de
Fraser--Swinney~\cite{FraserSwinney1986}, pero no su partición adaptativa
recursiva. \codepath{minimum_pairs} impide que colas con poco solapamiento sean
mínimos artificiales; el mínimo global sólo se usa mediante fallback explícito.

\codepath{false_nearest_neighbors} compara embeddings de dimensión \(m\) y
\(m+1\). En norma euclídea, un vecino es falso si
\[
 \frac{|\Delta x_{m+1}|}{R_m}>R_{\rm tol}
 \quad\text{o}\quad
 \frac{R_{m+1}}{\sigma_x}>A_{\rm tol}.
\]
La búsqueda usa \codepath{scipy.spatial.cKDTree}, ventana de Theiler y desempate
determinista. El costo típico por dimensión es \(\mathcal O(N\log N)\), aunque consultas
degeneradas pueden ser más costosas; la memoria crece con \(Nm\). La atribución
Kennel--Brown--Abarbanel es exacta para Euclídea
\cite{KennelEtAl1992}; Manhattan y Chebyshev se etiquetan como variantes
generalizadas. Cada dimensión conserva número de vecinos válidos y estado, y no
se selecciona con menos de las comparaciones mínimas declaradas.

\begin{verbatim}
from hidden_attractors.analysis.delay_embedding import (
    estimate_delay_mutual_information,
    false_nearest_neighbors,
    generalized_delay_embedding,
)

delay = estimate_delay_mutual_information(
    chua_states,
    column=0,
    max_lag=300,
    minimum_pairs=32,
    dt=0.01,
    time_unit="s",
)
fnn = false_nearest_neighbors(
    chua_states,
    column=0,
    delay=delay.selected_lag,
    max_dimension=8,
    theiler_window=50,
)
embedding = generalized_delay_embedding(
    chua_states,
    ((0, 0), (0, delay.selected_lag), (0, 2 * delay.selected_lag)),
    dt=0.01,
    time_unit="s",
)
\end{verbatim}
\codepath{tests/test_delay_embedding_advanced.py} cubre alineación, retardos
con signo, tiempos irregulares, ACF, información mutua, FNN, Theiler, métricas y
falta de vecinos. Son diagnósticos de muestra finita, no pruebas de caos.

\subsection{Matrices de recurrencia y RQA avanzada}
\label{sec:recurrence-advanced}

\codepath{auto_recurrence_matrix} construye
\[
 R_{ij}=\mathbf1\{\|x_i-x_j\|\le\varepsilon\};
\]
\codepath{cross_recurrence_matrix} usa dos trayectorias y produce una matriz
rectangular; \codepath{joint_recurrence_matrix} hace el AND de auto-recurrencias
sincronizadas. Estas construcciones siguen la geometría de Eckmann--Kamphorst--
Ruelle~\cite{EckmannRecurrence1987} y las convenciones de RQA se apoyan en la
revisión de Marwan et al.~\cite{MarwanEtAl2007}.

Las normas Euclídea, Manhattan y Chebyshev tienen rutas Numba y NumPy. Las
distancias se procesan por bloques, con costo \(\mathcal O(NMd)\), pero la matriz booleana
final sigue ocupando \(\mathcal O(NM)\) y se protege mediante \codepath{max_bytes}. El
umbral puede ser un radio fijo o una tasa objetivo. En el segundo caso se
materializan las distancias elegibles, se elige el menor estadístico de orden
que alcanza el conteo y se incluyen todos los empates; la tasa obtenida puede
superar la solicitada y la memoria se limita separadamente mediante
\codepath{max_threshold_bytes}. No hay aún salida sparse.

En recurrencia conjunta, una tupla de radios fija un umbral por componente. La
tasa objetivo actual selecciona un radio común sobre la distancia máxima antes
del AND; por ello exige escalas comparables y no equivale a todas las políticas
componentwise de RecurrenceAnalysis.jl. Para HAFO,
\codepath{theiler_window=w} excluye \(|i-j|\le w\); así \(w=0\) excluye la línea
de identidad, mientras que en RecurrenceAnalysis.jl \codepath{theiler=0} la
incluye. La matriz devuelta es no escribible para evitar que sus conteos y
metadatos diverjan por mutación posterior.

Sea \(P(l)\) el histograma de longitudes diagonales seleccionadas y \(P(v)\) el
de longitudes verticales. \codepath{recurrence_quantification_advanced}
calcula
\begin{align*}
 RR &= \frac{N_R}{N_{\rm elegible}},
 & DET &= \frac{\sum_{l\ge l_{\min}}lP(l)}{N_R},\\
 L &= \frac{\sum_{l\ge l_{\min}}lP(l)}{\sum_{l\ge l_{\min}}P(l)},
 & DIV &= \frac1{L_{\max}},\\
 ENTR &= -\sum_l p_l\log p_l,
 & LAM &= \frac{\sum_{v\ge v_{\min}}vP(v)}{N_R},\\
 TT &= \frac{\sum_{v\ge v_{\min}}vP(v)}{\sum_{v\ge v_{\min}}P(v)}.
\end{align*}
También devuelve \(L_{\max}\), \(V_{\max}\), entropía vertical y TREND clásico:
la pendiente de densidad diagonal multiplicada por \(1000\), con borde
declarado y el mayor cuadrado inicial en matrices rectangulares. La antigua
pendiente HAFO se conserva como \codepath{normalized_absolute_trend}, separada
y sin el factor \(1000\). Una distribución de líneas vacía produce
\codepath{NaN}; sin diagonal válida, \codepath{DIV=inf}.
El análisis RQA recorre diagonales y columnas en NumPy con costo
\(\mathcal O(NM)\); la matriz densa de entrada domina su memoria.

\begin{verbatim}
from hidden_attractors.analysis.recurrence_advanced import (
    auto_recurrence_matrix,
    cross_recurrence_matrix,
    joint_recurrence_matrix,
    recurrence_quantification_advanced,
)

rp = auto_recurrence_matrix(
    embedding.vectors,
    target_rate=0.03,
    metric="euclidean",
    theiler_window=20,
    backend="auto",
)
rqa = recurrence_quantification_advanced(
    rp,
    min_diagonal=2,
    min_vertical=2,
    trend_border=10,
)
cross = cross_recurrence_matrix(integer_embedding, fractional_embedding,
                                radius=0.2)
joint = joint_recurrence_matrix(x_embedding, y_embedding,
                                radius=(0.15, 0.20))
\end{verbatim}
Para comparar Chua entero y fraccionario deben alinearse tiempos físicos,
normalización, transitorio y observable. Una recurrencia entre estados
proyectados no mide la distancia entre historias completas de una FDE.
\codepath{tests/test_recurrence_advanced.py} verifica matrices exactas,
auto/cross/joint, empates de tasa, Theiler, paridad Numba--NumPy, métricas RQA,
inmutabilidad y guardas de memoria.

\subsection{Entropía de permutación de Bandt--Pompe}
\label{sec:permutation-entropy}

Para una serie escalar finita $x_0,\ldots,x_{N-1}$, una dimensión ordinal
$m\geq2$ y un retardo entero $\tau\geq1$, HAFO construye las ventanas forward
\[
 \mathbf y_s=(x_s,x_{s+\tau},\ldots,x_{s+(m-1)\tau}),
 \qquad s=0,\ldots,W-1,
 \qquad W=N-(m-1)\tau.
\]
Cada ventana se ordena de forma ascendente y estable; el índice temporal dentro
de la ventana rompe empates. La permutación resultante $\pi_s$ se codifica por
su rango de Lehmer lexicográfico
\[
 r(\pi_s)=\sum_{i=0}^{m-2}c_i(m-1-i)!,
 \qquad
 c_i=\#\{j>i:\pi_{s,j}<\pi_{s,i}\},
\]
de modo que $0\leq r<m!$. Esta convención reproduce el alfabeto ordinal de
Bandt y Pompe~\cite{BandtPompe2002} y hace auditable cada celda del histograma.
Si $n_r$ es el conteo del rango $r$ y $V=\sum_r n_r$ el número de ventanas
aceptadas,
\[
 p_r=\frac{n_r}{V},\qquad
 \widehat H_b=-\sum_{p_r>0}p_r\log_b p_r,
 \qquad
 \widehat H_{b,\mathrm{norm}}=
 \frac{\widehat H_b}{\log_b(m!)}.
\]
La normalización usa todo el espacio $m!$, no sólo los patrones observados. El
estimador es la entropía de Shannon plug-in de ventanas solapadas; por ello
conserva advertencias de tamaño finito, dependencia y patrones ausentes
\cite{UnakafovaKeller2013EfficientPermutation,
ReyEtAl2023PermutationAsymptotic}.

Los empates nunca se resuelven implícitamente. \codepath{stable_index} usa el
índice temporal como llave secundaria; \codepath{omit} descarta toda ventana
empatada y cambia el denominador $V$; \codepath{raise} rechaza la serie. Esta
decisión es especialmente importante con sensores cuantizados
\cite{TraversaroEtAl2018TiedValues}. El resultado guarda ventanas candidatas,
válidas, empatadas y omitidas, además de conteos, probabilidades, base,
muestreo, proyección, backend, referencias y huella de trayectoria.

\codepath{ordinal_pattern_distribution} y \codepath{permutation_entropy}
aceptan una serie NumPy o una componente explícita de
\codepath{TrajectoryInput}. El mismo postproceso sirve para una trayectoria
entera o fraccionaria, pero en el segundo caso el escalar observado no es el
estado hereditario completo. Definición de derivada, orden, terminal,
prehistoria y política de memoria permanecen en el contrato; compartir
$(m,\tau)$ no vuelve equivalentes ambos problemas.

Python, Numba y C producen exactamente el mismo histograma \codepath{uint64}.
El kernel C/OpenMP paraleliza ventanas, usa histogramas privados bajo un
presupuesto de 64 MiB y pasa a incrementos atómicos cuando $m!$ haría esa
estrategia impracticable. El coste directo es
$\mathcal O(Wm^2+m!)$ y la memoria del histograma es
$\mathcal O(m!)$; por ello HAFO impone $2\leq m\leq10$. El despacho
\codepath{auto} usa la política medida siguiente:
\begin{itemize}
  \item $m=2$ u $8$: C desde $131\,072$ ventanas;
  \item $3\leq m\leq7$: C desde $32\,768$ ventanas;
  \item $m=9$ o $10$: Numba; C sólo mediante selección explícita.
\end{itemize}
La última guarda evita la penalización atómica observada para histogramas de
$9!$ y $10!=3\,628\,800$ celdas. No se cruza Python--C dentro de cada ventana.

El ejemplo oficial compara un oscilador armónico ordinario integrado con RK4
y su problema componente a componente de Caputo $q=0.85$ resuelto por
ABM--PECE. Ambos usan la coordenada de posición, $m=4$, $\tau=2$ y base dos:
\begin{verbatim}
from hidden_attractors import TrajectoryInput, permutation_entropy

trajectory = TrajectoryInput.from_simulation_result(
    simulation,
    projection=("position", "velocity"),
)
result = permutation_entropy(
    trajectory,
    component="position",
    embedding_dimension=4,
    delay=2,
    tie_policy="stable_index",
    backend="auto",
)
print(result.entropy, result.normalized_entropy)
\end{verbatim}
\codepath{examples/permutation_entropy_integer_fractional.py} genera los dos
problemas completos y emite JSON estricto. La comparación es de API y
proyección finita: no afirma equivalencia física, tasa de entropía, caos,
atracción u ocultedad.

El caso independiente
\codepath{validation/wolfram/cases/permutation_entropy.wl} no lee el código
HAFO ni este reporte. Verifica diez identidades sobre cuatro fixtures: ventanas
$m=3$ con $\tau=1,2$, rangos Lehmer, empates \codepath{stable_index}/
\codepath{omit}, conteos, probabilidades, entropía base dos y normalización.
La ejecución viva y el artefacto persistido pasaron sus dos pruebas Python de
comparación; esto certifica concordancia finita de fórmulas, no comportamiento
asintótico.

\codepath{benchmarks/bench_permutation_entropy.py} separa JIT/compilación del
tiempo medido, rota el orden Python--Numba--C y distingue pipeline público de
kernel ordinal. La corrida local de cierre cubrió 21 cargas deterministas,
$m=2,\ldots,10$ y $4\,096$--$131\,072$ ventanas; todos los histogramas fueron
idénticos y la política anterior resultó respaldada para esas cargas. Los
cruces variables de $16\,384$ ventanas para $m=5,7$ se conservan como
oportunidades, no como umbrales universales. El benchmark es evidencia de
ingeniería del host, no evidencia dinámica.

\subsection{Suma y dimensión de correlación \texorpdfstring{\(q=2\)}{q=2}}
\label{sec:correlation-dimension-q2}

Para $N$ puntos ordenados y una ventana de Theiler $w\geq0$, HAFO usa los
pares no ordenados $i<j$ que satisfacen $j-i>w$. Su número exacto es
\[
 M_w=\frac{(N-w)(N-w-1)}{2},
\]
y la suma de correlación de orden dos se calcula como
\[
 C_2(r)=\frac{1}{M_w}
 \sum_{\substack{i<j\\j-i>w}}
 \mathbf 1\!\left\{d(x_i,x_j)<r\right\}.
\]
La desigualdad es estricta. Los radios deben ser positivos y estrictamente
crecientes; las métricas disponibles son Euclídea, Manhattan y Chebyshev. La
ventana de Theiler evita contar como vecinos geométricos muestras que son
cercanas sólo por continuidad temporal~\cite{GrassbergerProcaccia1983,
GrassbergerProcaccia1983Measuring,Theiler1986}.

\codepath{correlation_sum_curve} devuelve radios, conteos enteros, $C_2$,
denominador, backend efectivo y procedencia de muestreo/proyección. Python es
la referencia transparente; Numba compila el bucle de pares; el kernel
C/OpenMP usa acumuladores por hilo y búsqueda binaria sobre los radios. El
modo \codepath{auto} aplica un umbral conservador de pares y registra cualquier
fallback. No cruza lenguajes en cada paso del integrador porque opera sobre la
trayectoria ya producida.

\codepath{fit_correlation_dimension} ajusta
\[
 \log C_2(r)=b+D_2\log r
\]
únicamente en un \codepath{fit_radius_range} inclusivo declarado por quien
llama y sólo donde $0<C_2(r)<1$. Devuelve pendiente, intercepto, $R^2$,
error estándar de la regresión y pendientes locales. HAFO no selecciona una
región de escala automáticamente: esa decisión requiere análisis explícito y
la literatura muestra que su extracción en datos finitos es un problema por
sí mismo~\cite{DeshmukhEtAl2021}.

\begin{verbatim}
from hidden_attractors import estimate_correlation_dimension

result = estimate_correlation_dimension(
    sampled_states,
    radii,
    theiler_window=25,
    fit_radius_range=(0.02, 0.15),
    backend="auto",
    sampling="uniform dt=0.01 after transient",
    projection="selected state coordinates",
)
\end{verbatim}

La misma llamada acepta trayectorias de solvers enteros o de cualquiera de las
definiciones fraccionarias disponibles. En el segundo caso $D_2$ caracteriza
sólo la proyección suministrada: unas coordenadas finitas no representan en
general todo el estado hereditario. Una validación Wolfram independiente con
seis puntos exactos, $w=1$ y seis radios reproduce el denominador $10$, los
conteos $[0,4,6,6,8,10]$ y el ajuste declarado con diferencia máxima
$3.5527\times10^{-15}$ frente a una tolerancia de $5\times10^{-13}$.

El benchmark reproducible separa calentamiento y medición. En este host, con
cuatro hilos OpenMP, Numba fue ligeramente más rápido para $31\,626$ pares y
C/OpenMP para $124\,750$; Python fue mucho más lento en ambas cargas. Esos
tiempos calibran software local, no prueban superioridad universal ni aportan
evidencia de dimensión fractal, caos, atracción u ocultedad.

\subsection{Índices de alineamiento SALI, GALI y LDI para orden entero}
\label{sec:sali-gali-integer}

Sean $w_i\ne0$, $i=1,\ldots,k$, vectores de desviación y
$\widehat w_i=w_i/\lVert w_i\rVert_2$. El Smaller Alignment Index se define
por~\cite{Skokos2001AlignmentIndices,SkokosEtAl2004SALI}
\[
 \operatorname{SALI}=
 \min\!\left\{
  \lVert\widehat w_1+\widehat w_2\rVert_2,
  \lVert\widehat w_1-\widehat w_2\rVert_2
 \right\}.
\]
Con $W_k=[\widehat w_1\ \cdots\ \widehat w_k]$, el Generalized Alignment
Index es el volumen exterior~\cite{SkokosBountisAntonopoulos2007GALI}
\[
 \operatorname{GALI}_k
 =\lVert\widehat w_1\wedge\cdots\wedge\widehat w_k\rVert
 =\sqrt{\det(W_k^{\mathsf T}W_k)}
 =\prod_{j=1}^{k}\sigma_j(W_k).
\]
Para las mismas columnas normalizadas, el Linear Dependence Index coincide
algebraicamente:
\[
 \operatorname{LDI}_k=\prod_{j=1}^{k}\sigma_j(W_k)
 =\operatorname{GALI}_k.
\]
La igualdad permite usar SVD sin expandir productos exteriores. Para
$k=2$, si $c=\widehat w_1^{\mathsf T}\widehat w_2$, entonces
\[
 \operatorname{SALI}^2=2-2|c|,
 \qquad \operatorname{GALI}_2^2=1-c^2,
\]
y se obtiene la identidad exacta de control
\[
 4\operatorname{GALI}_2^2
 =\operatorname{SALI}^2\bigl(4-\operatorname{SALI}^2\bigr).
\]

La normalización se aplica a cada columna por separado. No se sustituye la
familia evolucionada por el factor $Q$ de una ortogonalización recurrente:
eso mantendría artificialmente independientes las direcciones y eliminaría el
alineamiento que los índices deben medir. QR sólo puede aplicarse sobre una
copia para evaluar el volumen o para construir el conjunto inicial.

Para un flujo ordinario autónomo, HAFO resuelve conjuntamente
\[
 \dot x=f(x),\qquad \dot V=J_f(x)V,
 \qquad V\in\mathbb R^{d\times k},
\]
mediante DOP853. El transitorio se integra sólo para el estado; después de cada
segmento las columnas tangentes se normalizan individualmente. Para un mapa,
\[
 x_{n+1}=F(n,x_n),\qquad V_{n+1}=D F(n,x_n)V_n,
\]
y el Jacobiano se evalúa en $x_n$ antes de actualizar el estado. El método
alternativo multiparticle propaga una referencia y $k$ vecinos, forma sus
separaciones y los reinserta a distancia $d_0$ tras cada intervalo. El estudio
reciente de Manda--Hillebrand--Skokos recomienda para sus controles
Hamiltonianos en doble precisión $d_0\simeq\sqrt{\epsilon}$, intervalo no
mayor que uno y error relativo de energía no mayor que
$\sqrt{\epsilon}$~\cite{MandaHillebrandSkokos2025GALI}; no se convierten esas
cifras en garantías universales.

\codepath{smaller_alignment_index},
\codepath{generalized_alignment_index} y
\codepath{linear_dependence_index} implementan el álgebra instantánea.
\codepath{alignment_indices_from_tangent_history} recibe exactamente la forma
\codepath{(n_samples, n_vectors, dimension)}. Las fachadas
\codepath{integer_flow_alignment_indices},
\codepath{integer_map_alignment_indices} e
\codepath{integer_system_alignment_indices} producen series completas y
conservan valores y logaritmos. Una máscara de censura distingue un volumen que
subdesbordó a cero de una familia de rango exactamente deficiente.
\codepath{AlignmentIndexResult} no duplica una serie LDI: para los mismos
vectores normalizados, el campo \codepath{gali} ya contiene la cantidad
equivalente.

NumPy/LAPACK SVD es la referencia. El backend Numba usa Householder sobre una
copia de cada matriz y acumula $\sum_j\log|R_{jj}|$; nunca usa ese $Q$ para
reiniciar las direcciones. No se añadió un kernel C para matrices pequeñas:
LAPACK ya es nativo y cruzar Python--C por muestra duplicaría interfaz sin
eliminar el costo dominante. Julia queda como oráculo por lotes opcional, no
como dependencia por paso. La formulación SVD y las tasas para mapas y flujos
se contrastaron también con el trabajo de 2026 de Rolim
Sales--Leonel--Antonopoulos~\cite{RolimSalesEtAl2026Alignment}.

El ejemplo Hamiltoniano ejecutable es
\codepath{examples/sali_gali_henon_heiles.py}. Usa
\[
 H=\tfrac12(p_x^2+p_y^2+x^2+y^2)+x^2y-\tfrac13y^3
\]
y su Jacobiano analítico para comparar, en una ventana finita, la ruta
variacional y la multiparticle:
\begin{verbatim}
python examples/sali_gali_henon_heiles.py
python benchmarks/bench_alignment_indices.py --output \
  validation/outputs/benchmarks/alignment_indices_numpy_numba_20260803.json
\end{verbatim}
La salida no adjudica la etiqueta ``caótico''. En particular, los índices
pueden comportarse de forma distinta en sistemas disipativos y no deben usarse
como clasificador universal aislado~\cite{MaLongZhu2016ChaosIndicators}.

En la ejecución local reproducible del ejemplo con duración 4, ambas rutas
conservaron una deriva relativa de energía de $8.42\times10^{-16}$; las máximas
diferencias variacional--multiparticle fueron $1.97\times10^{-8}$ para SALI y
$7.34\times10^{-9}$ para GALI. El benchmark caliente, sobre historiales
deterministas de 64, 512 y 4096 muestras, midió razones medianas
NumPy/Numba de 3.141, 2.122 y 2.018, respectivamente. La compilación de la
primera llamada Numba tomó 0.532 s; la peor discrepancia GALI fue
$1.56\times10^{-15}$ y la de log-volumen $3.55\times10^{-15}$. El SHA-256 del
JSON retenido fue
\codepath{FF4AEB083FDDB7594857F70DC841413CDDDBD36533CE31F6936CBD29BF2B13D8}.
Son cifras de ingeniería dependientes del host, no una superioridad universal.

El oráculo independiente
\codepath{validation/wolfram/cases/sali_gali_integer.wl} trabaja con 80 dígitos,
reconstruye MatrixPower/MatrixExp y compara Gram--Cauchy--Binet con SVD. Sus
fixtures son una rotación 3D, un mapa hiperbólico diagonal y un flujo diagonal
con soluciones cerradas. Prueban definiciones y propagación lineal finita; no
prueban caos no lineal, convergencia asintótica general, atracción u ocultedad.
El resumen promovido aprobó 12/12 comprobaciones; la comparación HAFO--Wolfram
alcanzó una diferencia global máxima de $1.7764\times10^{-15}$ y una diferencia
máxima de flujo de $8.8818\times10^{-16}$. Su SHA-256 es
\codepath{299CE6A9A0BB7A3C2CF920AE2F8F9C85A1A11D9244729D80302BDA4513ED6787}.

No existe en esta versión una fachada SALI/GALI/LDI fraccionaria. Reutilizar el
sistema tangente ODE no representa la perturbación completa de Caputo,
Riemann--Liouville, Grünwald--Letnikov, Caputo--Fabrizio,
Atangana--Baleanu u otro operador con memoria. Cada familia requiere una
ecuación variacional en espacio de historia y una regla de renormalización
publicada y verificada; por ello 
\codepath{q != 1} se rechaza y la capacidad permanece
\codepath{research_required}.

\subsection{Vectores covariantes de Lyapunov para orden entero}
\label{sec:clv-integer}

Los vectores covariantes de Lyapunov (CLV) describen direcciones del
desdoblamiento de Oseledets que son transportadas por el cociclo tangente. No
son, en general, ortogonales. Esta propiedad los distingue de los vectores de
Gram--Schmidt que produce el cálculo QR del espectro: estos últimos forman una
base numéricamente estable, mientras que los CLV recuperan direcciones
covariantes hacia delante y hacia atrás. HAFO implementa para $q=1$ el método
de Ginelli y colaboradores, junto con los contratos y controles numéricos
desarrollados en la literatura posterior
\cite{GinelliEtAl2007CLV,WolfeSamelson2007LyapunovVectors,
KuptsovParlitz2012CLV,FroylandEtAl2013CLVComparison,
GinelliEtAl2013CLVReview}.

Para un mapa $x_{n+1}=F(x_n)$ se define $M_n=DF(x_n)$. En un flujo,
$M_n=D\Phi_{\Delta t_n}(x(t_n))$ es el propagador tangente del segmento. La
pasada hacia delante aplica QR reducido con diagonal estrictamente positiva:
\[
 M_n Q_n=Q_{n+1}R_{n+1},
 \qquad Q_n^{\mathsf T}Q_n=I,
 \qquad R_{n+1,ii}>0.
\]
Esta convención elimina la ambigüedad de signo de cada factorización, aunque
la recta CLV final continúa siendo no orientada. Los exponentes de tiempo
finito se acumulan a partir de la diagonal,
\[
 \lambda_i^{(T)}=\frac{1}{T}
 \sum_n\log R_{n+1,ii},
\]
donde $T$ es la duración física observada en un flujo o el número de
iteraciones observadas en un mapa.

Sea $C_n$ una matriz triangular superior de coeficientes y
$V_n=Q_nC_n$ la matriz cuyas columnas son CLV. La covariancia proyectiva y la
identidad QR implican, columna a columna,
\[
 R_{n+1}C_n=C_{n+1}D_{n+1},
 \qquad
 C_n\ \propto\ R_{n+1}^{-1}C_{n+1},
 \qquad
 V_n=Q_nC_n,
\]
con $D_{n+1}$ diagonal y no singular que sólo fija las escalas. La pasada
hacia atrás resuelve el sistema triangular y normaliza cada columna después
de cada segmento; nunca forma una inversa explícita. Un horizonte futuro
adicional reduce la dependencia de la matriz terminal. De modo análogo, el
transitorio QR hacia delante permite que las subvariedades anidadas se
alineen antes de la ventana retenida. Ambos horizontes son parámetros
explícitos: la implementación todavía no decide automáticamente su
convergencia. El análisis reciente de eficiencia recomienda vigilar los
transitorios y advierte que una evolución hacia atrás excesiva puede degradar
subespacios centrales degenerados~\cite{DuPlessisHillebrandSkokos2026CLV}.

El contrato de arreglos evita confundir vectores por filas y por columnas:
\begin{itemize}
  \item las bases QR tienen forma
  \codepath{(samples, dimension, n_vectors)};
  \item los factores observados $R$ tienen forma
  \codepath{(samples - 1, n_vectors, n_vectors)} y los factores futuros
  \codepath{(future_segments, n_vectors, n_vectors)};
  \item la salida pública \codepath{vectors} usa
  \codepath{(samples, n_vectors, dimension)} y \codepath{coefficients} usa
  \codepath{(samples, n_vectors, n_vectors)}.
\end{itemize}
HAFO rechaza bases no ortonormales, factores que no sean triangulares
superiores con diagonal positiva y segmentos singulares o casi singulares.
También estima el espacio de trabajo antes de propagar y aplica por defecto
una guarda de 512 MiB. NumPy usa \codepath{scipy.linalg.solve_triangular}; el
backend Numba ejecuta la misma sustitución hacia atrás mediante bucles
compilados. Para flujos, la fachada integra conjuntamente
\[
 \dot x=f(x),\qquad \dot Y=J_f(x)Y
\]
con DOP853 y QR por segmentos. Para mapas usa la recurrencia exacta del
Jacobiano evaluado antes de actualizar el estado. En ambos casos se acepta un
Jacobiano analítico o, como respaldo declarado, diferencias centrales
relativas por componente.

El benchmark CLV separado generó historias QR deterministas de mapas lineales
no normales, exigió paridad NumPy--Numba antes de medir, separó el primer
llamado JIT y alternó los backends durante siete repeticiones en tres cargas.
Para 4096 segmentos observados, 1024 segmentos futuros, $d=8$ y cuatro
vectores, la pasada pública fue $64.42$ veces más rápida con Numba; en el
pipeline completo la razón fue $2.118$ y la reducción observada fue 52.8\%.
Sin embargo, la reconstrucción Numba ya calentada ocupó sólo 1.60\% del tiempo
end-to-end. Aun suponiendo coste cero para toda esa fase, el techo idealizado
de mejora sería apenas $1.0163$. No se implementa por ello un backend C para
esta pasada: primero se deberá perfilar una carga productiva y después exigir
paridad y mediciones C end-to-end repetidas. Esta conclusión es local al host
y cargas registrados, no una superioridad universal. El artefacto promovido
es
\codepath{validation/outputs/benchmarks/covariant_lyapunov_numpy_numba_20260803.json},
con SHA-256
\codepath{15A904AC4F67491015DF4A07CFC4353BBC5F2E19C4BC249121F415CCEE7F61BD}.

Para dos direcciones unitarias $\widehat v_i$ y $\widehat v_j$, la API
geométrica calcula el ángulo no orientado de sus rectas mediante
\[
 c_{ij}=\widehat v_i^{\mathsf T}\widehat v_j,
 \qquad
 \theta_{ij}=\operatorname{atan2}
 \left(\sqrt{\max(0,1-c_{ij}^{2})},|c_{ij}|\right)
 \in[0,\pi/2].
\]
La forma con \codepath{atan2} es más estable cerca de $0$ y $\pi/2$ que una
evaluación directa de \codepath{arccos}. Para grupos de CLV,
\codepath{covariant_lyapunov_angles} calcula el menor ángulo principal con
\codepath{scipy.linalg.subspace_angles} y puede devolver promedios móviles
declarados. Si los exponentes de tiempo finito son repetidos o casi
degenerados, se cierran las brechas espectrales que controlan la velocidad de
convergencia~\cite{Noethen2021CLVHilbert}; una columna CLV individual puede no
ser única o estar mal condicionada. En ese caso deben compararse los
subespacios covariantes, no forzarse una interpretación de direcciones
individuales.

Las fachadas públicas son
\codepath{integer_covariant_vectors_from_qr_history},
\codepath{integer_flow_covariant_lyapunov_vectors},
\codepath{integer_map_covariant_lyapunov_vectors},
\codepath{integer_system_covariant_lyapunov_vectors} y
\codepath{covariant_lyapunov_angles}. Por ejemplo, para el mapa clásico de
Hénon,
\[
 F(x,y)=(1-1.4x^2+y,\,0.3x),
 \qquad
 DF(x,y)=\begin{bmatrix}-2.8x&1\\0.3&0\end{bmatrix},
\]
se puede ejecutar:
\begin{verbatim}
from hidden_attractors import (
    covariant_lyapunov_angles,
    integer_map_covariant_lyapunov_vectors,
)

result = integer_map_covariant_lyapunov_vectors(
    henon_map, henon_map_jacobian, [0.0, 0.0],
    iterations=300,
    transient_iterations=400,
    forward_transient_iterations=250,
    backward_transient_iterations=250,
    n_vectors=2,
    seed=20260803,
    backend="auto",
    q=1.0,
)
angles = covariant_lyapunov_angles(
    result.vectors,
    coordinates=result.coordinates,
    pairs=((0, 1),),
    unoriented=True,
)
\end{verbatim}
El ejemplo reproducible completo se encuentra en
\codepath{examples/covariant_lyapunov_henon_map.py}. La ejecución cerrada con
300 iteraciones observadas, 400 de transitorio de estado y horizontes QR
hacia delante y hacia atrás de 250 iteraciones resolvió
\codepath{backend=auto} con Numba y produjo 301 muestras. Los exponentes de
tiempo finito por iteración fueron $[0.4178061335,-1.6217789378]$; el ángulo
entre las dos rectas tuvo mínimo, mediana y máximo de
$[0.0044154621,0.9948201635,1.5661450963]$ rad, y el máximo residuo de
covariancia proyectiva fue $2.1073424255\times10^{-8}$. Estas cifras describen
únicamente esa trayectoria y ventana finitas: no certifican caos,
hiperbolicidad, atracción ni ocultedad.

El oráculo independiente
\codepath{validation/wolfram/cases/covariant_lyapunov_integer.wl} trabaja con
80 dígitos, usa su propia ortogonalización de Gram--Schmidt modificada y no lee
el código HAFO, sus fórmulas ni el reporte. Verifica un mapa no normal 2D y un
flujo constante no normal 3D mediante soluciones exactas, recurrencia
triangular, covariancia proyectiva, exponentes y ángulos. Las 17/17
comprobaciones pasaron. La comparación obligatoria con el núcleo HAFO obtuvo
diferencias máximas de $3.0658\times10^{-16}$ en rectas,
$4.3133\times10^{-16}$ en covariancia y $3.3307\times10^{-16}$ en ángulos.
Los artefactos promovidos y sus SHA-256 son:
\begin{itemize}
  \item
  \codepath{validation/outputs/wolfram/covariant_lyapunov_integer_verified/covariant_lyapunov_integer_validation_summary.json}.
  \par\noindent SHA-256:
  \codepath{C2366BC0379B8C73863431CAC8B72742286869293DA6A5CC6C6428C5D9DDB211};
  \item
  \codepath{validation/outputs/wolfram/covariant_lyapunov_integer_verified/covariant_lyapunov_integer_comparison.json}.
  \par\noindent SHA-256:
  \codepath{24BFCFBECDEAC466F3CBFFF7F7AE4869792D507A6C53E09629D0975A2A729AF0}.
\end{itemize}
El alcance de este oráculo es álgebra finita de cociclos lineales
diagonalizables y consistencia entre implementaciones; no constituye un
teorema de convergencia CLV no lineal ni prueba caos, atracción u ocultedad.

La capacidad fraccionaria permanece deliberadamente
\codepath{research_required}. Incluso si se escribe formalmente, por ejemplo,
${}^{C}D_t^q\delta x=J_f(x(t))\delta x$, el sistema con memoria no induce en
general el cociclo finito y sin historia que presupone la recurrencia anterior.
Reiniciar, ortogonalizar o normalizar sólo $\delta x(t)$ puede borrar parte del
estado hereditario. Caputo, Riemann--Liouville, Grünwald--Letnikov,
Caputo--Fabrizio y Atangana--Baleanu además difieren en núcleo, prehistoria y
condiciones iniciales. Por ello las fachadas dinámicas exigen $q=1$ y rechazan
un orden fraccionario; los ángulos son sólo un postproceso geométrico y no
validan por sí mismos que un historial suministrado sea un historial CLV. La
extensión fraccionaria requiere definir y verificar primero el espacio de
historia, su operador tangente, la covariancia y una renormalización que no
destruya la memoria.

Comandos de reproducción del bloque:
\begin{verbatim}
python examples/covariant_lyapunov_henon_map.py \
  --iterations 300 --transient-iterations 400 \
  --forward-transient-iterations 250 \
  --backward-transient-iterations 250 --backend auto

python -m pytest tests/test_covariant_lyapunov.py \
  tests/test_covariant_lyapunov_wolfram.py \
  tests/test_covariant_lyapunov_example.py \
  tests/test_covariant_lyapunov_benchmark.py -q

python benchmarks/bench_covariant_lyapunov.py --repeats 7 \
  --output benchmark_outputs/hafo_covariant_lyapunov_benchmark.json
\end{verbatim}

\subsection{Comandos reproducibles para el análisis avanzado}

\mbox{}\par
\begin{verbatim}
cd version_2
python -m pytest tests/test_basin_uncertainty.py \
  tests/test_delay_embedding_advanced.py \
  tests/test_recurrence_advanced.py \
  tests/test_correlation_dimension.py \
  tests/test_native_correlation_sum.py \
  tests/test_correlation_dimension_example.py \
  tests/test_analysis_contracts.py \
  tests/test_permutation_entropy.py \
  tests/test_native_permutation_entropy.py \
  tests/test_permutation_entropy_wolfram.py \
  tests/test_permutation_entropy_example.py \
  tests/test_alignment_indices.py \
  tests/test_sali_gali_integer_wolfram.py \
  tests/test_sali_gali_henon_heiles_example.py \
  tests/test_covariant_lyapunov.py \
  tests/test_covariant_lyapunov_wolfram.py \
  tests/test_covariant_lyapunov_example.py \
  tests/test_covariant_lyapunov_benchmark.py -q
\end{verbatim}
El comando valida fórmulas y contratos sintéticos. No ejecuta por sí mismo una
campaña Chua, no demuestra fractalidad asintótica y no certifica caos,
atracción u ocultedad.


\section{Catálogo reproducible de gráficas}
\label{sec:catalogo-graficas}

Esta sección cubre todas las funciones públicas cuyo nombre comienza con
\codepath{plot_} o \codepath{render_} en
\codepath{hidden_attractors.plotting}: 33 funciones en la versión documentada.
Cada imagen fue producida llamando a la librería, no dibujada manualmente. Las
entradas tampoco son curvas pedagógicas ni matrices sintéticas construidas a
mano: proceden de cálculos numéricos reproducibles del sistema registrado
\codepath{chua-nonsmooth}. El caso
\codepath{chua_nonsmooth_real_system_catalog_20260729} es una reintegración
con la implementación canónica de la biblioteca.

El vector usado es
\[
(\alpha,\beta,\gamma,m_0,m_1,a_1,a_2,\rho)=
(8.4562,12.0732,0.0052,-0.1768,-1.1468,0.4,-1.5585,1).
\]
Según la salida, el contrato numérico corresponde a una trayectoria
\(q=1\) con \codepath{efork_q1}, una continuación entera en \(\lambda\), un
barrido de \(\beta\) con RK4, exponentes finito-temporales con QR--Benettin,
controles integrados desde vecindades de los equilibrios o una clasificación
de cuenca por tiempo finito. El manifiesto conserva para cada insumo el paso,
horizonte, condición inicial, malla, umbrales y semilla exactos.

Estos son resultados numéricos reales de un sistema matemático definido, no
mediciones experimentales ni una nueva campaña de validación. Ninguna imagen
certifica por sí sola caos, estabilidad asintótica, ocultedad, convergencia de
un integrador o desempeño científico.

El catálogo completo se regenera desde la raíz de \codepath{version_2} con:
\begin{verbatim}
python figure_scripts/generate_plot_catalog_examples.py
\end{verbatim}
El archivo
\codepath{docs/assets/generated_plot_catalog/catalog_results.json} conserva
la firma inspeccionada, el comando individual, la imagen representativa y el
estado de cada función, así como las 62 imágenes y sus hashes SHA-256. El
generador compara su lista con
\codepath{hidden_attractors.plotting.__all__} y falla si la API y el catálogo
dejan de coincidir.

\newcommand{\plotcatalogentry}[4]{%
  \bigskip
  \noindent\textbf{Función.}\ \codepath{#1}\par
  \noindent #2\par
  \noindent\textbf{Llamada directa.}\ \codepath{#3}\par
  \noindent\textbf{Comando del ejemplo.}\
  \texttt{python}\ \codepath{figure_scripts/generate_plot_catalog_examples.py}\
  \texttt{--only}\ \codepath{#1}\par
  \begin{center}
    \includegraphics[
      width=0.82\textwidth,
      height=0.48\textheight,
      keepaspectratio
    ]{#4}\\[-1mm]
    {\footnotesize Sistema \codepath{chua-nonsmooth} con el vector de
    parámetros declarado arriba. La entrada procede de la reintegración
    canónica y del contrato específico registrado en
    \codepath{catalog_results.json}. Es una salida numérica reproducible de la
    API, no una medición ni nueva evidencia de validación; por sí sola no
    demuestra caos, estabilidad, ocultedad o convergencia.}
  \end{center}
}

\newcommand{\plotcatalogextra}[3]{%
  \medskip
  \noindent\textbf{#2} (\codepath{#1}).\par
  \begin{center}
    \includegraphics[
      width=0.76\textwidth,
      height=0.39\textheight,
      keepaspectratio
    ]{#3}\\[-1mm]
    {\footnotesize Salida adicional calculada para
    \codepath{chua-nonsmooth}; parámetros, contrato numérico, procedencia y
    hash se registran en \codepath{catalog_results.json}. No constituye una
    medición experimental, nueva validación ni prueba autónoma de una
    propiedad dinámica.}
  \end{center}
}

\subsection{Trayectorias, series y espectros}

\plotcatalogentry
  {plot_phase_space}
  {Representa una trayectoria en dos o tres coordenadas, con color temporal
  opcional. La entrada usual tiene columnas \((t,x,y,z)\).}
  {plot_phase_space(trajectory, output_path, dims=("x","y","z"))}
  {assets/generated_plot_catalog/examples/08_plot_phase_space.png}

\plotcatalogentry
  {plot_phase_projections}
  {Construye las proyecciones \(x\)-\(y\), \(x\)-\(z\) y \(y\)-\(z\) de una
  misma trayectoria.}
  {plot_phase_projections(trajectory, output_path)}
  {assets/generated_plot_catalog/examples/07_plot_phase_projections.png}

\plotcatalogentry
  {plot_time_series}
  {Grafica las componentes seleccionadas contra el tiempo de muestreo.}
  {plot_time_series(trajectory, output_path, columns=("x","y","z"))}
  {assets/generated_plot_catalog/examples/10_plot_time_series.png}

\plotcatalogentry
  {plot_trajectory_overlay}
  {Superpone dos o más trayectorias etiquetadas para una comparación visual
  bajo la misma proyección.}
  {plot_trajectory_overlay(trajectories, labels, title="Comparacion", output_path=output_path)}
  {assets/generated_plot_catalog/examples/12_plot_trajectory_overlay.png}

\plotcatalogentry
  {plot_spectrum}
  {Representa un \codepath{SpectrumResult} calculado previamente, con
  frecuencia en las unidades declaradas.}
  {plot_spectrum(spectrum, output_path, x_units="Hz")}
  {assets/generated_plot_catalog/examples/09_plot_spectrum.png}

\plotcatalogentry
  {plot_trajectory_spectra}
  {Calcula y exporta el espectro de las componentes de una trayectoria. Acepta
  FFT y los alias \codepath{psd}, \codepath{welch} y
  \codepath{psd_welch}.}
  {plot_trajectory_spectra(trajectory, output_dir, method="fft", prefix="case")}
  {assets/generated_plot_catalog/examples/11_plot_trajectory_spectra.png}

\plotcatalogentry
  {plot_lyapunov_convergence}
  {Muestra la evolución finito-temporal de los exponentes contenidos en un
  \codepath{LyapunovResult}; no estima los exponentes por sí sola.}
  {plot_lyapunov_convergence(result, output_path)}
  {assets/generated_plot_catalog/examples/03_plot_lyapunov_convergence.png}

\plotcatalogentry
  {plot_attractor_trajectories}
  {Genera vistas de una trayectoria y de los equilibrios suministrados usando
  el contrato de configuración del workflow.}
  {plot_attractor_trajectories(trajectory, equilibria, config, output_dir)}
  {assets/generated_plot_catalog/examples/24_plot_attractor_trajectories.png}

\plotcatalogentry
  {plot_flexible_attractor_and_projections}
  {Exporta una vista espacial y proyecciones configurables para un caso
  identificado.}
  {plot_flexible_attractor_and_projections(trajectory, equilibria, config, output_dir, "case")}
  {assets/generated_plot_catalog/examples/25_plot_flexible_attractor_and_projections.png}

\plotcatalogentry
  {plot_timeseries_data}
  {Exporta las series temporales configuradas y el CSV correspondiente con
  columnas \(t,x,y,z\).}
  {plot_timeseries_data(trajectory, config, output_dir, "case")}
  {assets/generated_plot_catalog/examples/26_plot_timeseries_data.png}

\subsection{Bifurcación, cuencas y estabilidad}

\plotcatalogentry
  {plot_bifurcation_diagram}
  {Representa puntos ya postprocesados contra el parámetro; no continúa ramas
  ni identifica automáticamente bifurcaciones.}
  {plot_bifurcation_diagram(points, output_path, parameter_label="p", observable_label="x")}
  {assets/generated_plot_catalog/examples/06_plot_bifurcation_diagram.png}

\plotcatalogentry
  {plot_basin_slices}
  {Representa una o varias matrices de clases sobre planos declarados. Los
  colores reflejan las etiquetas recibidas, no una prueba global de cuenca.}
  {plot_basin_slices(slices, "system_id", output_dir)}
  {assets/generated_plot_catalog/examples/13_plot_basin_slices.png}

\plotcatalogentry
  {plot_basin_slice_file}
  {Exporta un único corte de cuenca a partir de dos ejes y una matriz de
  clases.}
  {plot_basin_slice_file("xy", u, v, classes, "E0", "system_id", output_dir)}
  {assets/generated_plot_catalog/examples/14_plot_basin_slice_file.png}

\plotcatalogentry
  {plot_matignon_equilibria}
  {Ubica autovalores de los Jacobianos de los equilibrios respecto de la
  frontera angular de Matignon para el orden \(q\).}
  {plot_matignon_equilibria(system, equilibria, 0.9, output_dir)}
  {assets/generated_plot_catalog/examples/15_plot_matignon_equilibria.png}

\subsection{Transferencia de Lur'e y función descriptiva}

\plotcatalogentry
  {plot_lure_nyquist_describing_function}
  {Compara el lugar de Nyquist de \(W_q\) con la curva
  \(-1/N(A)\) y marca una semilla armónica suministrada.}
  {plot_lure_nyquist_describing_function(system.lure, seed, output_path, q=1.0)}
  {assets/generated_plot_catalog/examples/04_plot_lure_nyquist_describing_function.png}

\plotcatalogentry
  {plot_lure_transfer_components}
  {Representa componentes reales, imaginarias y magnitudes de la transferencia
  de Lur'e alrededor de una semilla.}
  {plot_lure_transfer_components(system.lure, seed, output_path, q=1.0)}
  {assets/generated_plot_catalog/examples/05_plot_lure_transfer_components.png}

\plotcatalogentry
  {plot_nyquist_transfer}
  {Representa valores complejos de transferencia sobre una malla de
  frecuencias y candidatos \((A,\omega,k)\).}
  {plot_nyquist_transfer(omega, W, candidates, config, output_dir)}
  {assets/generated_plot_catalog/examples/16_plot_nyquist_transfer.png}

\plotcatalogentry
  {plot_describing_function}
  {Representa la función descriptiva del sistema y los candidatos entregados
  por el usuario.}
  {plot_describing_function(system, candidates, config, output_dir)}
  {assets/generated_plot_catalog/examples/17_plot_describing_function.png}

\plotcatalogentry
  {plot_harmonic_residual_map}
  {Construye el mapa del residuo
  \(\lvert1+N(A)W(\mathrm{i}\omega)\rvert\). Un mínimo visual no reemplaza la
  verificación numérica de cierre.}
  {plot_harmonic_residual_map(system, candidates, config, output_dir)}
  {assets/generated_plot_catalog/examples/18_plot_harmonic_residual_map.png}

El ejemplo real de estas gráficas declara
\codepath{transfer_convention="opposite_sign"} y
\codepath{harmonic_condition="1_plus_WN"}; usa entonces
\[
W_{\mathrm{code}}(s)=c^\mathsf{T}(P-sI)^{-1}b,\qquad
\left|1+N(A)W_{\mathrm{code}}(\mathrm{i}\omega)\right|,
\qquad W_{\mathrm{code}}(\mathrm{i}\omega_0)=-\frac{1}{k}.
\]
La pareja normalizada también aceptada por la API es
\codepath{standard + 1_minus_WN}. Como
\(W_{\mathrm{report}}=-W_{\mathrm{code}}\), en ella se representa
\(\lvert1-N(A)W_{\mathrm{report}}\rvert\) y el cierre queda en
\(W_{\mathrm{report}}=1/k\). Las funciones infieren la clave complementaria
si se proporciona solamente una y rechazan parejas de signo incoherentes,
salvo habilitación explícita.

\subsection{Continuación y controles de vecindad}

\plotcatalogentry
  {plot_integer_lure_continuation}
  {Representa los pasos de una continuación de Lur'e de orden entero ya
  calculada.}
  {plot_integer_lure_continuation(steps, output_path)}
  {assets/generated_plot_catalog/examples/02_plot_integer_lure_continuation.png}

\plotcatalogentry
  {plot_integer_hiddenness_controls}
  {Compara la trayectoria objetivo con controles enteros suministrados desde
  vecindades de equilibrios.}
  {plot_integer_hiddenness_controls(trajectory, probes, output_path)}
  {assets/generated_plot_catalog/examples/01_plot_integer_hiddenness_controls.png}

\plotcatalogentry
  {plot_continuation_eta}
  {Exporta norma final y amplitud pico a pico contra el parámetro de
  continuación \(\lambda\), además de comparaciones cuando hay varios pasos.
  El nombre público y los archivos con sufijo \codepath{_eta} se conservan
  únicamente por compatibilidad.}
  {plot_continuation_eta(cont_steps, config, output_dir)}
  {assets/generated_plot_catalog/examples/19_plot_continuation_eta.png}

\plotcatalogentry
  {plot_continuation_first_last_comparison}
  {Compara geométricamente las trayectorias del primer y último paso
  disponibles.}
  {plot_continuation_first_last_comparison(cont_steps, config, output_dir)}
  {assets/generated_plot_catalog/examples/20_plot_continuation_first_last_comparison.png}

\plotcatalogentry
  {plot_continuation_timeseries_comparison}
  {Compara las series temporales del primer y último paso de la continuación.}
  {plot_continuation_timeseries_comparison(cont_steps, config, output_dir)}
  {assets/generated_plot_catalog/examples/21_plot_continuation_timeseries_comparison.png}

\plotcatalogentry
  {plot_continuation_progression}
  {Muestra la progresión de estados de entrada, salida y trayectorias a lo
  largo de los pasos suministrados.}
  {plot_continuation_progression(cont_steps, config, output_dir)}
  {assets/generated_plot_catalog/examples/22_plot_continuation_progression.png}

\plotcatalogentry
  {plot_continuation_tracking}
  {Exporta por separado la norma rastreada y el estado discreto de cada paso.}
  {plot_continuation_tracking(cont_steps, config, output_dir)}
  {assets/generated_plot_catalog/examples/23_plot_continuation_tracking.png}

\plotcatalogentry
  {plot_neighborhood_control_spheres}
  {Representa la trayectoria, equilibrios y procedencia de sondas sobre
  esferas de control declaradas.}
  {plot_neighborhood_control_spheres(trajectory, probes, equilibria, config, output_dir)}
  {assets/generated_plot_catalog/examples/27_plot_neighborhood_control_spheres.png}

\plotcatalogentry
  {plot_sphere_test_results}
  {Resume los registros de una prueba de esfera para un equilibrio y radio
  específicos.}
  {plot_sphere_test_results("E0", equilibrium, radius, probe_runs, output_dir)}
  {assets/generated_plot_catalog/examples/28_plot_sphere_test_results.png}

\subsection{Renderizadores unificados}

\plotcatalogentry
  {render_attractor}
  {Renderiza una trayectoria y sus equilibrios en el almacén de figuras
  configurado, junto con metadatos.}
  {render_attractor(trajectory, equilibria, config, run_id="case")}
  {assets/generated_plot_catalog/examples/29_render_attractor.png}

\plotcatalogentry
  {render_basin}
  {Renderiza una matriz de cuenca y sus ejes en PNG, PDF y metadatos JSON.}
  {render_basin(grid_x, grid_y, basin_grid, config, run_id="case")}
  {assets/generated_plot_catalog/examples/30_render_basin.png}

\plotcatalogentry
  {render_nyquist}
  {Renderiza transferencia, función descriptiva y candidatos complejos con
  metadatos de procedencia.}
  {render_nyquist(freqs, W, N, candidates, config, run_id="case")}
  {assets/generated_plot_catalog/examples/31_render_nyquist.png}

\plotcatalogentry
  {render_matignon}
  {Renderiza autovalores y frontera de estabilidad para el orden declarado.}
  {render_matignon(eigenvalues, q, config, run_id="case")}
  {assets/generated_plot_catalog/examples/32_render_matignon.png}

\plotcatalogentry
  {render_all_plots}
  {Orquesta cualquier subconjunto compatible de atractor, cuenca, Nyquist y
  Matignon; no define un quinto diagnóstico matemático.}
  {render_all_plots(trajectory=trajectory, equilibria=equilibria, config=config, run_id="case")}
  {assets/generated_plot_catalog/examples/33_render_all_plots.png}

\subsection{Registro visual completo de funciones multisalida}

Las fichas anteriores conservan la correspondencia uno a uno entre las 33
funciones, sus llamadas directas y sus comandos individuales. Diez funciones
producen más de una imagen. El mismo comando
\codepath{python figure_scripts/generate_plot_catalog_examples.py --only FUNCTION_NAME}
indicado en cada ficha regenera su imagen principal y las salidas
adicionales siguientes. Con ellas, el inventario completo contiene 62 PNG:
33 representativos y 29 adicionales.

\plotcatalogextra
  {plot_trajectory_spectra}
  {Espectro de la componente 1}
  {assets/generated_plot_catalog/examples/11_plot_trajectory_spectra__02_catalog_fft_component_1.png}

\plotcatalogextra
  {plot_trajectory_spectra}
  {Espectro de la componente 2}
  {assets/generated_plot_catalog/examples/11_plot_trajectory_spectra__03_catalog_fft_component_2.png}

\plotcatalogextra
  {plot_nyquist_transfer}
  {Ampliación del entorno de cierre de Nyquist}
  {assets/generated_plot_catalog/examples/16_plot_nyquist_transfer__02_fig01b_nyquist_zoom_x.png}

\plotcatalogextra
  {plot_nyquist_transfer}
  {Partes real e imaginaria de la transferencia}
  {assets/generated_plot_catalog/examples/16_plot_nyquist_transfer__03_transfer_real_imag.png}

\plotcatalogextra
  {plot_continuation_eta}
  {Amplitud pico a pico contra \(\lambda\)}
  {assets/generated_plot_catalog/examples/19_plot_continuation_eta__02_continuation_amplitude_vs_eta.png}

\plotcatalogextra
  {plot_continuation_eta}
  {Comparación espacial del primer y último paso}
  {assets/generated_plot_catalog/examples/19_plot_continuation_eta__03_continuation_first_last_comparison.png}

\plotcatalogextra
  {plot_continuation_eta}
  {Proyecciones del primer y último paso}
  {assets/generated_plot_catalog/examples/19_plot_continuation_eta__04_continuation_first_last_projections.png}

\plotcatalogextra
  {plot_continuation_eta}
  {Comparación de series temporales}
  {assets/generated_plot_catalog/examples/19_plot_continuation_eta__05_continuation_timeseries_comparison_x.png}

\plotcatalogextra
  {plot_continuation_eta}
  {Progresión de la continuación}
  {assets/generated_plot_catalog/examples/19_plot_continuation_eta__06_continuation_progression.png}

\plotcatalogextra
  {plot_continuation_first_last_comparison}
  {Proyecciones del primer y último paso}
  {assets/generated_plot_catalog/examples/20_plot_continuation_first_last_comparison__02_continuation_first_last_projections.png}

\plotcatalogextra
  {plot_continuation_tracking}
  {Estado discreto por paso}
  {assets/generated_plot_catalog/examples/23_plot_continuation_tracking__02_continuation_tracking_status.png}

\plotcatalogextra
  {plot_attractor_trajectories}
  {Proyección \(xy\) de la trayectoria}
  {assets/generated_plot_catalog/examples/24_plot_attractor_trajectories__02_attractor_xy.png}

\plotcatalogextra
  {plot_attractor_trajectories}
  {Proyección \(xz\) de la trayectoria}
  {assets/generated_plot_catalog/examples/24_plot_attractor_trajectories__03_attractor_xz.png}

\plotcatalogextra
  {plot_attractor_trajectories}
  {Proyección \(yz\) de la trayectoria}
  {assets/generated_plot_catalog/examples/24_plot_attractor_trajectories__04_attractor_yz.png}

\plotcatalogextra
  {plot_flexible_attractor_and_projections}
  {Vista flexible en el plano \(xy\)}
  {assets/generated_plot_catalog/examples/25_plot_flexible_attractor_and_projections__02_catalog_flexible_xy.png}

\plotcatalogextra
  {plot_flexible_attractor_and_projections}
  {Vista flexible en el plano \(xz\)}
  {assets/generated_plot_catalog/examples/25_plot_flexible_attractor_and_projections__03_catalog_flexible_xz.png}

\plotcatalogextra
  {plot_flexible_attractor_and_projections}
  {Vista flexible en el plano \(yz\)}
  {assets/generated_plot_catalog/examples/25_plot_flexible_attractor_and_projections__04_catalog_flexible_yz.png}

\plotcatalogextra
  {plot_timeseries_data}
  {Serie temporal de la componente \(y\)}
  {assets/generated_plot_catalog/examples/26_plot_timeseries_data__02_catalog_series_timeseries_y.png}

\plotcatalogextra
  {plot_timeseries_data}
  {Serie temporal de la componente \(z\)}
  {assets/generated_plot_catalog/examples/26_plot_timeseries_data__03_catalog_series_timeseries_z.png}

\plotcatalogextra
  {plot_timeseries_data}
  {Series temporales conjuntas \(x,y,z\)}
  {assets/generated_plot_catalog/examples/26_plot_timeseries_data__04_catalog_series_timeseries_xyz.png}

\plotcatalogextra
  {render_attractor}
  {Renderizado del atractor en el plano \(xy\)}
  {assets/generated_plot_catalog/examples/29_render_attractor__02_chua-nonsmooth_attractor_xy.png}

\plotcatalogextra
  {render_attractor}
  {Renderizado del atractor en el plano \(xz\)}
  {assets/generated_plot_catalog/examples/29_render_attractor__03_chua-nonsmooth_attractor_xz.png}

\plotcatalogextra
  {render_attractor}
  {Renderizado del atractor en el plano \(yz\)}
  {assets/generated_plot_catalog/examples/29_render_attractor__04_chua-nonsmooth_attractor_yz.png}

\plotcatalogextra
  {render_all_plots}
  {Renderizado conjunto del atractor en \(xy\)}
  {assets/generated_plot_catalog/examples/33_render_all_plots__02_chua-nonsmooth_attractor_xy.png}

\plotcatalogextra
  {render_all_plots}
  {Renderizado conjunto del atractor en \(xz\)}
  {assets/generated_plot_catalog/examples/33_render_all_plots__03_chua-nonsmooth_attractor_xz.png}

\plotcatalogextra
  {render_all_plots}
  {Renderizado conjunto del atractor en \(yz\)}
  {assets/generated_plot_catalog/examples/33_render_all_plots__04_chua-nonsmooth_attractor_yz.png}

\plotcatalogextra
  {render_all_plots}
  {Renderizado conjunto de la cuenca finito-temporal}
  {assets/generated_plot_catalog/examples/33_render_all_plots__05_chua-nonsmooth_basin.png}

\plotcatalogextra
  {render_all_plots}
  {Renderizado conjunto de Nyquist}
  {assets/generated_plot_catalog/examples/33_render_all_plots__06_chua-nonsmooth_nyquist.png}

\plotcatalogextra
  {render_all_plots}
  {Renderizado conjunto de Matignon}
  {assets/generated_plot_catalog/examples/33_render_all_plots__07_chua-nonsmooth_matignon.png}

\noindent
El nombre \codepath{plot_continuation_eta} y los sufijos
\codepath{_eta} de algunas rutas se mantienen sólo por compatibilidad. Las
etiquetas visibles usan el parámetro efectivo \(\lambda\). El manifiesto
generado exige 33 resultados exitosos y enumera cada una de estas 62 imágenes
en \codepath{produced_outputs}.


\section{API de la librería}

\subsection{Importaciones principales}

Para un usuario nuevo, el patrón recomendado es importar desde
\codepath{hidden_attractors} cuando el símbolo esté reexportado:

\begin{verbatim}
from hidden_attractors import (
    ChuaParameters,
    chua_nonsmooth_parameters,
    equilibria_nonsmooth,
    rhs_nonsmooth,
    get_system,
    register_system,
    find_lure_harmonic_seed,
    compute_trajectory_metrics,
    compute_boundedness_metrics,
    compute_fft_psd,
    detect_poincare_crossings,
    zero_one_test,
    estimate_time_series_lyapunov,
    load_trajectory_csv,
)
\end{verbatim}

Las extensiones experimentales se importan desde sus módulos para que el nivel
de estabilidad y la semántica del cálculo permanezcan visibles:

\begin{verbatim}
from hidden_attractors.fractional import (
    FractionalProblem,
    solve_fractional_problem,
    solve_fractional_system,
    lubich_convolution_quadrature,
    hadamard_convolution_quadrature,
    integrate_caputo_hadamard_abm,
)
from hidden_attractors.analysis.basin_uncertainty import basin_entropy
from hidden_attractors.analysis.delay_embedding import generalized_delay_embedding
from hidden_attractors.analysis.recurrence_advanced import auto_recurrence_matrix
\end{verbatim}

\begin{center}
\footnotesize
\begin{longtable}{@{}L{0.24\textwidth}L{0.21\textwidth}L{0.31\textwidth}L{0.16\textwidth}@{}}
\toprule
Grupo & Símbolos típicos & Fundamento matemático & Nivel \\
\midrule
Modelos Chua & \codepath{ChuaParameters}, \codepath{rhs_nonsmooth}, \codepath{rhs_arctan} & Campo vectorial de Chua, no linealidad por saturación o arctan. & \STABLE\\
Equilibrios y Jacobianos & \codepath{equilibria_nonsmooth}, \codepath{jacobian_nonsmooth}, \codepath{equilibria_arctan}, \codepath{jacobian_arctan} & Ecuaciones algebraicas \(F(E)=0\), linealización y Matignon. & \STABLE\\
Registro de sistemas & \codepath{ChaoticSystem}, \codepath{LureSystem}, \codepath{register_system}, \codepath{get_system} & Contrato para campo, parámetros, equilibrios, Jacobiano y forma de Lur'e. & \STABLE\\
Cuencas & \codepath{CLASS_LABELS}, \codepath{TARGET_CLASS_IDS}, \codepath{is_target_class} & Etiquetas finitas de destino y regla \TARGET{}. & \STABLE\\
IO portable & \codepath{load_trajectory_csv} & Lectura reproducible de trayectorias CSV sin depender del árbol de validación. & \STABLE\\
Semillas & \codepath{HarmonicSeed}, \codepath{find_lure_harmonic_seed}, \codepath{find_omega_gain_candidates} & Balance armónico, Nyquist y función descriptiva. & \EXPERIMENTAL\\
Núcleo fraccionario & \codepath{FractionalProblem}; \codepath{solve_fractional_system}; \codepath{lubich_convolution_quadrature}; \codepath{hadamard_convolution_quadrature}; \codepath{integrate_caputo_hadamard_abm} & Definición, discretización, condición inicial, memoria y backend registrados por separado. Operadores muestreados y solvers no se intercambian. & \EXPERIMENTAL\\
Análisis & \codepath{compute_trajectory_metrics}, \codepath{compute_fft_psd}, \codepath{zero_one_test}, \codepath{detect_poincare_crossings}, \codepath{integer_system_lyapunov_exponents}, \codepath{estimate_time_series_lyapunov} & Caracterización finito-temporal de estados y señales, independiente de una búsqueda de ocultedad. & \EXPERIMENTAL\\
Análisis avanzado & \codepath{basin_entropy}; \codepath{uncertainty_fraction}; \codepath{generalized_delay_embedding}; \codepath{false_nearest_neighbors}; \codepath{auto_recurrence_matrix}; RQA & Incertidumbre de cuencas, reconstrucción empírica y recurrencia sobre datos muestreados enteros o fraccionarios. & \EXPERIMENTAL\\
Plotting & \codepath{plot_phase_space}, \codepath{plot_time_series}, \codepath{plot_lure_nyquist_describing_function} & Visualización de trayectorias, transferencia, cuencas y diagnósticos. & \EXPERIMENTAL\\
Workflows & \codepath{NumericalContract}, \codepath{ContinuationPlan}, \codepath{WorkflowInputSpec} & Protocolos reproducibles, especificación de solver y evidencia por etapas. & \EXPERIMENTAL\\
Native & \begin{tabular}[t]{@{}l@{}}\codepath{FractionalChuaBackend}\\\codepath{BasinBackend}\end{tabular} & Integración y clasificación C Chua-específicas. & \INTERNAL\\
\bottomrule
\end{longtable}
\end{center}

\subsection{Lyapunov reconstruido desde una serie temporal escalar}

Cuando sólo se dispone de una señal escalar uniformemente muestreada, la API
\codepath{estimate_time_series_lyapunov} usa el backend opcional
\codepath{nolds} sin requerir el campo vectorial, el Jacobiano ni un
integrador. La salida estructurada contiene el mayor exponente estimado por el
método de Rosenstein, un espectro finito reconstruido por Eckmann, la dimensión
de Kaplan--Yorke, unidades, parámetros, versión del backend, diagnóstico del
ajuste y advertencias.

\begin{verbatim}
from hidden_attractors import estimate_time_series_lyapunov

result = estimate_time_series_lyapunov(
    signal,
    sample_interval=0.01,
    observable="x",
    random_seed=0,
)
print(result.largest_exponent)
print(result.spectrum)
print(result.kaplan_yorke_dimension)
\end{verbatim}

El intervalo de muestreo se pasa como \codepath{tau}, de modo que los
exponentes quedan expresados en unidades de tiempo inversas. El ajuste
polinomial es el predeterminado. Si se solicita RANSAC, la semilla se fija
localmente bajo un bloqueo que protege el estado aleatorio global. La
implementación limita de forma explícita la memoria cuadrática de la matriz de
distancias de Rosenstein. Estas cantidades dependen de la ventana, el muestreo,
la reconstrucción y los parámetros del ajuste: son diagnósticos de datos
finitos, no certificados de caos, espectros asintóticos ni ocultedad.

\subsection{Cómo leer una función}

Una función de la librería debe leerse con cuatro preguntas:

\begin{enumerate}
  \item \textbf{Qué objeto matemático calcula.} Por ejemplo,
  \codepath{jacobian_nonsmooth} calcula una linealización regional, no una
  prueba global.
  \item \textbf{Qué entrada exige.} Por ejemplo, un integrador necesita \(q\),
  \(h\), \(t_f\), memoria y condición inicial.
  \item \textbf{Qué salida produce.} Por ejemplo, una semilla armónica produce
  \(A\), \(\omega\), ganancia y estado inicial.
  \item \textbf{Qué afirmación permite.} Por ejemplo, \codepath{is_target_class}
  permite contar contactos \TARGET{} bajo un muestreo, no probar una cuenca
  global exacta.
\end{enumerate}

\subsection{Inventario exhaustivo de la superficie pública}
\label{sec:api-publica-exhaustiva}

La superficie canónica importable desde \codepath{hidden_attractors} contiene
99 símbolos. La tabla siguiente los enumera todos, sin convertir constantes,
tipos de datos o envolturas de flujo en procedimientos matemáticos que no son.
Las 33 funciones públicas de \codepath{hidden_attractors.plotting} se
documentan por separado, una por una, en la
Sección~\ref{sec:catalogo-graficas}. Las rutas especializadas de análisis e
integración que no se reexportan en la raíz se describen en sus secciones
temáticas del presente informe. El módulo experimental
\codepath{hidden_attractors.fractional} se documenta función por función en la
Sección~\ref{sec:nucleo-numerico-fraccionario}; entropía de cuencas, delay y
recurrencia avanzados se mantienen como imports module-qualified y se describen
en las Secciones~\ref{sec:basin-entropy-uncertainty},
\ref{sec:delay-embedding-advanced} y \ref{sec:recurrence-advanced}.

\begin{center}
\footnotesize
\begin{longtable}{@{}L{0.16\textwidth}L{0.50\textwidth}L{0.28\textwidth}@{}}
\toprule
Grupo & Símbolos públicos reexportados & Papel y fundamento \\
\midrule
\endfirsthead
\toprule
Grupo & Símbolos públicos reexportados & Papel y fundamento \\
\midrule
\endhead
Estabilidad de API &
\codepath{EXPERIMENTAL},
\codepath{INTERNAL},
\codepath{LEGACY},
\codepath{STABLE},
\codepath{api_tier},
\codepath{assert_tier},
\codepath{get_tier},
\codepath{PUBLIC_API_STABLE},
\codepath{PUBLIC_API_EXPERIMENTAL},
\codepath{PUBLIC_API_TIERS}
&
Metadatos y comprobaciones del contrato de estabilidad. No calculan una
propiedad dinámica. \\
\addlinespace

Modelo Chua &
\codepath{ChuaParameters},
\codepath{chua_parameters},
\codepath{chua_arctan_wu2023_parameters},
\codepath{chua_nonsmooth_parameters},
\codepath{equilibria_arctan},
\codepath{equilibria_nonsmooth},
\codepath{jacobian_arctan},
\codepath{jacobian_nonsmooth},
\codepath{rhs_arctan},
\codepath{rhs_nonsmooth}
&
Parámetros, campos \(F(X)\), equilibrios \(F(E)=0\) y Jacobianos
\(J=\partial F/\partial X\) de los modelos definidos en las Secciones 5 y 6. \\
\addlinespace

Registro de sistemas &
\codepath{ChaoticSystem},
\codepath{LureSystem},
\codepath{check_system_capability},
\codepath{get_system},
\codepath{known_workflows},
\codepath{list_systems},
\codepath{register_system},
\codepath{requirements_for}
&
Representación reusable del campo autónomo y, cuando aplica, del
desdoblamiento \(PX+b\psi(r^\top X)\); además valida capacidades exigidas por
cada flujo. \\
\addlinespace

Clases de destino e I/O &
\codepath{CLASS_LABELS},
\codepath{TARGET_CLASS_IDS},
\codepath{class_label},
\codepath{is_target_class},
\codepath{load_trajectory_csv}
&
Convenciones finitas de clasificación y carga portable de trayectorias. Una
clase de destino es una etiqueta bajo el contrato declarado, no una prueba
global de cuenca. \\
\addlinespace

Análisis y resultados &
\codepath{BifurcationPoint},
\codepath{LyapunovComputationRequest},
\codepath{LyapunovComputationSummary},
\codepath{LyapunovResult},
\codepath{PoincareCrossingResult},
\codepath{RobustnessCase},
\codepath{SpectrumResult},
\codepath{TimeSeriesLyapunovResult},
\codepath{bifurcation_points_from_trajectories},
\codepath{bifurcation_summary},
\codepath{compute_boundedness_metrics},
\codepath{compute_fft_psd},
\codepath{compute_lyapunov_spectrum},
\codepath{compute_trajectory_metrics},
\codepath{detect_poincare_crossings},
\codepath{estimate_time_series_lyapunov},
\codepath{integer_qr_benettin_lyapunov_exponents},
\codepath{integer_system_lyapunov_exponents},
\codepath{kaplan_yorke_dimension},
\codepath{trajectory_metrics},
\codepath{trajectory_metrics_for_system},
\codepath{validate_lyapunov_method_request},
\codepath{zero_one_test}
&
Tipos estructurados y funciones de caracterización independientes. Sus
ecuaciones, discretizaciones, unidades, llamadas y límites se presentan en la
Sección 12. \\
\addlinespace

Semillas armónicas &
\codepath{HarmonicSeed},
\codepath{find_harmonic_seed},
\codepath{find_lure_harmonic_seed},
\codepath{find_lure_omega_gain_candidates},
\codepath{find_omega_gain_candidates},
\codepath{validate_fractional_order}
&
Solución del cierre de balance armónico y validación del orden. Son
constructores de semillas; no certifican por sí mismos caos ni ocultedad. \\
\addlinespace

Contratos de flujo &
\codepath{BasinSliceSpec},
\codepath{DestinationClassifierSpec},
\codepath{FullWorkflowContract},
\codepath{ContinuationPlan},
\codepath{ContinuationTrace},
\codepath{DynamicReference},
\codepath{FINAL_LABELS},
\codepath{HiddennessTestResult},
\codepath{IntegratorSpec},
\codepath{NumericalContract},
\codepath{OFFICIAL_STAGE_ORDER},
\codepath{ParameterSweepSpec},
\codepath{RobustnessCaseSpec},
\codepath{RobustnessVerdict},
\codepath{PROTOCOL_VERSION},
\codepath{PostContinuationDecision},
\codepath{SEED_FAMILIES},
\codepath{SoftPrecheckResult},
\codepath{SphereControlSpec},
\codepath{StageEnvelope},
\codepath{StrictRefinementSpec},
\codepath{TargetReferenceSpec},
\codepath{TrajectoryDiagnosticsSpec},
\codepath{WorkflowInputSpec},
\codepath{UnifiedSeedRecord}
&
Tipos, enumeraciones y envolturas que fijan entradas, solver, memoria,
continuación, muestreo, evidencia y salidas. Formalizan el protocolo, pero no
añaden una ecuación dinámica distinta. \\
\addlinespace

Ejecución de flujos &
\codepath{continue_integer_lure_seed},
\codepath{final_integer_lure_attractor},
\codepath{integer_lure_seed},
\codepath{integrate_integer_lure},
\codepath{run_integer_lure_hiddenness_controls},
\codepath{validate_full_workflow_system},
\codepath{load_config},
\codepath{save_effective_config},
\codepath{run_attractor_only_workflow},
\codepath{run_bifurcation_workflow},
\codepath{run_basin_workflow},
\codepath{run_simple_workflow}
&
Orquestación reproducible de los métodos anteriores. Las cuatro últimas rutas
permiten simular, analizar o construir bifurcaciones y cuencas sin iniciar una
búsqueda de atractor oculto. \\
\bottomrule
\end{longtable}
\end{center}

El inventario anterior se limita de forma deliberada a la API soportada. Los
auxiliares privados, scripts de validación y backends internos permanecen
auditables en el código y en \codepath{docs/api_reference.md}, pero no se
promueven como llamadas de usuario. La documentación de un símbolo experimental
module-qualified tampoco implica promoción a la raíz ni a la categoría
\STABLE; hace visible su contrato para evitar llamadas ambiguas.


\section{Mapa función-fundamento-validación}

\begin{center}
\scriptsize
\begin{longtable}{@{}L{0.17\textwidth}L{0.17\textwidth}L{0.22\textwidth}L{0.20\textwidth}L{0.14\textwidth}@{}}
\toprule
Función o módulo & Archivo fuente & Fundamento & Validación o evidencia & Lectura \\
\midrule
\endfirsthead
\toprule
Función o módulo & Archivo fuente & Fundamento & Validación o evidencia & Lectura \\
\midrule
\endhead
\codepath{ChuaParameters}, \codepath{chua_parameters} & \codepath{hidden_attractors/models/chua.py} & Parametrización del circuito de Chua. & Contrato local y pruebas de configuración. & Parámetros, no dinámica por sí solos. \\
\codepath{rhs_nonsmooth}, \codepath{rhs_arctan} & \codepath{models/chua.py} & Campo vectorial \(\CaputoD x=F(x)\). & Validación algebraica y pruebas unitarias de modelos. & Ecuación implementada. \\
\codepath{equilibria_nonsmooth}, \codepath{equilibria_arctan} & \codepath{models/chua.py} & Resolución de \(F(E)=0\). & Residuales algebraicos y comparación cruzada. & Equilibrios locales. \\
\codepath{jacobian_nonsmooth}, \codepath{jacobian_arctan} & \codepath{models/chua.py} & Linealización \(J=\partial F/\partial x\). & Diferencias finitas, Wolfram/Matlab cuando aplica. & Estabilidad local. \\
\codepath{ChaoticSystem}, \codepath{register_system} & \codepath{systems/} & Contrato abstracto de sistema autónomo. & Ejemplo \codepath{custom_system_definition.py} y capability checks. & Punto de extensión. \\
\codepath{LureSystem} & \codepath{systems/lure.py} & Split \(PX+b\psi(r^\top X)\). & Casos Chua y comparación Lur'e. & Necesario para DF/Nyquist. \\
\codepath{find_lure_harmonic_seed} & \codepath{seed_generation/} & \(1-N(A)\widehat W_q(\lambda)=0\). & Baseline entero y residuos complejos. & Semilla heurística. \\
\codepath{caputo_abm_integrate} & \codepath{integrations/abm.py} & ABM predictor-corrector de Caputo. & Reproducción Danca/contratos de historia completa. & Integración numérica. \\
\codepath{efork3_caputo_integrate} & \codepath{solvers/} & Runge--Kutta fraccionario EFORK y tablas manufacturadas Ghoreishi--Ghaffari. & Pruebas de referencia del solver. & Integración numérica experimental. \\
\codepath{integrate} & \codepath{integrations/selector.py} & Selector según \(q\) y backend. & Pruebas de compatibilidad. & Despacha método. \\
\codepath{FractionalProblem}, \codepath{solve_fractional_system} & \codepath{fractional/problem.py} & Contrato que separa derivada, método, orden, terminal, malla, condición inicial y memoria. & \codepath{test_fractional_problem.py} y pruebas específicas por solver. & Despacho reproducible; no eleva el estado de evidencia. \\
\codepath{grunwald_letnikov_derivative}, \codepath{native_grunwald_letnikov_derivative}, \codepath{fast_grunwald_letnikov_derivative} & \codepath{fractional/grunwald_letnikov.py}, \codepath{native_grunwald_letnikov.py}, \codepath{fast_grunwald_letnikov.py} & Convolución binomial directa, C/OpenMP o FFT. & Identidades GL/RL con Wolfram, convergencia finita y paridad de backends. & Operador muestreado; no solver por sí solo. \\
\codepath{lubich_convolution_quadrature} & \codepath{fractional/convolution_quadrature.py} & Pesos de \(\delta(\zeta)^q\) para BDF1/BDF2. & Alta precisión, monomios y paridad Python--Numba--FFT. & Operador experimental, sin correcciones de arranque. \\
\codepath{tempered_convolution_quadrature} & \codepath{fractional/tempered_convolution_quadrature.py} & CQ BDF1/BDF2 de RL templada y Caputo templada conjugada, con pesos \(e^{-\lambda kh}\omega_k\). & Solución manufacturada, 18 aserciones Wolfram independientes, paridad Python--Numba--FFT y benchmark portable. & Operador experimental por lotes; no solver ni historia rápida. \\
\codepath{tempered_fast_multistep_history} & \codepath{fractional/tempered_fast_history.py} & Fast Method II real con FBDF1/GNGF2, ventana local exacta e historia recurrente templada. & Pruebas directa--rápida y Python--Numba; 13/13 aserciones Wolfram independientes y comparador obligatorio con la API pública. & Operador muestral experimental; no solver FDE, y GNGF2 no es BDF2 fraccionaria. \\
\codepath{hadamard_convolution_quadrature} & \codepath{fractional/hadamard.py} & CQ en \(u=\log(t/a)\) para Hadamard y Caputo--Hadamard. & Fórmulas analíticas, tres backends y comparación Wolfram independiente de malla finita. & Operador sobre malla exponencial; no FDE. \\
\codepath{integrate_caputo_hadamard_abm} & \codepath{fractional/caputo_hadamard_solver.py} & Transformación exacta a un IVP Caputo en tiempo logarítmico y ABM/PECE. & Soluciones manufacturadas, paridad C--Python y caso ABM finito independiente. & Trayectoria finita experimental; no valida dinámica Chua. \\
\codepath{riemann_liouville_gl_derivative}, operadores templado, variable, conformable, CF y distribuido & \codepath{fractional/sampled_operators.py}, \codepath{caputo_fabrizio.py}, \codepath{distributed_order.py} & Kernels y discretizaciones no intercambiables. & Tests unitarios de fórmulas, límites y oráculos directos. & Operadores de datos con estado explícito. \\
\codepath{integrate_variable_order_caputo_type3_l1} & \codepath{fractional/variable_order_caputo_type3.py} & Solver Caputo tipo III, L1 de historia completa y corrector Picard. & Potencia manufacturada, reducción constante, paridad Numba--Python y oráculo Wolfram independiente. & Trayectoria finita experimental; no valida caos, atracción u ocultedad. \\
\codepath{compute_trajectory_metrics} & \codepath{analysis/trajectory.py} & Tiempo y estados explícitos, rangos, varianzas, acotamiento y resumen geométrico en dimensión arbitraria. & Tests de trayectorias \(1\)D, \(2\)D y estados puros \(4\)D. & Diagnóstico independiente. \\
\codepath{fft_spectrum}, \codepath{psd_welch} & \codepath{analysis/} & FFT/PSD post-transitorio. & Diagnósticos F5. & No certifica caos. \\
\codepath{integer_system_lyapunov_exponents} & \codepath{analysis/lyapunov.py} & QR/Benettin para \(q=1\). & Baseline entero. & Diagnóstico condicionado. \\
\codepath{compute_lyapunov_spectrum} & \codepath{analysis/lyapunov_api.py} & Despacho validado de métodos variacionales y de dinámica clonada según \(q\), memoria y Jacobiano. & Pruebas de contrato y entradas no finitas. & Diagnóstico condicionado. \\
\codepath{estimate_time_series_lyapunov} & \codepath{analysis/time_series_lyapunov.py} & Reconstrucción escalar Rosenstein/Eckmann y dimensión Kaplan--Yorke. & Series logística y sinusoidal, validación CLI y límites de memoria. & Diagnóstico de datos finitos. \\
\codepath{basin_entropy}, \codepath{uncertainty_fraction}, \codepath{estimate_uncertainty_exponent} & \codepath{analysis/basin_uncertainty.py} & Entropía por cajas, desacuerdo de destino y ajuste \(f(\varepsilon)\sim\varepsilon^\alpha\). & Cajas exactas, Wilson, políticas parciales y ajustes degenerados. & Diagnóstico de clasificación finita; no prueba Wada u ocultedad. \\
\codepath{generalized_delay_embedding}, estimadores de retardo y \codepath{false_nearest_neighbors} & \codepath{analysis/delay_embedding.py} & Vectores retardados, ACF, información mutua plug-in y FNN. & Alineación, tiempos irregulares, criterios y Theiler. & Reconstrucción empírica; no reconstruye automáticamente toda la historia FDE. \\
\codepath{auto_recurrence_matrix}, \codepath{cross_recurrence_matrix}, \codepath{joint_recurrence_matrix}, RQA avanzada & \codepath{analysis/recurrence_advanced.py} & Umbrales de distancia, Theiler, líneas diagonales/verticales y TREND. & Matrices exactas, paridad Numba--NumPy y guardas de memoria. & Diagnóstico denso de muestra finita. \\
\codepath{correlation_sum_curve}, \codepath{fit_correlation_dimension} & \codepath{analysis/correlation_dimension.py} y kernel C/OpenMP & Suma $q=2$, Theiler, conteo estricto y ajuste sobre rango explícito. & Paridad Python--Numba--C y oráculo Wolfram finito. & Diagnóstico de proyección; no certifica escala, fractalidad, caos u ocultedad. \\
\codepath{is_target_class} & \codepath{basins/classification.py} & Clase \TARGET{} en cuencas finitas. & Pruebas de hiddenness y basin runner. & Contactos bloquean ocultedad fuerte. \\
\codepath{WorkflowInputSpec} & \codepath{workflows/specs.py} & Contrato de solver, clasificador, referencia, cuencas y robustez. & Ejemplo \codepath{new_system_workflow_spec.py}. & Especificación reusable. \\
\codepath{StageEnvelope} & \codepath{workflows/protocol.py} & Envoltura uniforme de evidencia. & \codepath{validation/00_manifest/} y contrato oficial. & Trazabilidad. \\
\codepath{plot_*} & \codepath{plotting/} & Proyección visual de trayectorias y diagnósticos. & Galerías y smoke tests. & Figura auxiliar. \\
\bottomrule
\end{longtable}
\end{center}

\section{Ejemplos oficiales de la librería}

El ejemplo integral distribuible vive en
\codepath{version_2/examples/chua_integer_lure_reference/}. Los demás casos
científicos del repositorio se conservan como validaciones trazables y no como
promesas de la distribución PyPI.

\begin{center}
\footnotesize
\begin{longtable}{@{}L{0.26\textwidth}L{0.25\textwidth}L{0.39\textwidth}@{}}
\toprule
Ejemplo & Archivos activos & Papel en el reporte \\
\midrule
\codepath{examples/chua_integer_lure_reference/} & \codepath{run_example.py}, \codepath{reproducibility.yaml} & Referencia entera \(q=1\). Ejecuta localización directa, continuación y controles finitos de vecindad; el resultado registrado es compatible con un atractor oculto bajo el muestreo ensayado, no una prueba global de ocultedad. \\
\codepath{examples/dynamical_analysis_gallery.py} & API de análisis sobre una trayectoria suministrada por la persona usuaria. & FFT/PSD, bifurcación y métricas sin requerir una búsqueda de atractor oculto. \\
Núcleo fraccionario diverso & \codepath{examples/fractional_core_catalog.py} & Ejecuta operadores GL/RL, templado, variable, conformable, Caputo--Fabrizio, distribuido, CQ de Lubich/Hadamard y solvers manufacturados GL, conformable, ABC, Caputo templado y Caputo variable tipo III. \\
CQ templada manufacturada & \codepath{examples/tempered_convolution_quadrature.py} & Evalúa un sistema no lineal de dos componentes con solución analítica mediante CQ Caputo templada BDF2--FFT. Reporta residuo finito; no es un solver FDE ni evidencia dinámica. \\
Historia rápida templada sobre Chua entero & \codepath{examples/tempered_fast_history_chua.py} & Postprocesa una trayectoria Chua \(q=1\) uniformemente muestreada con GNGF2 Caputo templada y la contrasta con convolución directa. No es una solución Chua fraccionaria ni evidencia de caos, atracción u ocultedad. \\
Chua Caputo--Hadamard & \codepath{examples/caputo_hadamard_chua.py} & Integra una trayectoria Chua finita con ABM log y la compara con un control entero en la coordenada \(u\). No concluye caos, atracción, cuenca u ocultedad. \\
Análisis avanzado & \codepath{examples/chua_advanced_analysis.py} & Usa una trayectoria Chua entera real para delay/RQA y el flujo biestable exacto \(x'=x-x^3,\ y'=-y\) para entropía e incertidumbre de cuenca. \\
Dimensión de correlación & \codepath{examples/correlation_dimension_integer_fractional.py} & Aplica los mismos radios, Theiler y rango explícito a una trayectoria RK4 \(q=1\) y una Caputo ABM--PECE \(q=0.85\); no convierte las pendientes finitas en prueba dinámica. \\
\bottomrule
\end{longtable}
\end{center}

Los ejemplos pequeños de API siguen siendo útiles para aprendizaje, pero no son los casos científicos promovidos por este reporte. Los ejemplos nuevos producen registros JSON y están cubiertos por pruebas de ejecución dedicadas; conservar el contrato de software no transforma sus salidas en evidencia dinámica promovida.
\section{Caso científico 1: Chua entero de referencia}

El caso entero \(q=1\) es la línea base. Permite comprobar que el split de Lur'e, la transferencia, la función descriptiva, la continuación y los controles de cuenca funcionan sin la complejidad de memoria fraccionaria.

\subsection{Sistema y parámetros}

El sistema es~\eqref{eq:chua_nonsmooth} con \(q=1\). Parámetros:

\begin{center}
\small
\begin{tabular}{@{}L{0.24\textwidth}L{0.18\textwidth}L{0.18\textwidth}L{0.31\textwidth}@{}}
\toprule
Parámetro & Símbolo & Valor & Uso \\
\midrule
Razón de capacitores & \(\alpha\) & \(8.4562\) & Primera ecuación. \\
Inductancia normalizada & \(\beta\) & \(12.0732\) & Tercera ecuación. \\
Amortiguamiento resistivo & \(\gamma\) & \(0.0052\) & Disipación en \(z\). \\
Pendiente interna & \(m_0\) & \(-0.1768\) & Región \(|x|\le 1\). \\
Pendiente externa & \(m_1\) & \(-1.1468\) & Región \(|x|>1\). \\
Orden & \(q\) & \(1.0\) & Derivada ordinaria. \\
YAML & N/A & \codepath{examples/chua_integer_lure_reference/reproducibility.yaml} & Fuente del contrato del ejemplo. \\
\bottomrule
\end{tabular}
\end{center}

\subsection{Búsqueda y descarte de candidatos}

La búsqueda se hace con la forma Lur'e centrada y la función de transferencia racional del sistema entero. Las condiciones de fase y amplitud se resuelven directamente para obtener, en cada ejecución, \(\omega_0\), \(k\), \(a_0\) y la semilla inicial a partir de
\[
1-N(A)W_1(\ii\omega)=0.
\]
Los límites \(\omega_{\min}=10^{-5}\) y \(\omega_{\max}=50\) delimitan solamente el intervalo de admisibilidad de la solución; no definen una malla ni un barrido numérico de frecuencias.
La semilla resultante se transporta mediante continuación con \(\lambda=0,0.1,\ldots,1\). La simulación final usa \(h=0.01\), transitorio \(120\) y ventana conservada \(180\). El descarte se hace por controles de vecindades: si alguna bola alrededor de un equilibrio llega a \TARGET{}, el candidato no se promueve como oculto.

\begin{figure}[H]
\centering
\reportinclude[width=0.78\textwidth]{fig01_nyquist_df.pdf}
\caption{Chua entero: lectura Nyquist/DF del cruce armónico. La curva y el punto de balance validan la semilla, no la ocultedad.}
\end{figure}

\begin{figure}[H]
\centering
\reportinclude[width=0.78\textwidth]{fig01c_transfer_real_imag.pdf}
\caption{Chua entero: componentes real e imaginaria de la transferencia en el cruce armónico.}
\end{figure}

\begin{figure}[H]
\centering
\reportinclude[width=0.82\textwidth]{fig02d_continuation_story.pdf}
\caption{Chua entero: historia de continuación desde la semilla armónica hasta el sistema no lineal completo.}
\end{figure}

\begin{figure}[H]
\centering
\reportinclude[width=0.70\textwidth]{fig03_final_attractor.pdf}
\caption{Chua entero: trayectoria final post-continuación bajo el contrato del ejemplo.}
\end{figure}

\begin{figure}[H]
\centering
\reportinclude[width=0.78\textwidth]{fig05b_hiddenness_overview.pdf}
\caption{Chua entero: controles de vecindad alrededor de equilibrios. La lectura se limita a radios y muestras declaradas.}
\end{figure}

\begin{figure}[H]
\centering
\begin{minipage}{0.49\textwidth}
\centering
\reportinclude[width=\linewidth]{fig11a_fft_x.pdf}
\end{minipage}\hfill
\begin{minipage}{0.49\textwidth}
\centering
\reportinclude[width=\linewidth]{fig11d_psd_x.pdf}
\end{minipage}
\caption{Chua entero: FFT y PSD de \(x(t)\). Son diagnósticos espectrales, no criterios únicos de caos.}
\end{figure}

\begin{figure}[H]
\centering
\reportinclude[width=0.78\textwidth]{fig13_lyapunov_convergence.pdf}
\caption{Chua entero: convergencia de estimación Lyapunov bajo el contrato del ejemplo.}
\end{figure}
\section{Caso científico 2: Chua fraccionario no suave}

El caso fraccionario no suave conserva la saturación bilineal y usa derivada
de Caputo con \(q\) cercano a uno. En el repositorio coexisten tres registros
que no deben fusionarse:
\begin{enumerate}
  \item el control bibliográfico Danca2017, una implementación parcial que
  verifica ecuaciones, parámetros, equilibrios, Matignon y el contrato
  ABM--Caputo, pero no dispone de la condición inicial exacta publicada
  necesaria para reconstruir independientemente el atractor;
  \item el candidato oficial
  \codepath{m1_m1p2000_m0_m0p2000_branch_0}, evaluado en
  \codepath{validation/09_hiddenness_tests/}, que registró \(1305\) contactos
  desde vecindades de equilibrio y quedó rechazado como oculto bajo ese
  contrato; y
  \item la cadena no suave corregida de Paper07, con raíz BDF sesgada,
  continuación Caputo de historia única y evidencia compacta en
  \codepath{validation/paper07_chua/evidence/nonsmooth_corrected/}. Comparte
  los parámetros físicos nominales de Danca, pero su semilla y sus resultados
  pertenecen al método propuesto.
\end{enumerate}
El ejemplo que originó la raíz sesgada se conserva en
\codepath{examples/chua_nonsmooth_biased_hidden_attractor/} y su contrato en
\codepath{validation/paper07_chua/configs/chua_nonsmooth_biased_df_search.yaml}.

\subsection{Sistema}

\[
\begin{aligned}
{}^{C}D_t^q x &= \alpha(y-x-f(x)),\\
{}^{C}D_t^q y &= x-y+z,\\
{}^{C}D_t^q z &= -\beta y-\gamma z,
\end{aligned}
\qquad 0<q<1,
\]
con \(f\) dado por~\eqref{eq:nonsmooth_f}. El valor histórico usado en varios artefactos es \(q=0.9998\), cercano al caso entero pero con memoria explícita.

\subsection{Búsqueda, separación de registros y descarte}

La búsqueda del ejemplo se divide en etapas:
\begin{enumerate}
  \item búsqueda centrada de referencia para obtener ramas comparables;
  \item búsqueda por función descriptiva sesgada con \(A\), bias \(c\) y frecuencia \(\omega\);
  \item barrido \(m_1\in\{-1.1468,-1.20,-1.25\}\), \(m_0\in\{-0.1768,-0.20,-0.24\}\);
  \item continuación afín hacia el sistema original;
  \item simulación Caputo ABM con memoria completa;
  \item verificación por esferas y luego por un protocolo extendido de bolas.
\end{enumerate}

El criterio de descarte es explícito: un contacto \TARGET{} desde vecindades
locales de algún equilibrio bloquea una etiqueta fuerte de ocultedad para ese
registro. Esto rechaza al candidato oficial anterior, pero no debe trasladarse
por identidad de modelo a la cadena corregida: sus semillas, continuaciones y
nubes objetivo son diferentes. A la inversa, los sondeos favorables de
Paper07 tampoco rehabilitan el candidato rechazado ni completan los datos que
faltan en la referencia Danca.

\subsection{Reproducción bibliográfica y controles}

\begin{figure}[H]
\centering
\reportinclude[width=0.76\textwidth]{danca_frac_phase_3d.pdf}
\caption{Chua fraccionario no suave: trayectoria 3D de la reproducción bibliográfica Danca/ABM, usada como comparación dinámica y no como validación de ocultedad.}
\end{figure}

\begin{figure}[H]
\centering
\reportinclude[width=0.86\textwidth]{danca_frac_projections.pdf}
\caption{Chua fraccionario no suave: proyecciones de la reproducción bibliográfica.}
\end{figure}

\begin{figure}[H]
\centering
\reportinclude[width=0.86\textwidth]{danca_frac_time_series.pdf}
\caption{Chua fraccionario no suave: series temporales de la reproducción bibliográfica.}
\end{figure}

\begin{figure}[H]
\centering
\reportinclude[width=0.76\textwidth]{danca_frac_spectrum_x.pdf}
\caption{Chua fraccionario no suave: espectro de \(x(t)\) en la reproducción bibliográfica.}
\end{figure}

\subsection{Ruta centrada fraccionaria}

\begin{figure}[H]
\centering
\begin{minipage}{0.49\textwidth}
\centering
\reportinclude[width=\linewidth]{chua_frac_ns_fig01a_transfer_real.pdf}
\end{minipage}\hfill
\begin{minipage}{0.49\textwidth}
\centering
\reportinclude[width=\linewidth]{chua_frac_ns_fig01b_transfer_imag.pdf}
\end{minipage}
\caption{Chua fraccionario no suave: transferencia fraccionaria en la rama centrada.}
\end{figure}

\begin{figure}[H]
\centering
\reportinclude[width=0.82\textwidth]{chua_frac_ns_fig02_continuation_story.pdf}
\caption{Chua fraccionario no suave: continuación de la rama centrada.}
\end{figure}

\begin{figure}[H]
\centering
\reportinclude[width=0.68\textwidth]{chua_frac_ns_fig03_final_attractor.pdf}
\caption{Chua fraccionario no suave: trayectoria final de la rama centrada bajo contrato Caputo.}
\end{figure}

\begin{figure}[H]
\centering
\reportinclude[width=0.86\textwidth]{chua_frac_ns_fig03abc_final_attractor_projections.pdf}
\caption{Chua fraccionario no suave: proyecciones de la trayectoria final centrada.}
\end{figure}

\begin{figure}[H]
\centering
\reportinclude[width=0.86\textwidth]{chua_frac_ns_fig04_time_series.pdf}
\caption{Chua fraccionario no suave: series temporales de la rama centrada.}
\end{figure}

\begin{figure}[H]
\centering
\begin{minipage}{0.49\textwidth}
\centering
\reportinclude[width=\linewidth]{chua_frac_ns_fig11a_fft_x.pdf}
\end{minipage}\hfill
\begin{minipage}{0.49\textwidth}
\centering
\reportinclude[width=\linewidth]{chua_frac_ns_fig12_lyapunov_two_methods.pdf}
\end{minipage}
\caption{Chua fraccionario no suave: FFT y Lyapunov por dos métodos. Ambos son diagnósticos dependientes del contrato.}
\end{figure}

\begin{figure}[H]
\centering
\begin{minipage}{0.32\textwidth}
\centering
\reportinclude[width=\linewidth]{chua_frac_ns_fig13_E0_hiddenness_spherical_3d.pdf}
\end{minipage}\hfill
\begin{minipage}{0.32\textwidth}
\centering
\reportinclude[width=\linewidth]{chua_frac_ns_fig13_Ep_hiddenness_spherical_3d.pdf}
\end{minipage}\hfill
\begin{minipage}{0.32\textwidth}
\centering
\reportinclude[width=\linewidth]{chua_frac_ns_fig13_Em_hiddenness_spherical_3d.pdf}
\end{minipage}
\caption{Chua fraccionario no suave: sondeos esféricos alrededor de \(E_0\), \(E_+\) y \(E_-\) para la rama centrada.}
\end{figure}

\begin{figure}[H]
\centering
\reportinclude[width=0.78\textwidth]{chua_frac_ns_fig14_hiddenness_contact_heatmap.pdf}
\caption{Chua fraccionario no suave: mapa de contactos de vecindad para la rama centrada.}
\end{figure}

\subsection{Ruta sesgada propuesta}

La búsqueda de \(2430\) inicios seleccionó, dentro del par nominal y mediante
la regla \(c>0.05\) con residuo mínimo, la raíz BDF
\[
A=4.5778827210,\qquad
c=2.7763970033,\qquad
\omega=2.0402860511,
\]
con norma de residuo \(4.5490078475\times10^{-16}\). La reconstrucción
armónica sesgada de~\eqref{eq:biased_df}--\eqref{eq:biased_dc_harmonic_closure}
produjo la semilla
\[
X_{\rm sem}=
(7.3542797243,\;0.2909687788,\;-9.3160548187)^\top .
\]

La continuación transportó esa semilla por
\(\eta=0,0.05,\ldots,1\), es decir, \(21\) niveles dentro de una sola historia
causal de Caputo. El extremo registrado fue
\[
X_{\rm end}=
(0.8223107098,\;1.4652727040,\;1.2923082502)^\top .
\]
Desde \(X_{\rm end}\) comenzó un nuevo problema autónomo ABM--PECE de memoria
completa, con \(h=0.01\), horizonte \(300\) y descarte \(100\). Las
\(36/36\) perturbaciones de robustez, doce en cada uno de los radios
\(10^{-6}\), \(10^{-4}\) y \(10^{-2}\), hicieron contacto con la nube objetivo
según~\eqref{eq:d90_contact}. Esto respalda atracción local del conjunto, no
caos. La clasificación almacenada es
\codepath{regular_periodic_rejected}/\codepath{periodic_post_transient}; por
tanto, Paper07 documenta aquí un conjunto atrayente regular de apariencia
periódica y su cuenca finitamente muestreada, sin afirmación de atractor
caótico.

El sondeo volumétrico limpio completó \(17400\) problemas sin fallos numéricos
y aplicó la regla de completar el primer radio con contactos. La separación
entre vecindades locales y geometría macroscópica es:
\begin{center}
\small
\begin{tabular}{@{}L{0.28\textwidth}L{0.20\textwidth}r r@{}}
\toprule
Registro & Radios incluidos & Sondas & Contactos \\
\midrule
Contrato local & \(r\le0.01\) & \(7200\) & \(0\) \\
Auditoría acumulada sin contacto & \(r\le0.1\) & \(13800\) & \(0\) \\
Primer radio macroscópico con contacto & \(r=0.3\) & \(3600\) & \(37\) \\
\bottomrule
\end{tabular}
\end{center}
Los \(37\) contactos de \(r=0.3\) delimitan la geometría extendida y no se
ocultan en el total local. La afirmación compatible con la evidencia queda
restringida a las bolas locales \(r\le0.01\); no es una prueba global de
cuenca.

\begin{figure}[H]
\centering
\begin{minipage}{0.49\textwidth}
\centering
\reportinclude[width=\linewidth]{chua_frac_ns_biased_fig01a_transfer_real.pdf}
\end{minipage}\hfill
\begin{minipage}{0.49\textwidth}
\centering
\reportinclude[width=\linewidth]{chua_frac_ns_biased_fig01b_transfer_imag.pdf}
\end{minipage}
\caption{Chua fraccionario no suave sesgado: transferencia fraccionaria usada para generar semillas BDF.}
\end{figure}

\begin{figure}[H]
\centering
\reportinclude[width=0.82\textwidth]{chua_frac_ns_biased_fig02_continuation_story.pdf}
\caption{Chua fraccionario no suave sesgado: historia de continuación afín. La leyenda ``atractor caótico'' dentro de la figura es una etiqueta visual heredada; la clasificación canónica posterior es regular/periódica y no respalda caos.}
\end{figure}

\begin{figure}[H]
\centering
\reportinclude[width=0.68\textwidth]{chua_frac_ns_biased_fig03_attractor.pdf}
\caption{Chua fraccionario no suave sesgado: conjunto atrayente regular post-continuación; la clasificación registrada es periódica y no respalda caos.}
\end{figure}

\begin{figure}[H]
\centering
\reportinclude[width=0.86\textwidth]{chua_frac_ns_biased_fig04_projections.pdf}
\caption{Chua fraccionario no suave sesgado: proyecciones de fase del candidato.}
\end{figure}

\begin{figure}[H]
\centering
\reportinclude[width=0.86\textwidth]{chua_frac_ns_biased_fig04_timeseries.pdf}
\caption{Chua fraccionario no suave sesgado: series temporales del candidato.}
\end{figure}

\begin{figure}[H]
\centering
\begin{minipage}{0.49\textwidth}
\centering
\reportinclude[width=\linewidth]{chua_frac_ns_biased_fig11_fft.pdf}
\end{minipage}\hfill
\begin{minipage}{0.49\textwidth}
\centering
\reportinclude[width=\linewidth]{chua_frac_ns_biased_fig12_lyapunov.pdf}
\end{minipage}
\caption{Chua fraccionario no suave sesgado: FFT y Lyapunov como diagnósticos finito-temporales. La decisión dinámica canónica sigue siendo regular/periódica, sin promoción de caos.}
\end{figure}

\begin{figure}[H]
\centering
\begin{minipage}{0.32\textwidth}
\centering
\reportinclude[width=\linewidth]{chua_frac_ns_biased_fig13_E0_hiddenness_spherical_3d.pdf}
\end{minipage}\hfill
\begin{minipage}{0.32\textwidth}
\centering
\reportinclude[width=\linewidth]{chua_frac_ns_biased_fig13_Ep_hiddenness_spherical_3d.pdf}
\end{minipage}\hfill
\begin{minipage}{0.32\textwidth}
\centering
\reportinclude[width=\linewidth]{chua_frac_ns_biased_fig13_Em_hiddenness_spherical_3d.pdf}
\end{minipage}
\caption{Chua fraccionario no suave sesgado: sondeos esféricos alrededor de todos los equilibrios.}
\end{figure}

\begin{figure}[H]
\centering
\begin{minipage}{0.49\textwidth}
\centering
\reportinclude[width=\linewidth]{chua_frac_ns_biased_fig14_hiddenness_contact_heatmap.pdf}
\end{minipage}\hfill
\begin{minipage}{0.49\textwidth}
\centering
\reportinclude[width=\linewidth]{chua_frac_ns_biased_fig15_extended_heatmap.pdf}
\end{minipage}
\caption{Chua fraccionario no suave sesgado: contrato local de bolas hasta \(r=0.01\) y auditoría extendida; el primer contacto aparece en \(r=0.3\), después de completar ese radio.}
\end{figure}
\section{Caso científico 3: Chua fraccionario arctan}

En el paquete de validación registrado, el candidato c590 se promueve solo
como evidencia local limitada por radio: radios \(r\le0.3\), \(8400\) pruebas
finitas y cero contactos. Esta etiqueta no es prueba global de cuenca.

El caso arctan reemplaza la saturación por una función suave. Su interés es doble: evita superficies no diferenciables y permite una función descriptiva analítica cerrada. El modelo está vinculado a Wu et al.~2023~\cite{Wu2023ArctanChua}, pero el reporte separa estrictamente la reproducción bibliográfica de la metodología fraccionaria propuesta.

\subsection{Modelo}

\[
\begin{aligned}
\CaputoD x &= \alpha(y-x-f(x)),\\
\CaputoD y &= x-y+z,\\
\CaputoD z &= -\beta y-\gamma z,
\end{aligned}
\qquad
f(x)=a_1x+a_2\arctan(\rho x).
\]
El Jacobiano es suave:
\[
J(x)=
\begin{bmatrix}
-\alpha\left(1+a_1+\dfrac{a_2\rho}{1+\rho^2x^2}\right)&\alpha&0\\
1&-1&1\\
0&-\beta&-\gamma
\end{bmatrix}.
\]

\subsection{Función descriptiva}

Para \(\psi(\sigma)=a_2\arctan(\rho\sigma)\), el primer armónico produce
\[
N_{\arctan}(A)=a_2\frac{2(\sqrt{1+\rho^2A^2}-1)}{\rho A^2}.
\]
Esta fórmula se implementa en la forma \codepath{LureSystem}. Su lectura es semilla heurística; no prueba ocultedad.

\subsection{Wu2023: reproducción bibliográfica y límite metodológico}

El caso Wu2023 usa \(\alpha=8.4562\), \(\beta=12.0732\), \(\gamma=0.0052\), \(a_1=0.4\), \(a_2=-1.5585\), \(\rho=1\), \(q=0.99\), \(h=0.01\). En el repositorio se conserva como contrato de validación en \codepath{validation/reference_cases/fractional_chua_arctan_wu2023/config.json} y como ejemplo reproducible en \codepath{examples/chua_arctan_wu2023/}.

La reproducción publicada no pudo promoverse como evidencia fraccionaria completa de ocultedad por una razón de contrato: el ADM usado en la ruta bibliográfica local no acumula memoria Caputo completa (\codepath{memory_policy: none_local_adm}) y la DF/NC publicada se evalúa con transferencia entera \({\left(\ii\omega I-P\right)}^{-1}b\), con \(q_{seed}=1\). Por eso las figuras de Wu2023 se reportan como reproducción bibliográfica y comparación dinámica; no sustituyen una búsqueda fraccionaria Lur'e-Caputo con historia completa ni pruebas de vecindad de todos los equilibrios.

\begin{figure}[H]
\centering
\begin{minipage}{0.49\textwidth}
\centering
\reportinclude[width=\linewidth]{wu2023_chua_fractional_arctan_q099_x0_plus_phase3d.pdf}
\end{minipage}\hfill
\begin{minipage}{0.49\textwidth}
\centering
\reportinclude[width=\linewidth]{wu2023_chua_fractional_arctan_q099_x0_minus_phase3d.pdf}
\end{minipage}
\caption{Wu2023 arctan: trayectorias 3D para las condiciones iniciales reportadas con signo opuesto.}
\end{figure}

\begin{figure}[H]
\centering
\begin{minipage}{0.49\textwidth}
\centering
\reportinclude[width=\linewidth]{wu2023_chua_fractional_arctan_q099_x0_plus_projections.pdf}
\end{minipage}\hfill
\begin{minipage}{0.49\textwidth}
\centering
\reportinclude[width=\linewidth]{wu2023_chua_fractional_arctan_q099_x0_minus_projections.pdf}
\end{minipage}
\caption{Wu2023 arctan: proyecciones de fase de la reproducción bibliográfica.}
\end{figure}

\begin{figure}[H]
\centering
\begin{minipage}{0.49\textwidth}
\centering
\reportinclude[width=\linewidth]{wu2023_chua_fractional_arctan_q099_x0_plus_timeseries.pdf}
\end{minipage}\hfill
\begin{minipage}{0.49\textwidth}
\centering
\reportinclude[width=\linewidth]{wu2023_chua_fractional_arctan_q099_x0_minus_timeseries.pdf}
\end{minipage}
\caption{Wu2023 Chua arctan: series temporales de la reproducción bibliográfica; este resultado no forma parte de la validación promovida de ocultedad, cuya evidencia c590 está limitada a radios locales.}
\end{figure}

\subsection{Candidato dinámico c590}

El candidato \codepath{chua_arctan_c590_q09999_seed9} usa
\[
\alpha=21.8493569066,\quad \beta=19.0818408409,\quad \gamma=0.00737801198,
\]
\[
a_1=0.04228979344,\quad a_2=-3.3367815123,\quad \rho=1.7984259333,\quad q=0.9999.
\]
La semilla de referencia es
\[
(5.8642449791,\;1.5847111486,\;3.2155806478).
\]

La búsqueda se documenta como una ruta dinámica: exploración global de \(2400\) casos en orden entero, refinamiento local de \(1000\) casos, selección del índice 590, barrido de problemas de valor inicial Caputo en \(q\), extracción de estados interiores y reinicios sesgados. Cada corrida fraccionaria acumula memoria completa desde su propio tiempo inicial; la historia no se transfiere entre órdenes ni entre reinicios. Entre los estados acotados en los tres pasos, la semilla 9 se selecciona por su apoyo robusto en la prueba 0--1 para \(h=0.0025\) y \(h=0.005\).

El resumen vigente usa radios
\(10^{-5},10^{-4},10^{-3},10^{-2},0.03,0.1,0.3,1,2\). Para los radios locales
hasta \(r=0.3\) registra \(8400\) pruebas y cero contactos alrededor de todos
los equilibrios: \(8396\) trayectorias completan el horizonte y cuatro llegan
anticipadamente a un equilibrio. El candidato c590 se promueve como
\codepath{hiddenness_supported_under_tested_local_radii}. El primer radio con
contacto es \(r=1\): las \(2400\) trayectorias completan el horizonte y \(22\)
hacen contacto con la nube objetivo. La auditoría en \(r=2\) se conserva
separadamente: \(588\) contactos en \(2700\) pruebas, con \(2569\)
trayectorias completas y \(131\) que alcanzan el umbral de divergencia
declarado. En conjunto, \(r=1\) y \(r=2\) aportan \(610\) contactos en
\(5100\) sondeos. El contacto en \(r=1\) fija el inicio observado de la
geometría macroscópica de cuenca; no contradice la ausencia de contactos en
las superficies locales ya completadas, pero impide extender esa lectura más
allá de \(r=0.3\).

\begin{figure}[H]
\centering
\reportinclude[width=0.74\textwidth]{chua_frac_arctan_c590_fig00_initial_seed.pdf}
\caption{Arctan c590: semilla de búsqueda, trayectoria objetivo y equilibrios.}
\end{figure}

\begin{figure}[H]
\centering
\reportinclude[width=0.68\textwidth]{chua_frac_arctan_c590_fig03_attractor.pdf}
\caption{Arctan c590: trayectoria objetivo bajo ABM Caputo de memoria completa.}
\end{figure}

\begin{figure}[H]
\centering
\reportinclude[width=0.86\textwidth]{chua_frac_arctan_c590_fig04_timeseries.pdf}
\caption{Arctan c590: series temporales de la trayectoria objetivo.}
\end{figure}

\begin{figure}[H]
\centering
\begin{minipage}{0.49\textwidth}
\centering
\reportinclude[width=\linewidth]{chua_frac_arctan_c590_fig11_fft.pdf}
\end{minipage}\hfill
\begin{minipage}{0.49\textwidth}
\centering
\reportinclude[width=\linewidth]{chua_frac_arctan_c590_fig12_lyapunov_two_methods.pdf}
\end{minipage}
\caption{Arctan c590: FFT y estimación Lyapunov por dos métodos. Su lectura es diagnóstica y depende de validación del método.}
\end{figure}

\begin{figure}[H]
\centering
\begin{minipage}{0.32\textwidth}
\centering
\reportinclude[width=\linewidth]{chua_frac_arctan_c590_fig13_E0_hiddenness_spherical_3d.pdf}
\end{minipage}\hfill
\begin{minipage}{0.32\textwidth}
\centering
\reportinclude[width=\linewidth]{chua_frac_arctan_c590_fig13_Ep_hiddenness_spherical_3d.pdf}
\end{minipage}\hfill
\begin{minipage}{0.32\textwidth}
\centering
\reportinclude[width=\linewidth]{chua_frac_arctan_c590_fig13_Em_hiddenness_spherical_3d.pdf}
\end{minipage}
\caption{Arctan c590: sondeos esféricos alrededor de \(E_0\), \(E_+\) y \(E_-\).}
\end{figure}

\begin{figure}[H]
\centering
\reportinclude[width=0.78\textwidth]{chua_frac_arctan_c590_fig14_hiddenness_contact_heatmap.pdf}
\caption{Arctan c590: mapa de contactos de vecindad. La afirmación promovida usa radios locales hasta \(r=0.3\); el primer contacto aparece en \(r=1\) y la auditoría de \(r=2\) se conserva como geometría macroscópica extendida.}
\end{figure}
\section{Validación oficial y evidencia promovida}

La fuente normativa de conteos y estados es
\codepath{validation/freeze_audit/}; este reporte no fija una cifra que pueda
quedar obsoleta frente al artefacto generado.

\subsection{Matriz resumida de afirmaciones defendibles}

La tabla a continuación resume las afirmaciones científicamente defendibles para los tres casos Chua del reporte y para la metodología general. La fuente normativa sigue siendo \codepath{validation/references/thesis_claims_matrix.md}; el reporte no debe promover resultados por encima de esa matriz ni de los resúmenes vigentes.

\begin{center}
\footnotesize
\begin{longtable}{@{}L{0.21\textwidth}L{0.18\textwidth}L{0.24\textwidth}L{0.27\textwidth}@{}}
\toprule
Caso & Estado defendible & Evidencia principal & Límite de lectura \\
\midrule
Chua entero de referencia & \codepath{reproducido} & YAML del ejemplo, figuras de Lur'e/continuacion, controles de vecindad. & Chua entero valida la ruta entera; no valida resultados fraccionarios. \\
Chua fraccionario no suave, reproducción Danca & Parcial; no promovida & Contrato bibliográfico: ecuaciones, parámetros, equilibrios, Matignon y ABM--Caputo. & La referencia no aporta la condición inicial exacta ni los datos suficientes para reconstruir independientemente su atractor. \\
Candidato oficial fraccionario no suave \codepath{branch_0} & \codepath{rechazado} & \codepath{validation/09_hiddenness_tests/}; \(1305\) contactos registrados. & El contacto desde vecindades locales de \(E_+\) y \(E_-\) bloquea ocultedad para este candidato. \\
Paper07 no suave corregido & Regular y periódico; compatible localmente & Raíz BDF, semilla, 21 niveles con historia única, \(36/36\) controles y \(0/7200\) contactos para \(r\le0.01\). & No hay afirmación de caos ni prueba global; \(37/3600\) contactos aparecen en \(r=0.3\). \\
Chua arctan fraccionario c590 & promovido localmente & \codepath{validation/chua_fractional_arctan_c590/}; \(8400\) pruebas sin contactos hasta \(r=0.3\). & El primer contacto es \(22/2400\) en \(r=1\); \(r=2\) conserva \(588/2700\) contactos y \(131\) divergencias como auditoría extendida. \\
Metodología Lur'e/Caputo & validada como framework & Modelos, API, semillas, integradores, verificación, manifiestos y reporte. & Alcance restringido a sistemas compatibles con Lur'e escalar. \\
\bottomrule
\end{longtable}
\end{center}

\subsection{Validación Wolfram y trazabilidad de Paper07}

La comprobación algebraica independiente usa los casos
\codepath{validation/wolfram/cases/chua_fractional_saturation.wl},
\codepath{validation/wolfram/cases/chua_fractional_arctan_c590.wl} y
\codepath{validation/wolfram/cases/hadamard_fractional_operator.wl}. El resumen
oficial no suave
\codepath{validation/02_algebraic_validation/algebraic_validation_validation_summary.json}
registra la comparación Python--Wolfram de equilibrios, Jacobianos,
autovalores/Matignon, equivalencia de Lur'e y transferencia. Para c590,
\codepath{validation/chua_fractional_arctan_c590/algebraic_validation/}
conserva el resumen aprobado del álgebra y la consistencia del candidato
registrado. Estas comprobaciones certifican fórmulas y parametrización; no
certifican por sí solas la selección dinámica de la semilla, caos, cuenca u
ocultedad.

El caso Hadamard recalcula la transformación logarítmica, identidades
Hadamard/Caputo--Hadamard, CQ BDF1/BDF2, el límite \(q=1\) y una integración
ABM manufacturada sin leer la implementación Python ni el texto del reporte.
Los artefactos retenidos están bajo
\codepath{validation/outputs/wolfram/hadamard_fractional_operator/}:
\begin{verbatim}
hadamard_fractional_operator_validation_summary.json
hadamard_python_comparison.json
\end{verbatim}
Ambos registran \codepath{passed=true}, una peor diferencia de operador de
\(3.019806626980426\times10^{-14}\) y una diferencia máxima de estado ABM de
\(8.881784197001252\times10^{-16}\). Esa evidencia verifica fórmulas y
resultados discretos finitos; no establece convergencia general, estabilidad,
caos, atracción u ocultedad.

La trazabilidad del artículo se organiza desde
\codepath{validation/paper07_chua/manifest.json} hacia
\codepath{validation/paper07_chua/evidence_manifest.json}, que fija tamaños y
huellas SHA-256, y de allí a los contratos, resúmenes y tablas bajo
\codepath{validation/paper07_chua/evidence/}. La verificación
\begin{verbatim}
python validation/paper07_chua/scripts/sync_paper07_evidence.py --verify
\end{verbatim}
comprueba esa proyección compacta sin depender de los directorios de salida
ignorados. Así se mantienen separados resultados propuestos, controles
bibliográficos y procedencia de exploración.

\subsection{Protocolo de validación}

El protocolo Caputo oficial está definido por estos dos archivos:
\begin{flushleft}
\codepath{configs/validation_contract.json}\\
\codepath{validation/00_manifest/source_configs/unified_caputo_protocol.json}
\end{flushleft}
Cada etapa debe escribir un resumen con estado, entradas, salidas, métricas,
archivos y metadatos. El orden de lectura es vinculante:

\begin{verbatim}
numerical_contract -> algebraic_validation -> seed_generation -> soft_precheck
-> continuation -> post_continuation_filter -> dynamic_reference -> robustness
-> hiddenness_tests -> diagnostics
\end{verbatim}

\begin{center}
\footnotesize
\begin{longtable}{@{}L{0.24\textwidth}L{0.43\textwidth}L{0.25\textwidth}@{}}
\toprule
Etapa & Qué prueba & Qué no prueba \\
\midrule
\codepath{numerical_contract} & Solver, backend, memoria, \(q\), \(h\), horizonte y tolerancias. & Correctitud dinámica por sí sola. \\
\codepath{algebraic_validation} & RHS, equilibrios, Jacobianos, Matignon, Lur'e. & Existencia de atractor. \\
\codepath{seed_generation} & Semillas uniformes y residuos. & Ocultedad. \\
\codepath{soft_precheck} & Admisibilidad diagnóstica. & Rechazo matemático fuerte. \\
\codepath{continuation} & Transporte numérico y fallos. & Continuación formal global. \\
\codepath{post_continuation_filter} & Duplicados, periodicidad operativa, no finitud. & Caos certificado. \\
\codepath{dynamic_reference} & Trayectoria larga, geometría, espectros. & Cuenca global. \\
\codepath{robustness} & Persistencia finita ante perturbaciones. & Invariancia matemática. \\
\codepath{hiddenness_tests} & Bolas, esferas y cortes alrededor de equilibrios. & No intersección global exacta. \\
\codepath{diagnostics} & FFT, PSD, Poincaré, Lyapunov y comparaciones. & Ocultedad verificada. \\
\bottomrule
\end{longtable}
\end{center}

La validación publicada de Lyapunov queda fuera de la suite rápida. La ruta
DK2018 larga se ejecuta con:

\begin{verbatim}
$env:RUN_PUBLISHED_LYAPUNOV=1
python -m pytest tests/test_efork_published_validation.py -q
\end{verbatim}

El estado registrado conserva discrepancias: Lorenz pasó y RF no satisfizo el
umbral registrado para \(\lambda_3\). Esto se documenta como
\codepath{published_benchmarks_pending_reproduced_discrepancy}, no como
fracaso silencioso ni como validación completa del método local.

\section{Cómo interpretar las figuras}

\begin{center}
\footnotesize
\resizebox{0.98\textwidth}{!}{%
\begin{tabular}{@{}L{0.22\textwidth}L{0.27\textwidth}L{0.25\textwidth}L{0.21\textwidth}@{}}
\toprule
Figura & Qué muestra & Lectura permitida & Lectura prohibida \\
\midrule
Espacio de fases & Nube post-transitorio en 3D. & Geometría y acotamiento aparente. & Prueba de caos u ocultedad. \\
Proyecciones & Cortes \(xy,xz,yz\). & Simetría, escala y estructura. & Atracción global. \\
Series temporales & Evolución por coordenada. & Transitorios, saturación, divergencia. & Existencia global. \\
FFT/PSD & Energía espectral. & Periodicidad candidata o espectro ancho. & Certificación de caos. \\
Poincaré & Cruces geométricos. & Dispersión de secciones. & Mapa exacto Caputo. \\
Lyapunov & Espectro estimado bajo método. & Sensibilidad numérica. & Validación automática. \\
Cuencas & Etiquetas por malla, bola o esfera. & Contactos \TARGET{} o ausencia finita de contactos. & No intersección global exacta. \\
Robustez & Superposición de corridas. & Persistencia bajo perturbaciones finitas. & Prueba matemática. \\
Nyquist/DF & Cruce armónico. & Semilla \(A,\omega\). & Existencia de atractor. \\
\bottomrule
\end{tabular}
}
\end{center}

\section{Reproducibilidad del reporte}

El reporte no debe depender de scripts sueltos de raíz ni de carpetas
históricas. El bundle conservado se encuentra en
\codepath{library_figures/by_report/unified_chua_fractional/}. El generador
mantenido resuelve recursivamente los \codepath{input/include} del LaTeX,
sincroniza el bundle y comprueba que cada \codepath{reportinclude} esté
declarado y presente:

\begin{verbatim}
cd version_2
python -m tools.unified_report_assets
\end{verbatim}

Las figuras del catálogo PNG siguen su pipeline separado. La prueba
\codepath{tests/test_unified_report_figure_generator.py} fija el descubrimiento
recursivo y la verificación del manifiesto.

Compilación del PDF

\begin{verbatim}
cd version_2/docs
latexmk -pdf -interaction=nonstopmode `
  -halt-on-error reporte_unificado_chua_fraccionario.tex
\end{verbatim}

Ejecución rápida de las rutas distribuibles de usuario:

\begin{verbatim}
cd version_2
python examples/chua_integer_lure_reference/run_example.py --quick
python examples/dynamical_analysis_gallery.py
python examples/fractional_core_catalog.py
python examples/caputo_hadamard_chua.py
python examples/chua_advanced_analysis.py
python examples/correlation_dimension_integer_fractional.py
python examples/tempered_convolution_quadrature.py
python examples/tempered_fast_history_chua.py
python -m pytest tests/test_fractional_core_examples.py -q
python -m pytest tests/test_correlation_dimension_example.py -q
python -m pytest tests/test_hadamard_wolfram.py -q
python -m pytest tests/test_tempered_convolution_quadrature.py \
  tests/test_tempered_convolution_quadrature_example.py \
  tests/test_tempered_convolution_quadrature_wolfram.py \
  tests/test_tempered_convolution_quadrature_benchmark.py -q
python -m pytest tests/test_tempered_fast_history.py \
  tests/test_tempered_fast_history_edges.py \
  tests/test_tempered_fast_history_example.py \
  tests/test_tempered_fast_history_wolfram.py \
  tests/test_tempered_fast_history_benchmark.py -q
\end{verbatim}

Las campañas científicas descritas en las secciones anteriores se leen desde
sus paquetes de evidencia bajo \codepath{validation/}; no se presentan como
ejemplos públicos ni como capacidades adicionales de la distribución.

\begin{center}
\footnotesize
\begin{longtable}{@{}L{0.25\textwidth}L{0.30\textwidth}L{0.35\textwidth}@{}}
\toprule
Ruta & Entrada & Alcance \\
\midrule
Chua entero & \codepath{examples/chua_integer_lure_reference/run_example.py} & Referencia reproducible del pipeline y candidato oculto bajo los controles finitos de vecindad ensayados. \\
Caracterización & \codepath{examples/dynamical_analysis_gallery.py} & Métricas y figuras sobre una trayectoria, sin búsqueda de ocultedad. \\
Catálogo fraccionario & \codepath{examples/fractional_core_catalog.py} & Operadores sobre funciones conocidas y un solver manufacturado; evidencia algebraico-numérica finita. \\
Caputo--Hadamard Chua & \codepath{examples/caputo_hadamard_chua.py} & Demostración de solver sobre el campo Chua y control entero en coordenada logarítmica; sin claim dinámico. \\
Análisis avanzado & \codepath{examples/chua_advanced_analysis.py} & Delay/RQA sobre Chua entero y métricas de cuenca sobre un flujo biestable exactamente clasificable. \\
Dimensión de correlación & \codepath{examples/correlation_dimension_integer_fractional.py} & Curva \(q=2\) y ajuste explícito comunes a dos trayectorias pequeñas; diagnóstico de proyección finita. \\
\bottomrule
\end{longtable}
\end{center}
\section{Limitaciones}

\begin{itemize}
  \item Toda conclusión numérica es finito-temporal.
  \item El resultado depende de la definición de derivada, \(q\), terminal,
  prehistoria o condición inicial, coordenada de malla, paso, horizonte,
  transitorio, memoria, backend y tolerancias.
  \item Un operador sobre muestras no es un solver FDE; una trayectoria finita
  tampoco prueba estabilidad, caos, atracción u ocultedad.
  \item Las figuras son evidencia auxiliar; no reemplazan contratos ni
  archivos de origen.
  \item DF/Nyquist son generadores heurísticos de semillas y no pruebas de existencia u ocultedad.
  \item Matignon decide estabilidad local de equilibrios linealizados, no
  ocultedad.
  \item Lyapunov, FFT, PSD, Poincaré, entropía de cuenca, embedding y RQA son
  diagnósticos si el método no está validado bajo el contrato correspondiente.
  En una FDE, proximidad de estados proyectados no implica proximidad de
  historias completas.
  \item Los comparadores y campañas específicas viven bajo
  \codepath{validation/}; no forman parte de la API pública instalada.
\end{itemize}

\appendix

\section{Comandos útiles}

\begin{center}
\footnotesize
\begin{longtable}{@{}L{0.34\textwidth}L{0.43\textwidth}L{0.15\textwidth}@{}}
\toprule
Comando & Uso & Estado \\
\midrule
\cmdbreak{hidden-attractors}{--help} & Muestra la ayuda general del CLI unificado. & Usuario \\
\cmdbreak{hidden-attractors}{inspect systems} & Lista o inspecciona sistemas registrados. & Usuario \\
\cmdbreak{hidden-attractors}{inspect candidates --source <ruta>} & Lee candidatos desde un JSON explícito sin depender del checkout. & Usuario \\
\cmdbreak{hidden-attractors}{inspect workflow-requirements} & Explica requisitos de workflows. & Usuario \\
\cmdbreak{hidden-attractors}{run -p chua\_integer} & Ejecuta el preset entero de referencia. & Usuario \\
\cmdbreak{hidden-attractors}{run -p chua\_fractional} & Ejecuta el preset fraccionario no suave. & Usuario \\
\cmdbreak{hidden-attractors}{run -c path/to/config.yaml} & Ejecuta un YAML definido por el usuario. & Usuario \\
\cmdbreak{hidden-attractors}{init -e chua\_fractional} & Copia un YAML editable. & Usuario \\
\cmdbreak{hidden-attractors}{seed lure-centered --help} & Ayuda para semillas Lur'e centradas por función descriptiva. & Experimental \\
\cmdbreak{hidden-attractors}{seed lure-biased --help} & Ayuda para semillas Lur'e sesgadas/BDF & Experimental \\
\cmdbreak{hidden-attractors}{continuation run --help} & Continuación de un candidato bajo contrato declarado. & Experimental \\
\cmdbreak{hidden-attractors}{continuation multiparameter --help} & Continuación multiparámetro. & Experimental \\
\cmdbreak{hidden-attractors}{protocol <stage>} & Ejecuta una etapa del protocolo oficial. & Alfa \\
\cmdbreak{hidden-attractors}{hiddenness sphere-controls --help} & Pruebas de vecindades/esferas alrededor de equilibrios. & Alfa \\
\cmdbreak{hidden-attractors}{robustness overlay --help} & Barridos de robustez bajo perturbaciones. & Alfa \\
\cmdbreak{hidden-attractors}{bifurcation run --help} & Barrido de bifurcación. & Alfa \\
\cmdbreak{hidden-attractors}{lyapunov compute --help} & Cálculo variacional de exponentes de Lyapunov. & Experimental \\
\cmdbreak{hidden-attractors}{lyapunov spectrum --help} & Lyapunov reconstruido desde una serie temporal escalar. & Integrado (API experimental) \\
\cmdbreak{hidden-attractors}{chaos-test --help} & Pruebas diagnósticas de caos. & Alfa \\
\cmdbreak{python examples/}{fractional\_core\_catalog.py} & Ejecuta el catálogo de operadores fraccionarios sobre funciones conocidas. & Experimental \\
\cmdbreak{python examples/}{caputo\_hadamard\_chua.py} & Ejecuta la demostración Chua con ABM Caputo--Hadamard. & Experimental \\
\cmdbreak{python examples/}{chua\_advanced\_analysis.py} & Ejecuta delay/RQA Chua y métricas de la cuenca biestable. & Experimental \\
\begin{tabular}[t]{@{}l@{}}\texttt{python examples/}\\\texttt{correlation\_dimension\_}\\\texttt{integer\_fractional.py}\end{tabular} & Compara el contrato finito \(q=2\) sobre trayectorias entera y fraccionaria. & Experimental \\
\begin{tabular}[t]{@{}l@{}}\texttt{python examples/}\\\texttt{tempered\_convolution\_}\\\texttt{quadrature.py}\end{tabular} & Ejecuta el caso manufacturado no lineal CQ Caputo templado BDF2--FFT. & Experimental \\
\begin{tabular}[t]{@{}l@{}}\texttt{python examples/}\\\texttt{tempered\_fast\_history\_}\\\texttt{chua.py}\end{tabular} & Ejecuta el postprocesamiento GNGF2 Caputo templado sobre muestras Chua enteras y su contraste directo. & Experimental \\
\begin{tabular}[t]{@{}l@{}}\texttt{python -m pytest}\\\texttt{tests/test\_fractional\_core\_}\\\texttt{examples.py -q}\end{tabular} & Prueba de humo de los tres ejemplos del núcleo fraccionario. & Desarrollo \\
\begin{tabular}[t]{@{}l@{}}\texttt{python -m pytest}\\\texttt{tests/test\_hadamard\_}\\\texttt{wolfram.py -q}\end{tabular} & Verifica el caso Wolfram independiente y su comparación finita con Python. & Desarrollo \\
\begin{tabular}[t]{@{}l@{}}\texttt{python -m}\\\texttt{tools.unified\_report\_assets}\end{tabular} & Sincroniza y verifica el bundle canónico del reporte. & Editorial \\
\cmdbreak{hidden-attractors}{validate contract} & Revisa estrictamente el contrato de validación registrado. & Usuario \\
\bottomrule
\end{longtable}
\end{center}

Los comandos independientes antiguos no forman parte de la superficie pública instalada en esta versión. Se reemplazan por subcomandos agrupados bajo \codepath{hidden-attractors}:

\noindent\textbf{Legacy/deprecated:}
\texttt{hidden-attractors-check-validation}. Reemplazo:
\texttt{hidden-attractors validate contract}.

\noindent\textbf{Legacy/deprecated:}
\texttt{hidden-attractors-protocol}. Reemplazo:
\texttt{hidden-attractors protocol}.

\subsection{Salidas esperadas de comandos mínimos}

En la siguiente tabla se resumen las salidas o efectos esperados al ejecutar los comandos de prueba mínimos de la librería.

\begin{center}
\footnotesize
\begin{longtable}{@{}L{0.34\textwidth}L{0.28\textwidth}L{0.30\textwidth}@{}}
\toprule
Comando & Archivos o salida esperada & Interpretación \\
\midrule
\codepath{hidden-attractors inspect systems} &
Salida terminal. &
Lista los sistemas registrados y permite verificar que la instalación cargó el registro. \\
\codepath{hidden-attractors run -p chua\_integer} &
Resumen JSON, tablas CSV y figuras si el preset activa exportación. &
Control entero \(q=1\); valida la ruta de uso sin memoria fraccionaria. \\
\codepath{python examples/quickstart\_equilibria.py} &
Resumen corto de ejemplo. &
Prueba de humo del ejemplo correspondiente; no constituye validación completa. \\
\codepath{python examples/fractional\_core\_catalog.py} &
JSON en la salida estándar. &
Compara contratos y endpoints finitos de operadores; no demuestra convergencia global. \\
\codepath{python examples/caputo\_hadamard\_chua.py} &
JSON con trayectoria fraccionaria y control entero. &
Demuestra la llamada y sus metadatos; no prueba caos, atracción u ocultedad. \\
\codepath{python examples/chua\_advanced\_analysis.py} &
JSON con delay, FNN, recurrencia, RQA y métricas de cuenca. &
Diagnósticos finitos de dos sistemas declarados; no son evidencia promovida. \\
\codepath{python examples/correlation\_dimension\_integer\_fractional.py} &
JSON con dos curvas, rango de ajuste y procedencia de trayectoria. &
Diagnóstico finito común; no prueba dimensión fractal, caos u ocultedad. \\
\codepath{python examples/tempered\_fast\_history\_chua.py} &
JSON con paridad rápida--directa, calibración de pesos y alcance explícito. &
Postprocesa muestras Chua enteras; no resuelve una FDE ni prueba caos, atracción u ocultedad. \\
\cmdbreak{hidden-attractors}{validate contract} &
Diagnóstico de contrato. &
Revisa estrictamente la estructura de evidencia registrada. \\
\bottomrule
\end{longtable}
\end{center}

La opción \codepath{--quick} debe usarse sólo si el ejemplo la implementa.

\section{Bibliografía}

\begin{thebibliography}{99}
\sloppy
\RaggedRight\relax

\bibitem{Chua1986}
R. N. Madan and L. O. Chua.
\emph{Chaos in Chua's Circuit}.
IEE Proceedings D:\@ Control Theory and Applications, 1986.

\bibitem{Caputo1967}
M. Caputo.
\emph{Linear Models of Dissipation whose Q is almost Frequency Independent-II}.
Geophysical Journal International, 1967.

\bibitem{Matignon1996}
D. Matignon.
\emph{Stability Results for Fractional Differential Equations with Applications to Control Processing}.
Proceedings of the Computational Engineering in Systems Applications
Multiconference, Lille, France, pp. 963--968, 1996.

\bibitem{Lorenz1963}
E. N. Lorenz.
\emph{Deterministic Nonperiodic Flow}.
Journal of the Atmospheric Sciences, vol. 20, no. 2, pp. 130--141, 1963.
DOI:\@ \href{https://doi.org/10.1175/1520-0469(1963)020<0130:DNF>2.0.CO;2}{doi}.

\bibitem{Danca2017}
M. F. Danca.
\emph{Hidden Chaotic Attractors in Fractional-Order Systems}.
Nonlinear Dynamics, vol. 89, no. 1, pp. 577--586, 2017.
DOI:\@ \href{https://doi.org/10.1007/s11071-017-3472-7}{doi}.

\bibitem{LeonovKuznetsov2013}
G. A. Leonov and N. V. Kuznetsov.
\emph{Hidden Attractors in Dynamical Systems: From Hidden Oscillations in Hilbert-Kolmogorov, Aizerman, and Kalman Problems to Hidden Chaotic Attractor in Chua Circuits}.
International Journal of Bifurcation and Chaos, vol. 23, no. 1, 1330002,
2013, DOI:\@ \href{https://doi.org/10.1142/S0218127413300024}{doi}.

\bibitem{KuznetsovEtAl2017}
N. V. Kuznetsov, O. A. Kuznetsova, G. A. Leonov, T. N. Mokaev, and
N. V. Stankevich.
\emph{Hidden Attractors Localization in Chua Circuit via the Describing Function Method}.
IFAC-PapersOnLine, vol. 50, no. 1, pp. 2651--2656, 2017.
DOI:\@ \href{https://doi.org/10.1016/j.ifacol.2017.08.470}{doi}.

\bibitem{GenesioTesi1993}
R. Genesio, A. Tesi, and F. Villoresi.
\emph{A Frequency Approach for Analyzing and Controlling Chaos in Nonlinear Circuits}.
IEEE Transactions on Circuits and Systems I, 1993.

\bibitem{DiethelmFordFreed2002}
K. Diethelm, N. J. Ford, and A. D. Freed.
\emph{A Predictor-Corrector Approach for the Numerical Solution of Fractional Differential Equations}.
Nonlinear Dynamics, 2002, DOI:\@ \href{https://doi.org/10.1023/A:1016592219341}{doi}.

\bibitem{Podlubny1999}
I. Podlubny.
\emph{Fractional Differential Equations}.
Academic Press, 1999.

\bibitem{Ghoreishi2023}
F. Ghoreishi, R. Ghaffari, and N. Saad.
\emph{Fractional Order Runge-Kutta Methods}.
Preprint \mbox{arXiv:2210.13138}, 2023.

\bibitem{SciPy2020}
P. Virtanen et al.
\emph{SciPy 1.0: Fundamental Algorithms for Scientific Computing in Python}.
Nature Methods, vol. 17, pp. 261--272, 2020.
\par\noindent DOI:\@ \href{https://doi.org/10.1038/s41592-019-0686-2}{doi}.

\bibitem{CooleyTukey1965}
J. W. Cooley and J. W. Tukey.
\emph{An Algorithm for the Machine Calculation of Complex Fourier Series}.
Mathematics of Computation, vol. 19, no. 90, pp. 297--301, 1965.
\par\noindent DOI:\@ \href{https://doi.org/10.1090/S0025-5718-1965-0178586-1}{doi}.

\bibitem{Shannon1948}
C. E. Shannon.
\emph{A Mathematical Theory of Communication}.
Bell System Technical Journal, vol. 27, pp. 379--423 and 623--656, 1948.

\bibitem{Welch1967}
P. D. Welch.
\emph{The Use of Fast Fourier Transform for the Estimation of Power Spectra:
A Method Based on Time Averaging over Short, Modified Periodograms}.
IEEE Transactions on Audio and Electroacoustics, vol. 15, no. 2,
pp. 70--73, 1967.
\par\noindent DOI:\@ \href{https://doi.org/10.1109/TAU.1967.1161901}{doi}.

\bibitem{Henon1982}
M. Hénon.
\emph{On the Numerical Computation of Poincaré Maps}.
Physica D: Nonlinear Phenomena, vol. 5, nos. 2--3, pp. 412--414, 1982.
\par\noindent DOI:\@ \href{https://doi.org/10.1016/0167-2789(82)90034-3}{doi}.

\bibitem{GottwaldMelbourne2004}
G. A. Gottwald and I. Melbourne.
\emph{A New Test for Chaos in Deterministic Systems}.
Proceedings of the Royal Society A, vol. 460, no. 2042,
pp. 603--611, 2004.
\par\noindent DOI:\@ \href{https://doi.org/10.1098/rspa.2003.1183}{doi}.

\bibitem{GottwaldMelbourne2009}
G. A. Gottwald and I. Melbourne.
\emph{On the Implementation of the 0--1 Test for Chaos}.
SIAM Journal on Applied Dynamical Systems, vol. 8, no. 1,
pp. 129--145, 2009.
\par\noindent DOI:\@ \href{https://doi.org/10.1137/080718851}{doi}.

\bibitem{Benettin1980a}
G. Benettin, L. Galgani, A. Giorgilli, and J.-M. Strelcyn.
\emph{Lyapunov Characteristic Exponents for Smooth Dynamical Systems and for Hamiltonian Systems; A Method for Computing All of Them. Part 1: Theory}.
Meccanica, 1980.

\bibitem{Benettin1980b}
G. Benettin, L. Galgani, A. Giorgilli, and J.-M. Strelcyn.
\emph{Lyapunov Characteristic Exponents for Smooth Dynamical Systems and for Hamiltonian Systems; A Method for Computing All of Them. Part 2: Numerical Application}.
Meccanica, 1980.

\bibitem{Wolf1985}
A. Wolf, J. B. Swift, H. L. Swinney, and J. A. Vastano.
\emph{Determining Lyapunov Exponents from a Time Series}.
Physica D, vol. 16, no. 3, pp. 285--317, 1985.
\par\noindent DOI:\@ \href{https://doi.org/10.1016/0167-2789(85)90011-9}{doi}.

\bibitem{DancaKuznetsov2018}
M.-F. Danca and N. V. Kuznetsov.
\emph{Matlab Code for Lyapunov Exponents of Fractional-Order Systems}.
International Journal of Bifurcation and Chaos, vol. 28, no. 5,
1850067, 2018.
\par\noindent DOI:\@ \href{https://doi.org/10.1142/S0218127418500670}{doi}.

\bibitem{Fischer2020}
C. Fischer, K. L. A. Zourmba, and A. Mohamadou.
\emph{Lyapunov Exponents Spectrum Estimation of Fractional Order Nonlinear
Systems Using Cloned Dynamics}.
Applied Numerical Mathematics, vol. 154, pp. 187--204, 2020.
\par\noindent DOI:\@ \href{https://doi.org/10.1016/j.apnum.2020.03.027}{doi}.

\bibitem{Takens1981}
F. Takens.
\emph{Detecting Strange Attractors in Turbulence}.
In \emph{Dynamical Systems and Turbulence, Warwick 1980},
Lecture Notes in Mathematics, vol. 898, pp. 366--381, Springer, 1981.
\par\noindent DOI:\@ \href{https://doi.org/10.1007/BFb0091924}{doi}.

\bibitem{Rosenstein1993}
M. T. Rosenstein, J. J. Collins, and C. J. De Luca.
\emph{A Practical Method for Calculating Largest Lyapunov Exponents from
Small Data Sets}.
Physica D, vol. 65, nos. 1--2, pp. 117--134, 1993.
\par\noindent DOI:\@ \href{https://doi.org/10.1016/0167-2789(93)90009-P}{doi}.

\bibitem{Eckmann1986}
J.-P. Eckmann, S. O. Kamphorst, D. Ruelle, and S. Ciliberto.
\emph{Liapunov Exponents from Time Series}.
Physical Review A, vol. 34, no. 6, pp. 4971--4979, 1986.
\par\noindent DOI:\@ \href{https://doi.org/10.1103/PhysRevA.34.4971}{doi}.

\bibitem{Frederickson1983}
P. Frederickson, J. L. Kaplan, E. D. Yorke, and J. A. Yorke.
\emph{The Liapunov Dimension of Strange Attractors}.
Journal of Differential Equations, vol. 49, no. 2, pp. 185--207, 1983.
\par\noindent DOI:\@ \href{https://doi.org/10.1016/0022-0396(83)90011-6}{doi}.

\bibitem{RichmanMoorman2000}
J. S. Richman and J. R. Moorman.
\emph{Physiological Time-Series Analysis Using Approximate Entropy and
Sample Entropy}.
American Journal of Physiology---Heart and Circulatory Physiology,
vol. 278, no. 6, pp. H2039--H2049, 2000.
\par\noindent DOI:\@ \href{https://doi.org/10.1152/ajpheart.2000.278.6.H2039}{doi}.

\bibitem{GrassbergerProcaccia1983}
P. Grassberger and I. Procaccia.
\emph{Characterization of Strange Attractors}.
Physical Review Letters, vol. 50, no. 5, pp. 346--349, 1983.
\par\noindent DOI:\@ \href{https://doi.org/10.1103/PhysRevLett.50.346}{doi}.

\bibitem{GrassbergerProcaccia1983Measuring}
P. Grassberger and I. Procaccia.
\emph{Measuring the Strangeness of Strange Attractors}.
Physica D, vol. 9, nos. 1--2, pp. 189--208, 1983.
\par\noindent DOI:\@ \href{https://doi.org/10.1016/0167-2789(83)90298-1}{doi}.

\bibitem{Theiler1986}
J. Theiler.
\emph{Spurious Dimension from Correlation Algorithms Applied to Limited Time-Series Data}.
Physical Review A, vol. 34, no. 3, pp. 2427--2432, 1986.
\par\noindent DOI:\@ \href{https://doi.org/10.1103/PhysRevA.34.2427}{doi}.

\bibitem{DeshmukhEtAl2021}
V. Deshmukh, E. Bradley, J. Garland, and J. D. Meiss.
\emph{Toward Automated Extraction and Characterization of Scaling Regions in Dynamical Systems}.
Chaos, vol. 31, 123102, 2021.
\par\noindent DOI:\@ \href{https://doi.org/10.1063/5.0069365}{doi}.

\bibitem{Hurst1951}
H. E. Hurst.
\emph{Long-Term Storage Capacity of Reservoirs}.
Transactions of the American Society of Civil Engineers,
vol. 116, pp. 770--799, 1951.

\bibitem{Peng1994}
C.-K. Peng, S. V. Buldyrev, S. Havlin, M. Simons, H. E. Stanley,
and A. L. Goldberger.
\emph{Mosaic Organization of DNA Nucleotides}.
Physical Review E, vol. 49, no. 2, pp. 1685--1689, 1994.
\par\noindent DOI:\@ \href{https://doi.org/10.1103/PhysRevE.49.1685}{doi}.

\bibitem{BandtPompe2002}
C. Bandt and B. Pompe.
\emph{Permutation Entropy: A Natural Complexity Measure for Time Series}.
Physical Review Letters, vol. 88, no. 17, 174102, 2002.
\par\noindent DOI:\@ \href{https://doi.org/10.1103/PhysRevLett.88.174102}{doi}.

\bibitem{UnakafovaKeller2013EfficientPermutation}
V. A. Unakafova and K. Keller.
\emph{Efficiently Measuring Complexity on the Basis of Real-World Data}.
Entropy, vol. 15, no. 10, pp. 4392--4415, 2013.
\par\noindent DOI:\@ \href{https://doi.org/10.3390/e15104392}{doi}.

\bibitem{TraversaroEtAl2018TiedValues}
F. Traversaro, F. O. Redelico, M. R. Risk, A. C. Frery, and O. A. Rosso.
\emph{Bandt--Pompe Symbolization Dynamics for Time Series with Tied Values:
A Data-Driven Approach}.
Chaos, vol. 28, 075502, 2018.
\par\noindent DOI:\@ \href{https://doi.org/10.1063/1.5022021}{doi}.

\bibitem{ReyEtAl2023PermutationAsymptotic}
A. A. Rey, A. C. Frery, J. Gambini, and M. M. Lucini.
\emph{The Asymptotic Distribution of the Permutation Entropy}.
Chaos, vol. 33, 113108, 2023.
\par\noindent DOI:\@ \href{https://doi.org/10.1063/5.0171508}{doi}.

\bibitem{Higuchi1988}
T. Higuchi.
\emph{Approach to an Irregular Time Series on the Basis of the Fractal Theory}.
Physica D, vol. 31, no. 2, pp. 277--283, 1988.
\par\noindent DOI:\@ \href{https://doi.org/10.1016/0167-2789(88)90081-4}{doi}.

\bibitem{TavazoeiHaeri2009}
M. S. Tavazoei and M. Haeri.
\emph{A Proof for Non Existence of Periodic Solutions in Time Invariant Fractional Order Systems}.
Automatica, vol. 45, no. 8, pp. 1886--1890, 2009.
\par\noindent DOI:\@ \href{https://doi.org/10.1016/j.automatica.2009.04.001}{doi}.

\bibitem{Machado2015FDF}
J. Tenreiro Machado.
\emph{Fractional Order Describing Functions}.
Signal Processing, vol. 107, pp. 389--394, 2015.
\par\noindent DOI:\@ \href{https://doi.org/10.1016/j.sigpro.2014.05.012}{doi}.

\bibitem{GuanXie2025}
X. Guan and Y. Xie.
\emph{A Review on Methods for Localization of Hidden Attractors}.
Nonlinear Dynamics, vol. 113, pp. 22223--22255, 2025.
\par\noindent DOI:\@ \href{https://doi.org/10.1007/s11071-025-11327-5}{doi}.

\bibitem{Wu2023ArctanChua}
X. Wu, L. Fu, S. He, Z. Yao, H. Wang, and J. Han.
\emph{Hidden Attractors in a New Fractional-Order Chua System with Arctan Nonlinearity and Its DSP Implementation}.
Results in Physics, vol. 52, 106866, 2023.
\par\noindent DOI:\@ \href{https://doi.org/10.1016/j.rinp.2023.106866}{doi}.

\bibitem{Lubich1986}
C. Lubich.
\emph{Discretized Fractional Calculus}.
SIAM Journal on Mathematical Analysis, vol. 17, no. 3, pp. 704--719, 1986.
\par\noindent DOI:\@ \href{https://doi.org/10.1137/0517050}{doi}.

\bibitem{Lubich2004}
C. Lubich.
\emph{Convolution Quadrature Revisited}.
BIT Numerical Mathematics, vol. 44, pp. 503--514, 2004.
\par\noindent DOI:\@ \href{https://doi.org/10.1023/B:BITN.0000046813.23911.2D}{doi}.

\bibitem{JinLiZhou2017}
B. Jin, B. Li, and Z. Zhou.
\emph{Correction of High-Order BDF Convolution Quadrature for Fractional Evolution Equations}.
SIAM Journal on Scientific Computing, vol. 39, no. 6, pp. A3129--A3152, 2017.
\par\noindent DOI:\@ \href{https://doi.org/10.1137/17M1118816}{doi}.

\bibitem{SabzikarEtAl2015}
F. Sabzikar, M. M. Meerschaert, and J. Chen.
\emph{Tempered Fractional Calculus}.
Journal of Computational Physics, vol. 293, pp. 14--28, 2015.
\par\noindent DOI:\@ \href{https://doi.org/10.1016/j.jcp.2014.04.024}{doi}.

\bibitem{ChenDeng2015}
M. Chen and W. Deng.
\emph{Discretized Fractional Substantial Calculus}.
ESAIM: Mathematical Modelling and Numerical Analysis, vol. 49, no. 2,
pp. 373--394, 2015.
\par\noindent DOI:\@ \href{https://doi.org/10.1051/m2an/2014037}{doi}.

\bibitem{GuoEtAl2019Tempered}
L. Guo, F. Zeng, I. Turner, K. Burrage, and G. E. Karniadakis.
\emph{Efficient Multistep Methods for Tempered Fractional Calculus:
Algorithms and Simulations}.
SIAM Journal on Scientific Computing, vol. 41, no. 4,
pp. A2510--A2535, 2019.
\par\noindent DOI:\@ \href{https://doi.org/10.1137/18M1230153}{doi}.

\bibitem{TrefethenWeideman2014}
L. N. Trefethen and J. A. C. Weideman.
\emph{The Exponentially Convergent Trapezoidal Rule}.
SIAM Review, vol. 56, no. 3, pp. 385--458, 2014.
\par\noindent DOI:\@ \href{https://doi.org/10.1137/130932132}{doi}.

\bibitem{LiDengZhao2019}
C. Li, W. Deng, and L. Zhao.
\emph{Well-Posedness and Numerical Algorithm for the Tempered Fractional
Differential Equations}.
Discrete and Continuous Dynamical Systems -- B, vol. 24, no. 4,
pp. 1989--2015, 2019.
\par\noindent DOI:\@ \href{https://doi.org/10.3934/dcdsb.2019026}{doi}.

\bibitem{SamkoRoss1993}
S. G. Samko and B. Ross.
\emph{Integration and Differentiation to a Variable Fractional Order}.
Integral Transforms and Special Functions, vol. 1, no. 4, pp. 277--300, 1993.
\par\noindent DOI:\@ \href{https://doi.org/10.1080/10652469308819027}{doi}.

\bibitem{KhalilEtAl2014}
R. Khalil, M. Al Horani, A. Yousef, and M. Sababheh.
\emph{A New Definition of Fractional Derivative}.
Journal of Computational and Applied Mathematics, vol. 264, pp. 65--70, 2014.
\par\noindent DOI:\@ \href{https://doi.org/10.1016/j.cam.2014.01.002}{doi}.

\bibitem{CaputoFabrizio2015}
M. Caputo and M. Fabrizio.
\emph{A New Definition of Fractional Derivative without Singular Kernel}.
Progress in Fractional Differentiation and Applications, vol. 1, no. 2,
pp. 73--85, 2015.
\par\noindent DOI:\@ \href{https://doi.org/10.12785/pfda/010201}{doi}.

\bibitem{DiethelmEtAl2020}
K. Diethelm, R. Garrappa, A. Giusti, and M. Stynes.
\emph{Why Fractional Derivatives with Nonsingular Kernels Should Not Be Used}.
Fractional Calculus and Applied Analysis, vol. 23, pp. 610--634, 2020.
\par\noindent DOI:\@ \href{https://doi.org/10.1515/fca-2020-0032}{doi}.

\bibitem{CaputoDistributed2001}
M. Caputo.
\emph{Distributed Order Differential Equations Modelling Dielectric Induction and Diffusion}.
Fractional Calculus and Applied Analysis, vol. 4, pp. 421--442, 2001.

\bibitem{DiethelmFord2009}
K. Diethelm and N. J. Ford.
\emph{Numerical Analysis for Distributed-Order Differential Equations}.
Journal of Computational and Applied Mathematics, vol. 225, pp. 96--104, 2009.
\par\noindent DOI:\@ \href{https://doi.org/10.1016/j.cam.2008.07.018}{doi}.

\bibitem{HuEtAl2014Distributed}
Z. Hu, F. Liu, V. Anh, and I. Turner.
\emph{Numerical Methods for the Time Distributed-Order Superdiffusion Equation}.
ANZIAM Journal, vol. 55, pp. C464--C478, 2014.
\par\noindent DOI:\@ \href{https://doi.org/10.21914/ANZIAMJ.V55I0.7888}{doi}.

\bibitem{LinXu2007L1}
Y. Lin and C. Xu.
\emph{Finite Difference/Spectral Approximations for the Time-Fractional Diffusion Equation}.
Journal of Computational Physics, vol. 225, no. 2, pp. 1533--1552, 2007.
\par\noindent DOI:\@ \href{https://doi.org/10.1016/j.jcp.2007.02.001}{doi}.

\bibitem{DiethelmFord2004MultiOrder}
K. Diethelm and N. J. Ford.
\emph{Multi-Order Fractional Differential Equations and Their Numerical Solution}.
Applied Mathematics and Computation, vol. 154, no. 3, pp. 621--640, 2004.
\par\noindent DOI:\@ \href{https://doi.org/10.1016/S0096-3003(03)00739-2}{doi}.

\bibitem{RenSun2014MultiTerm}
J. Ren and Z.-Z. Sun.
\emph{Efficient and Stable Numerical Methods for Multi-Term Time Fractional Sub-Diffusion Equations}.
East Asian Journal on Applied Mathematics, vol. 4, no. 3, pp. 242--266, 2014.
\par\noindent DOI:\@ \href{https://doi.org/10.4208/EAJAM.181113.280514A}{doi}.

\bibitem{SheLiSun2022TransformedL1}
M. She, D. Li, and H.-W. Sun.
\emph{A Transformed L1 Method for Solving the Multi-Term Time-Fractional Diffusion Problem}.
Mathematics and Computers in Simulation, vol. 193, pp. 584--606, 2022.
\par\noindent DOI:\@ \href{https://doi.org/10.1016/j.matcom.2021.11.005}{doi}.

\bibitem{ZakyMachado2020Atomic}
M. A. Zaky and J. A. Tenreiro Machado.
\emph{Multi-Dimensional Spectral Tau Methods for Distributed-Order Fractional Diffusion Equations}.
Computers \& Mathematics with Applications, vol. 79, no. 2, pp. 476--488, 2020.
\par\noindent DOI:\@ \href{https://doi.org/10.1016/j.camwa.2019.07.008}{doi}.

\bibitem{AtanganaBaleanu2016}
A. Atangana and D. Baleanu.
\emph{New Fractional Derivatives with Nonlocal and Non-Singular Kernel: Theory and Application to Heat Transfer Model}.
Thermal Science, vol. 20, no. 2, pp. 763--769, 2016.
\par\noindent DOI:\@ \href{https://doi.org/10.2298/TSCI160111018A}{doi}.

\bibitem{YadavEtAl2019}
S. Yadav, R. K. Pandey, and A. K. Shukla.
\emph{Numerical Approximations of Atangana--Baleanu Caputo Derivative and Its Application}.
Chaos, Solitons \& Fractals, vol. 118, pp. 58--64, 2019.
\par\noindent DOI:\@ \href{https://doi.org/10.1016/j.chaos.2018.11.009}{doi}.

\bibitem{LeeKimJang2024}
S. Lee, H. Kim, and B. Jang.
\emph{A Novel Numerical Method for Solving Nonlinear Fractional-Order Differential Equations and Its Applications}.
Fractal and Fractional, vol. 8, no. 1, art. 65, 2024.
\par\noindent DOI:\@ \href{https://doi.org/10.3390/fractalfract8010065}{doi}.

\bibitem{Matusiak2020}
M. Matusiak.
\emph{Fast Evaluation of Grunwald--Letnikov Variable Fractional-Order Differentiation and Integration Based on the FFT Convolution}.
In \emph{Advanced, Contemporary Control}, AISC 1196, pp. 879--890, Springer, 2020.
\par\noindent DOI:\@ \href{https://doi.org/10.1007/978-3-030-50936-1_74}{doi}.

\bibitem{JaradEtAl2012}
F. Jarad, T. Abdeljawad, and D. Baleanu.
\emph{Caputo-Type Modification of the Hadamard Fractional Derivatives}.
Advances in Difference Equations, 2012:142, 2012.
\par\noindent DOI:\@ \href{https://doi.org/10.1186/1687-1847-2012-142}{doi}.

\bibitem{YinEtAl2024}
B. Yin, G. Zhang, Y. Liu, and H. Li.
\emph{Convolution Quadrature for Hadamard Fractional Calculus and Correction Methods for the Subdiffusion with Singular Source Terms}.
Communications in Nonlinear Science and Numerical Simulation, vol. 138,
108221, 2024.
\par\noindent DOI:\@ \href{https://doi.org/10.1016/j.cnsns.2024.108221}{doi}.

\bibitem{Zheng2021}
X. Zheng.
\emph{Logarithmic Transformation Between (Variable-Order) Caputo and Caputo--Hadamard Fractional Problems and Applications}.
Applied Mathematics Letters, vol. 121, 107366, 2021.
\par\noindent DOI:\@ \href{https://doi.org/10.1016/j.aml.2021.107366}{doi}.

\bibitem{GreenLiuYan2021}
C. W. H. Green, Y. Liu, and Y. Yan.
\emph{Numerical Methods for Caputo--Hadamard Fractional Differential Equations with Graded and Non-Uniform Meshes}.
Mathematics, vol. 9, no. 21, 2728, 2021.
\par\noindent DOI:\@ \href{https://doi.org/10.3390/math9212728}{doi}.

\bibitem{TavaresEtAl2016}
D. Tavares, R. Almeida, and D. F. M. Torres.
\emph{Caputo Derivatives of Fractional Variable Order: Numerical Approximations}.
Communications in Nonlinear Science and Numerical Simulation, vol. 35,
pp. 69--87, 2016.
\par\noindent DOI:\@ \href{https://doi.org/10.1016/j.cnsns.2015.10.027}{doi}.

\bibitem{FangEtAl2020}
Z. W. Fang, H. W. Sun, and H. Wang.
\emph{A Fast Method for Variable-Order Caputo Fractional Derivative with Applications to Time-Fractional Diffusion Equations}.
Computers \& Mathematics with Applications, vol. 80, pp. 1443--1458, 2020.
\par\noindent DOI:\@ \href{https://doi.org/10.1016/j.camwa.2020.07.009}{doi}.

\bibitem{DazaEtAl2016}
A. Daza, A. Wagemakers, B. Georgeot, D. Guery-Odelin, and
M. A. F. Sanjuan.
\emph{Basin Entropy: A New Tool to Analyze Uncertainty in Dynamical Systems}.
Scientific Reports, vol. 6, 31416, 2016.
\par\noindent DOI:\@ \href{https://doi.org/10.1038/srep31416}{doi}.

\bibitem{McDonaldEtAl1985}
S. W. McDonald, C. Grebogi, E. Ott, and J. A. Yorke.
\emph{Structure and Crises of Fractal Basin Boundaries}.
Physics Letters A, vol. 107, pp. 51--54, 1985.
\par\noindent DOI:\@ \href{https://doi.org/10.1016/0375-9601(85)90193-8}{doi}.

\bibitem{FraserSwinney1986}
A. M. Fraser and H. L. Swinney.
\emph{Independent Coordinates for Strange Attractors from Mutual Information}.
Physical Review A, vol. 33, pp. 1134--1140, 1986.
\par\noindent DOI:\@ \href{https://doi.org/10.1103/PhysRevA.33.1134}{doi}.

\bibitem{KennelEtAl1992}
M. B. Kennel, R. Brown, and H. D. I. Abarbanel.
\emph{Determining Embedding Dimension for Phase-Space Reconstruction Using a Geometrical Construction}.
Physical Review A, vol. 45, pp. 3403--3411, 1992.
\par\noindent DOI:\@ \href{https://doi.org/10.1103/PhysRevA.45.3403}{doi}.

\bibitem{EckmannRecurrence1987}
J.-P. Eckmann, S. O. Kamphorst, and D. Ruelle.
\emph{Recurrence Plots of Dynamical Systems}.
Europhysics Letters, vol. 4, no. 9, pp. 973--977, 1987.
\par\noindent DOI:\@ \href{https://doi.org/10.1209/0295-5075/4/9/004}{doi}.

\bibitem{MarwanEtAl2007}
N. Marwan, M. C. Romano, M. Thiel, and J. Kurths.
\emph{Recurrence Plots for the Analysis of Complex Systems}.
Physics Reports, vol. 438, pp. 237--329, 2007.
\par\noindent DOI:\@ \href{https://doi.org/10.1016/j.physrep.2006.11.001}{doi}.

\bibitem{Skokos2001AlignmentIndices}
Ch. Skokos.
\emph{Alignment Indices: A New, Simple Method for Determining the Ordered or Chaotic Nature of Orbits}.
Journal of Physics A: Mathematical and General, vol. 34, no. 47,
pp. 10029--10043, 2001.
\par\noindent DOI:\@ \href{https://doi.org/10.1088/0305-4470/34/47/309}{doi}.

\bibitem{SkokosEtAl2004SALI}
Ch. Skokos, Ch. Antonopoulos, T. C. Bountis, and M. N. Vrahatis.
\emph{Detecting Order and Chaos in Hamiltonian Systems by the SALI Method}.
Journal of Physics A: Mathematical and General, vol. 37, no. 24,
pp. 6269--6284, 2004.
\par\noindent DOI:\@ \href{https://doi.org/10.1088/0305-4470/37/24/006}{doi}.

\bibitem{SkokosBountisAntonopoulos2007GALI}
Ch. Skokos, T. C. Bountis, and Ch. Antonopoulos.
\emph{Geometrical Properties of Local Dynamics in Hamiltonian Systems: The Generalized Alignment Index (GALI) Method}.
Physica D: Nonlinear Phenomena, vol. 231, no. 1, pp. 30--54, 2007.
\par\noindent DOI:\@ \href{https://doi.org/10.1016/j.physd.2007.04.004}{doi}.

\bibitem{MandaHillebrandSkokos2025GALI}
B. M. Manda, M. Hillebrand, and Ch. Skokos.
\emph{Efficient Detection of Chaos through the Computation of the Generalized Alignment Index (GALI) by the Multi-Particle Method}.
Communications in Nonlinear Science and Numerical Simulation, vol. 143,
108635, 2025.
\par\noindent DOI:\@ \href{https://doi.org/10.1016/j.cnsns.2025.108635}{doi}.

\bibitem{RolimSalesEtAl2026Alignment}
M. Rolim Sales, E. D. Leonel, and C. G. Antonopoulos.
\emph{On the Behavior of Linear Dependence, Smaller, and Generalized Alignment Indices in Discrete and Continuous Chaotic Systems}.
Chaos, Solitons \& Fractals, vol. 205, 117884, 2026.
\par\noindent DOI:\@ \href{https://doi.org/10.1016/j.chaos.2026.117884}{doi}.

\bibitem{MaLongZhu2016ChaosIndicators}
D.-Z. Ma, Z.-C. Long, and Y. Zhu.
\emph{Application of Indicators for Chaos in Chaotic Circuit Systems}.
International Journal of Bifurcation and Chaos, vol. 26, no. 11,
1650182, 2016.
\par\noindent DOI:\@ \href{https://doi.org/10.1142/S0218127416501820}{doi}.

\bibitem{GinelliEtAl2007CLV}
F. Ginelli, P. Poggi, A. Turchi, H. Chat\'{e}, R. Livi, and A. Politi.
\emph{Characterizing Dynamics with Covariant Lyapunov Vectors}.
Physical Review Letters, vol. 99, no. 13, 130601, 2007.
\par\noindent DOI:\@ \href{https://doi.org/10.1103/PhysRevLett.99.130601}{doi}.

\bibitem{WolfeSamelson2007LyapunovVectors}
C. L. Wolfe and R. M. Samelson.
\emph{An Efficient Method for Recovering Lyapunov Vectors from Singular Vectors}.
Tellus A, vol. 59, no. 3, pp. 355--366, 2007.
\par\noindent DOI:\@ \href{https://doi.org/10.1111/j.1600-0870.2007.00234.x}{doi}.

\bibitem{KuptsovParlitz2012CLV}
P. V. Kuptsov and U. Parlitz.
\emph{Theory and Computation of Covariant Lyapunov Vectors}.
Journal of Nonlinear Science, vol. 22, no. 5, pp. 727--762, 2012.
\par\noindent DOI:\@ \href{https://doi.org/10.1007/s00332-012-9126-5}{doi}.

\bibitem{FroylandEtAl2013CLVComparison}
G. Froyland, T. H\"{u}ls, G. P. Morriss, and T. M. Watson.
\emph{Computing Covariant Lyapunov Vectors, Oseledets Vectors, and Dichotomy Projectors: A Comparative Numerical Study}.
Physica D: Nonlinear Phenomena, vol. 247, no. 1, pp. 18--39, 2013.
\par\noindent DOI:\@ \href{https://doi.org/10.1016/j.physd.2012.12.005}{doi}.

\bibitem{GinelliEtAl2013CLVReview}
F. Ginelli, H. Chat\'{e}, R. Livi, and A. Politi.
\emph{Covariant Lyapunov Vectors}.
Journal of Physics A: Mathematical and Theoretical, vol. 46, no. 25,
254005, 2013.
\par\noindent DOI:\@ \href{https://doi.org/10.1088/1751-8113/46/25/254005}{doi}.

\bibitem{Noethen2021CLVHilbert}
F. Noethen.
\emph{Computing Covariant Lyapunov Vectors in Hilbert Spaces}.
Journal of Computational Dynamics, vol. 8, no. 3, pp. 325--352, 2021.
\par\noindent DOI:\@ \href{https://doi.org/10.3934/jcd.2021014}{doi}.

\bibitem{DuPlessisHillebrandSkokos2026CLV}
J.-J. du Plessis, M. Hillebrand, and Ch. Skokos.
\emph{On the Efficient Numerical Computation of Covariant Lyapunov Vectors}.
Physica D: Nonlinear Phenomena, vol. 493, 135237, 2026.
\par\noindent DOI:\@ \href{https://doi.org/10.1016/j.physd.2026.135237}{doi}.

\end{thebibliography}

\end{document}
